\documentclass{amsart}
% For dealing with references we use the comment environment
\usepackage{verbatim}
\newenvironment{reference}{\comment}{\endcomment}
%\newenvironment{reference}{}{}
\newenvironment{slogan}{\comment}{\endcomment}
\newenvironment{history}{\comment}{\endcomment}

% For commutative diagrams we use Xy-pic
\usepackage[all]{xy}

% We use 2cell for 2-commutative diagrams.
\xyoption{2cell}
\UseAllTwocells

% We use multicol for the list of chapters between chapters
\usepackage{multicol}

% This is generally recommended for better output
\usepackage{lmodern}
\usepackage[T1]{fontenc}

% For cross-file-references

% Package for hypertext links:
\usepackage{hyperref}

% For any local file, say "hello.tex" you want to link to please

% Theorem environments.
%
\theoremstyle{plain}
\newtheorem{theorem}[subsection]{Theorem}
\newtheorem{proposition}[subsection]{Proposition}
\newtheorem{lemma}[subsection]{Lemma}

\theoremstyle{definition}
\newtheorem{definition}[subsection]{Definition}
\newtheorem{example}[subsection]{Example}
\newtheorem{exercise}[subsection]{Exercise}
\newtheorem{situation}[subsection]{Situation}

\theoremstyle{remark}
\newtheorem{remark}[subsection]{Remark}
\newtheorem{remarks}[subsection]{Remarks}

\numberwithin{equation}{subsection}

% Macros
%
\def\lim{\mathop{\mathrm{lim}}\nolimits}
\def\colim{\mathop{\mathrm{colim}}\nolimits}
\def\Spec{\mathop{\mathrm{Spec}}}
\def\Hom{\mathop{\mathrm{Hom}}\nolimits}
\def\Ext{\mathop{\mathrm{Ext}}\nolimits}
\def\SheafHom{\mathop{\mathcal{H}\!\mathit{om}}\nolimits}
\def\SheafExt{\mathop{\mathcal{E}\!\mathit{xt}}\nolimits}
\def\Sch{\mathit{Sch}}
\def\Mor{\mathop{\mathrm{Mor}}\nolimits}
\def\Ob{\mathop{\mathrm{Ob}}\nolimits}
\def\Sh{\mathop{\mathit{Sh}}\nolimits}
\def\NL{\mathop{N\!L}\nolimits}
\def\CH{\mathop{\mathrm{CH}}\nolimits}
\def\proetale{{pro\text{-}\acute{e}tale}}
\def\etale{{\acute{e}tale}}
\def\QCoh{\mathit{QCoh}}
\def\Ker{\mathop{\mathrm{Ker}}}
\def\Im{\mathop{\mathrm{Im}}}
\def\Coker{\mathop{\mathrm{Coker}}}
\def\Coim{\mathop{\mathrm{Coim}}}

% Boxtimes
%
\DeclareMathSymbol{\boxtimes}{\mathbin}{AMSa}{"02}

%
% Macros for moduli stacks/spaces
%
\def\QCohstack{\mathcal{QC}\!\mathit{oh}}
\def\Cohstack{\mathcal{C}\!\mathit{oh}}
\def\Spacesstack{\mathcal{S}\!\mathit{paces}}
\def\Quotfunctor{\mathrm{Quot}}
\def\Hilbfunctor{\mathrm{Hilb}}
\def\Curvesstack{\mathcal{C}\!\mathit{urves}}
\def\Polarizedstack{\mathcal{P}\!\mathit{olarized}}
\def\Complexesstack{\mathcal{C}\!\mathit{omplexes}}
% \Pic is the operator that assigns to X its picard group, usage \Pic(X)
% \Picardstack_{X/B} denotes the Picard stack of X over B
% \Picardfunctor_{X/B} denotes the Picard functor of X over B
\def\Pic{\mathop{\mathrm{Pic}}\nolimits}
\def\Picardstack{\mathcal{P}\!\mathit{ic}}
\def\Picardfunctor{\mathrm{Pic}}
\def\Deformationcategory{\mathcal{D}\!\mathit{ef}}

\setlength{\emergencystretch}{3em}
\begin{document}
\title[Stacks Project, Tag 05PF]{1512 --- The isomorphism locus of a surjection from a finitely presented base-flat relatively pure sheaf is a closed base subscheme}
\maketitle
\noindent Copyright (C) 2005--2025 Johan de Jong. Stacks Project authors. Source: \href{https://stacks.math.columbia.edu/tag/05PF}{Tag 05PF}.

\noindent Editorial difficulty note (not author text). Difficulty H2: This global representability theorem combines henselian projective etale models, universal purity and constructible relative assassins to reduce arbitrary base changes to a coefficient ideal in a projective module. The closed universal base locus is obtained without Noetherian hypotheses and its finite-presentation assertion needs limit control..

\noindent Editorial provenance note (not author text). History: excerpt from revision \texttt{a0444\allowbreak 6e57e\allowbreak c1fbc\allowbreak 252a8\allowbreak 71afc\allowbreak ec775\allowbreak 2fb28\allowbreak 07b14}, flat.tex, lines 6670--6750. Reference presentation was adapted; original-author context and core support blocks were retained or added. The mathematical wording is unchanged. Licensed under GNU FDL 1.2 or later; no invariant sections or cover texts. See ../licenses/COPYING. Context blocks reproduce author definitions and supporting results; they are not additional counted problems.

\begin{definition}
\label{divisors-definition-weakly-associated}
Let $X$ be a scheme.
Let $\mathcal{F}$ be a quasi-coherent sheaf on $X$.
\begin{enumerate}
\item We say $x \in X$ is {\it weakly associated} to $\mathcal{F}$
if the maximal ideal $\mathfrak m_x$ is weakly associated to the
$\mathcal{O}_{X, x}$-module $\mathcal{F}_x$.
\item We denote $\text{WeakAss}(\mathcal{F})$ the set of weakly associated
points of $\mathcal{F}$.
\item The {\it weakly associated points of $X$} are the weakly associated
points of $\mathcal{O}_X$.
\end{enumerate}
\end{definition}

\begin{definition}
\label{algebra-definition-weakly-associated}
Let $R$ be a ring. Let $M$ be an $R$-module.
A prime $\mathfrak p$ of $R$ is {\it weakly associated} to $M$
if there exists an element $m \in M$ such that $\mathfrak p$ is minimal
among the prime ideals containing the annihilator
$\text{Ann}(m) = \{f \in R \mid fm = 0\}$.
The set of all such primes is denoted $\text{WeakAss}_R(M)$
or $\text{WeakAss}(M)$.
\end{definition}

\begin{definition}
\label{more-algebra-definition-auto-ass}
A ring $R$ is said to be {\it auto-associated} if $R$ is local and its
maximal ideal $\mathfrak m$ is weakly associated to $R$.
\end{definition}

\begin{definition}
\label{definition-one-step-devissage}
Let $S$ be a scheme.
Let $X$ be locally of finite type over $S$.
Let $\mathcal{F}$ be a quasi-coherent $\mathcal{O}_X$-module of finite type.
Let $s \in S$ be a point.
A {\it one step d\'evissage of $\mathcal{F}/X/S$ over $s$}
is given by morphisms of schemes over $S$
$$
\xymatrix{
X & Z \ar[l]_i \ar[r]^\pi & Y
}
$$
and a quasi-coherent $\mathcal{O}_Z$-module $\mathcal{G}$ of finite type
such that
\begin{enumerate}
\item $X$, $S$, $Z$ and $Y$ are affine,
\item $i$ is a closed immersion of finite presentation,
\item $\mathcal{F} \cong i_*\mathcal{G}$,
\item $\pi$ is finite, and
\item the structure morphism $Y \to S$ is smooth with
geometrically irreducible fibres of
dimension $\dim(\text{Supp}(\mathcal{F}_s))$.
\end{enumerate}
In this case we say $(Z, Y, i, \pi, \mathcal{G})$ is a one step
d\'evissage of $\mathcal{F}/X/S$ over $s$.
\end{definition}

\begin{definition}
\label{definition-one-step-devissage-at-x}
Let $S$ be a scheme.
Let $X$ be locally of finite type over $S$.
Let $\mathcal{F}$ be a quasi-coherent $\mathcal{O}_X$-module of finite type.
Let $x \in X$ be a point with image $s$ in $S$.
A {\it one step d\'evissage of $\mathcal{F}/X/S$ at $x$}
is a system $(Z, Y, i, \pi, \mathcal{G}, z, y)$, where
$(Z, Y, i, \pi, \mathcal{G})$ is a one step d\'evissage of
$\mathcal{F}/X/S$ over $s$ and
\begin{enumerate}
\item $\dim_x(\text{Supp}(\mathcal{F}_s)) = \dim(\text{Supp}(\mathcal{F}_s))$,
\item $z \in Z$ is a point with $i(z) = x$ and $\pi(z) = y$,
\item we have $\pi^{-1}(\{y\}) = \{z\}$,
\item the extension $\kappa(y)/\kappa(s)$ is purely
transcendental.
\end{enumerate}
\end{definition}

\begin{definition}
\label{definition-complete-devissage}
Let $S$ be a scheme.
Let $X$ be locally of finite type over $S$.
Let $\mathcal{F}$ be a quasi-coherent $\mathcal{O}_X$-module of finite type.
Let $s \in S$ be a point.
A {\it complete d\'evissage of $\mathcal{F}/X/S$ over $s$} is given by a
diagram
$$
\xymatrix{
X & Z_1 \ar[l]^{i_1} \ar[d]^{\pi_1} \\
& Y_1 & Z_2 \ar[l]^{i_2} \ar[d]^{\pi_2} \\
& & Y_2 & Z_3 \ar[l] \ar[d] \\
& & & ... & ... \ar[l] \ar[d] \\
& & & & Y_n
}
$$
of schemes over $S$, finite type quasi-coherent $\mathcal{O}_{Z_k}$-modules
$\mathcal{G}_k$, and $\mathcal{O}_{Y_k}$-module maps
$$
\alpha_k :
\mathcal{O}_{Y_k}^{\oplus r_k}
\longrightarrow
\pi_{k, *}\mathcal{G}_k,
\quad
k = 1, \ldots, n
$$
satisfying the following properties:
\begin{enumerate}
\item $(Z_1, Y_1, i_1, \pi_1, \mathcal{G}_1)$ is a one step
d\'evissage of $\mathcal{F}/X/S$ over $s$,
\item the map $\alpha_k$ induces an isomorphism
$$
\kappa(\xi_k)^{\oplus r_k} \longrightarrow
(\pi_{k, *}\mathcal{G}_k)_{\xi_k}
\otimes_{\mathcal{O}_{Y_k, \xi_k}} \kappa(\xi_k)
$$
where $\xi_k \in (Y_k)_s$ is the unique generic point,
\item for $k = 2, \ldots, n$ the system
$(Z_k, Y_k, i_k, \pi_k, \mathcal{G}_k)$
is a one step d\'evissage of $\Coker(\alpha_{k - 1})/Y_{k - 1}/S$
over $s$,
\item $\Coker(\alpha_n) = 0$.
\end{enumerate}
In this case we say that
$(Z_k, Y_k, i_k, \pi_k, \mathcal{G}_k, \alpha_k)_{k = 1, \ldots, n}$
is a complete d\'evissage of $\mathcal{F}/X/S$ over $s$.
\end{definition}

\begin{definition}
\label{definition-complete-devissage-at-x}
Let $S$ be a scheme.
Let $X$ be locally of finite type over $S$.
Let $\mathcal{F}$ be a quasi-coherent $\mathcal{O}_X$-module of finite type.
Let $x \in X$ be a point with image $s \in S$.
A {\it complete d\'evissage of $\mathcal{F}/X/S$ at $x$} is given by a
system
$$
(Z_k, Y_k, i_k, \pi_k, \mathcal{G}_k, \alpha_k, z_k, y_k)_{k = 1, \ldots, n}
$$
such that $(Z_k, Y_k, i_k, \pi_k, \mathcal{G}_k, \alpha_k)$ is a
complete d\'evissage of $\mathcal{F}/X/S$ over $s$, and such that
\begin{enumerate}
\item $(Z_1, Y_1, i_1, \pi_1, \mathcal{G}_1, z_1, y_1)$ is a one step
d\'evissage of $\mathcal{F}/X/S$ at $x$,
\item for $k = 2, \ldots, n$ the system
$(Z_k, Y_k, i_k, \pi_k, \mathcal{G}_k, z_k, y_k)$
is a one step d\'evissage of $\Coker(\alpha_{k - 1})/Y_{k - 1}/S$
at $y_{k - 1}$.
\end{enumerate}
\end{definition}

\begin{definition}
\label{definition-complete-devissage-algebra}
Let $R \to S$ be a finite type ring map.
Let $\mathfrak r$ be a prime of $R$.
Let $N$ be a finite $S$-module.
A {\it complete d\'evissage of $N/S/R$ over $\mathfrak r$}
is given by $R$-algebra maps
$$
\xymatrix{
& A_1 & & A_2 & & ... & & A_n \\
S \ar[ru] & & B_1 \ar[lu] \ar[ru] & & ... \ar[lu] \ar[ru] & &
... \ar[lu] \ar[ru] & & B_n \ar[lu]
}
$$
finite $A_i$-modules $M_i$ and $B_i$-module maps
$\alpha_i : B_i^{\oplus r_i} \to M_i$ such that
\begin{enumerate}
\item $S \to A_1$ is surjective and of finite presentation,
\item $B_i \to A_{i + 1}$ is surjective and of finite presentation,
\item $B_i \to A_i$ is finite,
\item $R \to B_i$ is smooth with geometrically irreducible fibres,
\item $N \cong M_1$ as $S$-modules,
\item $\Coker(\alpha_i) \cong M_{i + 1}$ as $B_i$-modules,
\item $\alpha_i : \kappa(\mathfrak p_i)^{\oplus r_i}
\to M_i \otimes_{B_i} \kappa(\mathfrak p_i)$ is an isomorphism
where $\mathfrak p_i = \mathfrak rB_i$, and
\item $\Coker(\alpha_n) = 0$.
\end{enumerate}
In this situation we say that
$(A_i, B_i, M_i, \alpha_i)_{i = 1, \ldots, n}$
is a complete d\'evissage of $N/S/R$ over $\mathfrak r$.
\end{definition}

\begin{definition}
\label{definition-complete-devissage-at-x-algebra}
Let $R \to S$ be a finite type ring map.
Let $\mathfrak q$ be a prime of $S$ lying over the prime $\mathfrak r$ of $R$.
Let $N$ be a finite $S$-module.
A {\it complete d\'evissage of $N/S/R$ at $\mathfrak q$} is given by a
complete d\'evissage $(A_i, B_i, M_i, \alpha_i)_{i = 1, \ldots, n}$
of $N/S/R$ over $\mathfrak r$ and prime ideals $\mathfrak q_i \subset B_i$
lying over $\mathfrak r$ such that
\begin{enumerate}
\item $\kappa(\mathfrak r) \subset \kappa(\mathfrak q_i)$ is purely
transcendental,
\item there is a unique prime $\mathfrak q'_i \subset A_i$
lying over $\mathfrak q_i \subset B_i$,
\item $\mathfrak q = \mathfrak q'_1 \cap S$ and
$\mathfrak q_i = \mathfrak q'_{i + 1} \cap A_i$,
\item $R \to B_i$ has relative dimension
$\dim_{\mathfrak q_i}(\text{Supp}(M_i \otimes_R \kappa(\mathfrak r)))$.
\end{enumerate}
\end{definition}

\begin{remark}
\label{remark-finite-presentation}
Note that the $R$-algebras $B_i$ for all $i$ and $A_i$ for $i \geq 2$
are of finite presentation over $R$. If $S$ is of finite presentation over
$R$, then it is also the case that $A_1$ is of finite presentation over
$R$. In this case all the ring maps in the complete d\'evissage are of
finite presentation. See
Algebra, Lemma \href{https://stacks.math.columbia.edu/tag/00F4}{lemma compose finite type (Tag 00F4)}.
Still assuming $S$ of finite presentation over $R$
the following are equivalent
\begin{enumerate}
\item $M$ is of finite presentation over $S$,
\item $M_1$ is of finite presentation over $A_1$,
\item $M_1$ is of finite presentation over $B_1$,
\item each $M_i$ is of finite presentation both as an $A_i$-module
and as a $B_i$-module.
\end{enumerate}
The equivalences (1) $\Leftrightarrow$ (2) and (2) $\Leftrightarrow$ (3)
follow from
Algebra, Lemma \href{https://stacks.math.columbia.edu/tag/0564}{lemma finite finitely presented extension (Tag 0564)}.
If $M_1$ is finitely presented, so is $\Coker(\alpha_1)$ (see
Algebra, Lemma \href{https://stacks.math.columbia.edu/tag/0519}{lemma extension (Tag 0519)})
and hence $M_2$, etc.
\end{remark}

\begin{remark}
\label{remark-same-notion}
Let $A \to B$ be a finite type ring map and let $N$ be a finite
$B$-module. Let $\mathfrak q$ be a prime of $B$ lying over the prime
$\mathfrak r$ of $A$. Set $X = \Spec(B)$, $S = \Spec(A)$ and
$\mathcal{F} = \widetilde{N}$ on $X$. Let $x$ be the point corresponding
to $\mathfrak q$ and let $s \in S$ be the point corresponding to
$\mathfrak p$. Then
\begin{enumerate}
\item if there exists a complete d\'evissage of $\mathcal{F}/X/S$
over $s$ then there exists a complete d\'evissage of
$N/B/A$ over $\mathfrak p$, and
\item there exists a complete d\'evissage of $\mathcal{F}/X/S$
at $x$ if and only if there exists a complete d\'evissage of
$N/B/A$ at $\mathfrak q$.
\end{enumerate}
There is just a small twist in that we omitted the condition on
the relative dimension in the formulation of ``a complete d\'evissage of
$N/B/A$ over $\mathfrak p$'' which is why the implication in (1)
only goes in one direction.
The notion of a complete d\'evissage at
$\mathfrak q$ does have this condition built in. In any case we will
only use that existence for $\mathcal{F}/X/S$
implies the existence for $N/B/A$.
\end{remark}

\begin{lemma}
\label{lemma-elementary-devissage-variant}
Let $f : X \to S$ be morphism of schemes which is locally of finite type.
Let $\mathcal{F}$ be a finite type quasi-coherent $\mathcal{O}_X$-module.
Let $x \in X$ with image $s = f(x)$ in $S$.
Then there exists a commutative diagram of pointed schemes
$$
\xymatrix{
(X, x) \ar[d]_f & (X', x') \ar[l]^g \ar[d] \\
(S, s) & (S', s') \ar[l] \\
}
$$
such that $(S', s') \to (S, s)$ and $(X', x') \to (X, x)$
are elementary \'etale neighbourhoods, and such that
$g^*\mathcal{F}/X'/S'$ has a one step d\'evissage at $x'$.
\end{lemma}

\begin{definition}
\label{algebra-definition-standard-smooth}
Let $R$ be a ring. Given integers $n \geq c \geq 0$ and
$f_1, \ldots, f_c \in R[x_1, \ldots, x_n]$ we say
$$
R[x_1, \ldots, x_n]/(f_1, \ldots, f_c)
$$
is a {\it standard smooth algebra over $R$} if the polynomial
$$
g =
\det
\left(
\begin{matrix}
\partial f_1/\partial x_1 &
\partial f_2/\partial x_1 &
\ldots &
\partial f_c/\partial x_1 \\
\partial f_1/\partial x_2 &
\partial f_2/\partial x_2 &
\ldots &
\partial f_c/\partial x_2 \\
\ldots & \ldots & \ldots & \ldots \\
\partial f_1/\partial x_c &
\partial f_2/\partial x_c &
\ldots &
\partial f_c/\partial x_c
\end{matrix}
\right)
$$
maps to an invertible element in $R[x_1, \ldots, x_n]/(f_1, \ldots, f_c)$.
We say an $R$-algebra $S$ is {\it standard smooth} or that
the ring map $R \to S$ is {\it standard smooth} if there exist
$n \geq c \geq 0$ and $f_1, \ldots, f_c \in R[x_1, \ldots, x_n]$
such that $R[x_1, \ldots, x_n]/(f_1, \ldots, f_c)$ is a standard smooth
algebra over $R$ and $S$ is isomorphic to
$R[x_1, \ldots, x_n]/(f_1, \ldots, f_c)$ as an $R$-algebra.
\end{definition}

\begin{definition}
\label{definition-shrink}
Let $S$, $X$, $\mathcal{F}$, $x$, $s$ be as in
Definition \ref{definition-one-step-devissage-at-x}.
Let $(Z, Y, i, \pi, \mathcal{G}, z, y)$ be a one step d\'evissage
of $\mathcal{F}/X/S$ at $x$. Let us define a
{\it standard shrinking} of this situation to be
given by standard opens $S' \subset S$, $X' \subset X$, $Z' \subset Z$,
and $Y' \subset Y$ such that $s \in S'$, $x \in X'$, $z \in Z'$, and
$y \in Y'$ and such that
$$
(Z', Y', i|_{Z'}, \pi|_{Z'}, \mathcal{G}|_{Z'}, z, y)
$$
is a one step d\'evissage of $\mathcal{F}|_{X'}/X'/S'$ at $x$.
\end{definition}

\begin{definition}
\label{definition-shrink-complete}
Let $S$, $X$, $\mathcal{F}$, $x$, $s$ be as in
Definition \ref{definition-complete-devissage-at-x}.
Consider a complete d\'evissage
$(Z_k, Y_k, i_k, \pi_k, \mathcal{G}_k, \alpha_k, z_k, y_k)_{k = 1, \ldots, n}$
of $\mathcal{F}/X/S$ at $x$. Let us define a
{\it standard shrinking} of this situation to be
given by standard opens $S' \subset S$, $X' \subset X$,
$Z'_k \subset Z_k$, and $Y'_k \subset Y_k$ such that $s_k \in S'$,
$x_k \in X'$, $z_k \in Z'$, and $y_k \in Y'$ and such that
$$
(Z'_k, Y'_k, i'_k, \pi'_k,
\mathcal{G}'_k, \alpha'_k, z_k, y_k)_{k = 1, \ldots, n}
$$
is a one step d\'evissage of $\mathcal{F}'/X'/S'$ at $x$ where
$\mathcal{G}'_k = \mathcal{G}_k|_{Z'_k}$ and
$\mathcal{F}' = \mathcal{F}|_{X'}$.
\end{definition}

\begin{definition}
\label{divisors-definition-relative-assassin}
Let $f : X \to S$ be a morphism of schemes.
Let $\mathcal{F}$ be a quasi-coherent $\mathcal{O}_X$-module.
The {\it relative assassin of $\mathcal{F}$ in $X$ over $S$}
is the set
$$
\text{Ass}_{X/S}(\mathcal{F}) =
\bigcup\nolimits_{s \in S} \text{Ass}_{X_s}(\mathcal{F}_s)
$$
where $\mathcal{F}_s = (X_s \to X)^*\mathcal{F}$ is the restriction
of $\mathcal{F}$ to the fibre of $f$ at $s$.
\end{definition}

\begin{definition}
\label{algebra-definition-relative-assassin}
Let $R \to S$ be a ring map. Let $N$ be an $S$-module.
The {\it relative assassin of $N$ over $S/R$} is the set
$$
\text{Ass}_{S/R}(N)
=
\{ \mathfrak q \subset S \mid
\mathfrak q \in \text{Ass}_S(N \otimes_R \kappa(\mathfrak p))
\text{ with }\mathfrak p = R \cap \mathfrak q\}.
$$
This is the set named $A$ in
Lemma \href{https://stacks.math.columbia.edu/tag/05GB}{lemma compare relative assassins (Tag 05GB)}.
\end{definition}

\begin{definition}
\label{definition-pure}
Let $f : X \to S$ be a morphism of schemes which is of finite type.
Let $\mathcal{F}$ be a finite type quasi-coherent $\mathcal{O}_X$-module.
\begin{enumerate}
\item Let $s \in S$. We say $\mathcal{F}$ is {\it pure along $X_s$}
if there is no impurity $(g : T \to S, t' \leadsto t, \xi)$
of $\mathcal{F}$ above $s$ with $(T, t) \to (S, s)$ an
elementary \'etale neighbourhood.
\item We say $\mathcal{F}$ is {\it universally pure along $X_s$}
if there does not exist any impurity of $\mathcal{F}$ above $s$.
\item We say that $X$ is {\it pure along $X_s$} if $\mathcal{O}_X$
is pure along $X_s$.
\item We say $\mathcal{F}$ is {\it universally $S$-pure}, or
{\it universally pure relative to $S$} if $\mathcal{F}$ is universally
pure along $X_s$ for every $s \in S$.
\item We say $\mathcal{F}$ is {\it $S$-pure}, or
{\it pure relative to $S$} if $\mathcal{F}$ is pure along $X_s$
for every $s \in S$.
\item We say that $X$ is {\it $S$-pure} or {\it pure relative to $S$}
if $\mathcal{O}_X$ is pure relative to $S$.
\end{enumerate}
\end{definition}

\begin{definition}
\label{definition-impurity}
In
Situation \ref{situation-pre-pure}
we say a diagram (\ref{equation-impurity}) defines an
{\it impurity of $\mathcal{F}$ above $s$}
if $\xi \in \text{Ass}_{X_T/T}(\mathcal{F}_T)$ and
$\overline{\{\xi\}} \cap X_t = \emptyset$. We will indicate
this by saying ``let $(g : T \to S, t' \leadsto t, \xi)$ be
an impurity of $\mathcal{F}$ above $s$''.
\end{definition}

\begin{situation}
\label{situation-pre-pure}
Here $S$, $X$ are schemes and $f : X \to S$ is a finite type morphism.
Also, $\mathcal{F}$ is a finite type quasi-coherent $\mathcal{O}_X$-module.
Finally $s$ is a point of $S$.
\end{situation}

\noindent
In this situation consider a morphism $g : T \to S$, a point $t \in T$
with $g(t) = s$, a specialization $t' \leadsto t$, and a point
$\xi \in X_T$ in the base change of $X$ lying over $t'$. Picture
\begin{equation}
\label{equation-impurity}
\vcenter{
\xymatrix{
\xi \ar@{|->}[d] & \\
t' \ar@{~>}[r] & t \ar@{|->}[r] & s
}
}
\quad\quad
\vcenter{
\xymatrix{
X_T \ar[d] \ar[r] & X \ar[d] \\
T \ar[r]^g & S
}
}
\end{equation}
Moreover, denote $\mathcal{F}_T$ the pullback of $\mathcal{F}$ to $X_T$.

\begin{definition}
\label{properties-definition-locally-projective}
Let $X$ be a scheme. Let $\mathcal{F}$ be a quasi-coherent
$\mathcal{O}_X$-module. We say $\mathcal{F}$ is {\it locally projective}
if for every affine open $U \subset X$ the $\mathcal{O}_X(U)$-module
$\mathcal{F}(U)$ is projective.
\end{definition}

\begin{situation}
\label{situation-iso}
Let $f : X \to S$ be a morphism of schemes.
Let $u : \mathcal{F} \to \mathcal{G}$ be a homomorphism of
quasi-coherent $\mathcal{O}_X$-modules. For any scheme $T$ over
$S$ we will denote $u_T : \mathcal{F}_T \to \mathcal{G}_T$ the
base change of $u$ to $T$, in other words, $u_T$ is the pullback
of $u$ via the projection morphism $X_T = X \times_S T \to X$.
In this situation we can consider the functor
\begin{equation}
\label{equation-iso}
F_{iso} : (\Sch/S)^{opp} \longrightarrow \textit{Sets}, \quad
T \longrightarrow \left\{
\begin{matrix}
\{*\} & \text{if} & u_T \text{ is an isomorphism}, \\
\emptyset & \text{else.} &
\end{matrix}
\right.
\end{equation}
There are variants $F_{inj}$, $F_{surj}$, $F_{zero}$ where we ask that
$u_T$ is injective, surjective, or zero.
\end{situation}

\begin{definition}
\label{algebra-definition-devissage}
Let $M$ be an $R$-module.  A {\it direct sum d\'evissage} of $M$ is a family
of submodules $(M_{\alpha})_{\alpha \in S}$, indexed by an ordinal $S$ and
increasing (with respect to inclusion), such that:
\begin{enumerate}
\item[(0)] $M_0 = 0$;
\item[(1)] $M = \bigcup_{\alpha} M_{\alpha}$;
\item[(2)] if $\alpha \in S$ is a limit ordinal, then $M_{\alpha} =
\bigcup_{\beta < \alpha} M_{\beta}$;
\item[(3)] if $\alpha + 1 \in S$, then $M_{\alpha}$ is a direct summand of
$M_{\alpha + 1}$.
\end{enumerate}
If moreover
\begin{enumerate}
\item[(4)] $M_{\alpha + 1}/M_{\alpha}$ is countably generated for
$\alpha + 1 \in S$,
\end{enumerate}
then $(M_{\alpha})_{\alpha \in S}$ is called a {\it Kaplansky d\'evissage}
of $M$.
\end{definition}

\begin{definition}
\label{algebra-definition-mittag-leffler-module}
Let $M$ be an $R$-module.  We say that $M$ is {\it Mittag-Leffler} if the
equivalent conditions of
Proposition \ref{algebra-proposition-ML-characterization}
hold.
\end{definition}

\begin{proposition}
\label{algebra-proposition-ML-characterization}
Let $M$ be an $R$-module.  Let $(M_i, f_{ij})$ be a directed system of finitely
presented $R$-modules, indexed by $I$, such that $M = \colim M_i$.  Let
$f_i:
M_i \to M$ be the canonical map.  The following are equivalent:
\begin{enumerate}
\item For every finitely presented $R$-module $P$ and module map $f: P
\to M$, there exists a finitely presented $R$-module $Q$ and a module
map $g: P \to Q$ such that $g$ and $f$ dominate each other, i.e.,
$\Ker(f \otimes_R \text{id}_N) = \Ker(g \otimes_R \text{id}_N)$
for every $R$-module $N$.
\item For each $i \in I$, there exists $j \geq i$ such that $f_{ij}: M_i
\to M_j$ dominates $f_i: M_i \to M$.
\item For each $i \in I$, there exists $j \geq i$ such that $f_{ij}: M_i
\to M_j$ factors through $f_{ik}: M_i \to M_k$ for all $k \geq
i$.
\item For every $R$-module $N$, the inverse system
$(\Hom_R(M_i, N), \Hom_R(f_{ij}, N))$ is Mittag-Leffler.
\item For $N = \prod_{s \in I} M_s$, the inverse system
$(\Hom_R(M_i, N), \Hom_R(f_{ij}, N))$ is Mittag-Leffler.
\end{enumerate}
\end{proposition}

\begin{definition}
\label{algebra-definition-domination}
Let $f: M \to N$ and $g: M \to M'$ be maps of $R$-modules.
Then we say $g$ {\it dominates} $f$ if for any $R$-module $Q$, we have $\Ker(f
\otimes_R \text{id}_Q) \subset \Ker(g \otimes_R \text{id}_Q)$.
\end{definition}

\begin{definition}
\label{algebra-definition-universally-injective}
Let $f: M \to N$ be a map of $R$-modules.  Then $f$ is called
{\it universally injective} if for every $R$-module $Q$, the map $f
\otimes_R \text{id}_Q: M \otimes_R Q \to N \otimes_R Q$
is injective.  A sequence $0 \to M_1 \to M_2 \to M_3
\to 0$ of $R$-modules is called {\it universally exact} if it is exact
and $M_1 \to M_2$ is universally injective.
\end{definition}

\begin{proposition}
\label{proposition-finite-type-flat-at-point}
Let $f : X \to S$ be a morphism of schemes.
Let $\mathcal{F}$ be a quasi-coherent sheaf on $X$.
Let $x \in X$ with image $s \in S$.
Assume that
\begin{enumerate}
\item $f$ is locally of finite presentation,
\item $\mathcal{F}$ is of finite type, and
\item $\mathcal{F}$ is flat at $x$ over $S$.
\end{enumerate}
Then there exists an elementary \'etale neighbourhood $(S', s') \to (S, s)$
and an open subscheme
$$
V \subset X \times_S \Spec(\mathcal{O}_{S', s'})
$$
which contains the unique point of
$X \times_S \Spec(\mathcal{O}_{S', s'})$ mapping to $x$
such that the pullback of $\mathcal{F}$ to $V$ is an $\mathcal{O}_V$-module
of finite presentation and flat over $\mathcal{O}_{S', s'}$.
\end{proposition}

\begin{proof}[First proof]
This proof is longer but does not use the existence of a complete d\'evissage.
The problem is local around $x$ and $s$, hence we may assume that $X$
and $S$ are affine. During the proof we will finitely many times replace
$S$ by an elementary \'etale neighbourhood of $(S, s)$. The goal is then to find
(after such a replacement) an open
$V \subset X \times_S \Spec(\mathcal{O}_{S, s})$ containing $x$
such that $\mathcal{F}|_V$ is flat over $S$ and finitely presented.
Of course we may also replace $S$ by $\Spec(\mathcal{O}_{S, s})$
at any point of the proof, i.e., we may assume $S$ is a local scheme.
We will prove the proposition by induction on the integer
$n = \dim_x(\text{Supp}(\mathcal{F}_s))$.

\medskip\noindent
We can choose
\begin{enumerate}
\item elementary \'etale neighbourhoods $g : (X', x') \to (X, x)$,
$e : (S', s') \to (S, s)$,
\item a commutative diagram
$$
\xymatrix{
X \ar[dd]_f & X' \ar[dd] \ar[l]^g & Z' \ar[l]^i \ar[d]^\pi \\
& & Y' \ar[d]^h \\
S & S' \ar[l]_e & S' \ar@{=}[l]
}
$$
\item a point $z' \in Z'$ with $i(z') = x'$, $y' = \pi(z')$, $h(y') = s'$,
\item a finite type quasi-coherent $\mathcal{O}_{Z'}$-module $\mathcal{G}$,
\end{enumerate}
as in
Lemma \ref{lemma-elementary-devissage}.
We are going to replace $S$ by $\Spec(\mathcal{O}_{S', s'})$, see
remarks in first paragraph of the proof. Consider the diagram
$$
\xymatrix{
X_{\mathcal{O}_{S', s'}} \ar[ddr]_f &
X'_{\mathcal{O}_{S', s'}} \ar[dd] \ar[l]^g &
Z'_{\mathcal{O}_{S', s'}} \ar[l]^i \ar[d]^\pi \\
& & Y'_{\mathcal{O}_{S', s'}} \ar[dl]^h \\
& \Spec(\mathcal{O}_{S', s'})
}
$$
Here we have base changed the schemes $X', Z', Y'$ over $S'$ via
$\Spec(\mathcal{O}_{S', s'}) \to S'$ and the scheme $X$ over $S$ via
$\Spec(\mathcal{O}_{S', s'}) \to S$. It is still the case that
$g$ is \'etale, see
Lemma \href{https://stacks.math.columbia.edu/tag/05H2}{lemma etale at point (Tag 05H2)}.
After replacing $X$ by $X_{\mathcal{O}_{S', s'}}$,
$X'$ by $X'_{\mathcal{O}_{S', s'}}$,
$Z'$ by $Z'_{\mathcal{O}_{S', s'}}$, and
$Y'$ by $Y'_{\mathcal{O}_{S', s'}}$
we may assume we have a diagram as
Lemma \ref{lemma-elementary-devissage}
where in addition $S = S'$ is a local scheme with closed point $s$. By
Lemmas \ref{lemma-devissage-finite-presentation} and
\ref{lemma-devissage-flat}
the result for $Y' \to S$, the sheaf $\pi_*\mathcal{G}$, and the
point $y'$ implies the result for $X \to S$, $\mathcal{F}$ and $x$.
Hence we may assume that $S$ is local and $X \to S$ is a smooth morphism
of affines with geometrically irreducible fibres of dimension $n$.

\medskip\noindent
The base case of the induction: $n = 0$.
As $X \to S$ is smooth with geometrically
irreducible fibres of dimension $0$ we see that $X \to S$ is an open
immersion, see
Descent, Lemma \href{https://stacks.math.columbia.edu/tag/02LC}{lemma universally injective etale open immersion (Tag 02LC)}.
As $S$ is local and the closed point is in the image of $X \to S$
we conclude that $X = S$. Thus we see that $\mathcal{F}$ corresponds
to a finite flat $\mathcal{O}_{S, s}$ module. In this case the result
follows from
Algebra, Lemma \href{https://stacks.math.columbia.edu/tag/00NZ}{lemma finite flat local (Tag 00NZ)}
which tells us that $\mathcal{F}$ is in fact finite free.

\medskip\noindent
The induction step. Assume the result holds whenever the dimension
of the support in the closed fibre is $< n$. Write $S = \Spec(A)$,
$X = \Spec(B)$ and $\mathcal{F} = \widetilde{N}$ for some $B$-module
$N$. Note that $A$ is a local ring; denote its maximal ideal $\mathfrak m$.
Then $\mathfrak p = \mathfrak mB$ is the unique minimal prime lying over
$\mathfrak m$ as $X \to S$ has geometrically irreducible fibres. Finally,
let $\mathfrak q \subset B$ be the prime corresponding to $x$. By
Lemma \ref{lemma-induction-step}
we can choose a map
$$
\alpha : B^{\oplus r} \to N
$$
such that $\kappa(\mathfrak p)^{\oplus r} \to N \otimes_B \kappa(\mathfrak p)$
is an isomorphism. Moreover, as $N_{\mathfrak q}$ is $A$-flat the lemma
also shows that $\alpha$ is injective and that
$\Coker(\alpha)_{\mathfrak q}$ is $A$-flat.
Set $Q = \Coker(\alpha)$. Note that the support of $Q/\mathfrak mQ$
does not contain $\mathfrak p$. Hence it is certainly the case that
$\dim_{\mathfrak q}(\text{Supp}(Q/\mathfrak mQ)) < n$.
Combining everything we know about $Q$ we see
that the induction hypothesis applies to $Q$. It follows that there exists
an elementary \'etale morphism $(S', s) \to (S, s)$ such that the conclusion
holds for $Q \otimes_A A'$ over $B \otimes_A A'$ where
$A' = \mathcal{O}_{S', s'}$. After replacing $A$ by $A'$ we have an
exact sequence
$$
0 \to B^{\oplus r} \to N \to Q \to 0
$$
(here we use that $\alpha$ is injective as mentioned above)
of finite $B$-modules and we also get an element
$g \in B$, $g \not \in \mathfrak q$ such that
$Q_g$ is finitely presented over $B_g$ and flat over $A$. Since localization
is exact we see that
$$
0 \to B_g^{\oplus r} \to N_g \to Q_g \to 0
$$
is still exact. As $B_g$ and $Q_g$ are flat over $A$ we conclude that
$N_g$ is flat over $A$, see
Algebra, Lemma \href{https://stacks.math.columbia.edu/tag/00HM}{lemma flat ses (Tag 00HM)},
and as $B_g$ and $Q_g$ are finitely presented over $B_g$ the same holds
for $N_g$, see
Algebra, Lemma \href{https://stacks.math.columbia.edu/tag/0519}{lemma extension (Tag 0519)}.
\end{proof}

\begin{proof}[Second proof]
We apply
Proposition \ref{proposition-existence-complete-at-x}
to find a commutative diagram
$$
\xymatrix{
(X, x) \ar[d] & (X', x') \ar[l]^g \ar[d] \\
(S, s) & (S', s') \ar[l]
}
$$
of pointed schemes such that the horizontal
arrows are elementary \'etale neighbourhoods
and such that $g^*\mathcal{F}/X'/S'$ has a complete
d\'evissage at $x$.
(In particular $S'$ and $X'$ are affine.) By
Morphisms, Lemma \href{https://stacks.math.columbia.edu/tag/02JZ}{lemma flat permanence (Tag 02JZ)}
we see that $g^*\mathcal{F}$ is flat at $x'$ over $S$ and by
Lemma \href{https://stacks.math.columbia.edu/tag/05B9}{lemma etale flat up down (Tag 05B9)}
we see that it is flat at $x'$ over $S'$. Via
Remark \ref{remark-same-notion}
we deduce that
$$
\Gamma(X', g^*\mathcal{F})/
\Gamma(X', \mathcal{O}_{X'})/
\Gamma(S', \mathcal{O}_{S'})
$$
has a complete d\'evissage at the prime of $\Gamma(X', \mathcal{O}_{X'})$
corresponding to $x'$. We may base change this complete
d\'evissage to the local ring $\mathcal{O}_{S', s'}$
of $\Gamma(S', \mathcal{O}_{S'})$ at the prime corresponding to
$s'$. Thus
Lemma \href{https://stacks.math.columbia.edu/tag/05I4}{lemma complete devissage flat finite type module (Tag 05I4)}
implies that
$$
\Gamma(X', \mathcal{F}')
\otimes_{\Gamma(S', \mathcal{O}_{S'})}
\mathcal{O}_{S', s'}
$$
is flat over $\mathcal{O}_{S', s'}$ and of finite presentation over
$\Gamma(X', \mathcal{O}_{X'})
\otimes_{\Gamma(S', \mathcal{O}_{S'})} \mathcal{O}_{S', s'}$.
In other words, the restriction of $\mathcal{F}$ to
$X' \times_{S'} \Spec(\mathcal{O}_{S', s'})$
is of finite presentation and flat over $\mathcal{O}_{S', s'}$.
Since the morphism
$X' \times_{S'} \Spec(\mathcal{O}_{S', s'})
\to X \times_S \Spec(\mathcal{O}_{S', s'})$
is \'etale
(Lemma \href{https://stacks.math.columbia.edu/tag/05H2}{lemma etale at point (Tag 05H2)})
its image $V \subset X \times_S \Spec(\mathcal{O}_{S', s'})$
is an open subscheme, and by \'etale descent the restriction
of $\mathcal{F}$ to $V$ is of finite presentation and flat over
$\mathcal{O}_{S', s'}$. (Results used:
Morphisms, Lemma \href{https://stacks.math.columbia.edu/tag/03WT}{lemma etale open (Tag 03WT)},
Descent, Lemma \href{https://stacks.math.columbia.edu/tag/05B0}{lemma finite presentation descends (Tag 05B0)}, and
Morphisms, Lemma \href{https://stacks.math.columbia.edu/tag/02JZ}{lemma flat permanence (Tag 02JZ)}.)
\end{proof}


\begin{lemma}
\label{lemma-finite-type-flat-at-point-local}
Let $f : X \to S$ be a morphism which is locally of finite presentation.
Let $\mathcal{F}$ be a quasi-coherent $\mathcal{O}_X$-module of finite type.
If $x \in X$ and $\mathcal{F}$ is flat at $x$ over $S$, then
$\mathcal{F}_x$ is an $\mathcal{O}_{X, x}$-module of finite presentation.
\end{lemma}

\begin{proof}
Let $s = f(x)$. By
Proposition \ref{proposition-finite-type-flat-at-point}
there exists an elementary \'etale neighbourhood $(S', s') \to (S, s)$
such that the pullback of $\mathcal{F}$ to
$X \times_S \Spec(\mathcal{O}_{S', s'})$ is of
finite presentation in a neighbourhood of the point $x' \in X_{s'} = X_s$
corresponding to $x$. The ring map
$$
\mathcal{O}_{X, x} \longrightarrow
\mathcal{O}_{X \times_S \Spec(\mathcal{O}_{S', s'}), x'}
=
\mathcal{O}_{X \times_S S', x'}
$$
is flat and local as a localization of an \'etale ring map. Hence
$\mathcal{F}_x$ is of finite presentation over $\mathcal{O}_{X, x}$
by descent, see
Algebra, Lemma \href{https://stacks.math.columbia.edu/tag/03C4}{lemma descend properties modules (Tag 03C4)}
(and also that a flat local ring map is faithfully flat, see
Algebra, Lemma \href{https://stacks.math.columbia.edu/tag/00HR}{lemma local flat ff (Tag 00HR)}).
\end{proof}


\begin{lemma}
\label{lemma-elementary-devissage}
Let $f : X \to S$ be morphism of schemes which is locally of finite type.
Let $\mathcal{F}$ be a finite type quasi-coherent $\mathcal{O}_X$-module.
Let $x \in X$ with image $s = f(x)$ in $S$.
Set $\mathcal{F}_s = \mathcal{F}|_{X_s}$ and
$n = \dim_x(\text{Supp}(\mathcal{F}_s))$.
Then we can construct
\begin{enumerate}
\item elementary \'etale neighbourhoods $g : (X', x') \to (X, x)$,
$e : (S', s') \to (S, s)$,
\item a commutative diagram
$$
\xymatrix{
X \ar[dd]_f & X' \ar[dd] \ar[l]^g & Z' \ar[l]^i \ar[d]^\pi \\
& & Y' \ar[d]^h \\
S & S' \ar[l]_e & S' \ar@{=}[l]
}
$$
\item a point $z' \in Z'$ with $i(z') = x'$, $y' = \pi(z')$, $h(y') = s'$,
\item a finite type quasi-coherent $\mathcal{O}_{Z'}$-module $\mathcal{G}$,
\end{enumerate}
such that the following properties hold
\begin{enumerate}
\item $X'$, $Z'$, $Y'$, $S'$ are affine schemes,
\item $i$ is a closed immersion of finite presentation,
\item $i_*(\mathcal{G}) \cong g^*\mathcal{F}$,
\item $\pi$ is finite and $\pi^{-1}(\{y'\}) = \{z'\}$,
\item the extension $\kappa(y')/\kappa(s')$ is purely transcendental,
\item $h$ is smooth of relative dimension $n$
with geometrically integral fibres.
\end{enumerate}
\end{lemma}

\begin{proof}
Let $V \subset S$ be an affine neighbourhood of $s$.
Let $U \subset f^{-1}(V)$ be an affine neighbourhood of $x$.
Then it suffices to prove the lemma for $f|_U : U \to V$ and
$\mathcal{F}|_U$. Hence in the rest of the proof we assume that
$X$ and $S$ are affine.

\medskip\noindent
First, suppose that $X_s = \text{Supp}(\mathcal{F}_s)$, in particular
$n = \dim_x(X_s)$. Apply
More on Morphisms,
Lemmas \ref{more-morphisms-lemma-local-local-structure-finite-type} and
\ref{more-morphisms-lemma-local-local-structure-finite-type-affine}.
This gives us a commutative diagram
$$
\xymatrix{
X \ar[dd] & X' \ar[l]^g \ar[d]^\pi \\
& Y' \ar[d]^h  \\
S & S' \ar[l]_e
}
$$
and point $x' \in X'$. We set $Z' = X'$, $i = \text{id}$, and
$\mathcal{G} = g^*\mathcal{F}$ to obtain a solution in this case.

\medskip\noindent
In general choose a closed immersion $Z \to X$ and a sheaf
$\mathcal{G}$ on $Z$ as in
Lemma \ref{lemma-sheaf-lives-on-subscheme}.
Applying the result of the previous paragraph to $Z \to S$ and
$\mathcal{G}$ we obtain a diagram
$$
\xymatrix{
X \ar[dd]_f & Z \ar[l] \ar[dd]_{f|_Z} & Z' \ar[l]^g \ar[d]^\pi \\
& & Y' \ar[d]^h \\
S & S \ar@{=}[l] & S' \ar[l]_e
}
$$
and point $z' \in Z'$ satisfying all the required properties.
We will use
Lemma \href{https://stacks.math.columbia.edu/tag/057R}{lemma lift etale (Tag 057R)}
to embed $Z'$ into a scheme \'etale over $X$. We cannot apply the lemma directly
as we want $X'$ to be a scheme over $S'$. Instead we
consider the morphisms
$$
\xymatrix{
Z' \ar[r] & Z \times_S S' \ar[r] & X \times_S S'
}
$$
The first morphism is \'etale by
Morphisms, Lemma \href{https://stacks.math.columbia.edu/tag/02GW}{lemma etale permanence (Tag 02GW)}.
The second is a closed immersion as a base change of a closed immersion.
Finally, as $X$, $S$, $S'$, $Z$, $Z'$ are all affine we may apply
Lemma \href{https://stacks.math.columbia.edu/tag/057R}{lemma lift etale (Tag 057R)}
to get an \'etale morphism of affine schemes $X' \to X \times_S S'$ such that
$$
Z' = (Z \times_S S') \times_{(X \times_S S')} X' = Z \times_X X'.
$$
As $Z \to X$ is a closed immersion of finite presentation, so is $Z' \to X'$.
Let $x' \in X'$ be the point corresponding to $z' \in Z'$.
Then the completed diagram
$$
\xymatrix{
X \ar[dd] & X' \ar[dd] \ar[l] & Z' \ar[l]^i \ar[d]^\pi \\
& & Y' \ar[d]^h \\
S & S' \ar[l]_e & S' \ar@{=}[l]
}
$$
is a solution of the original problem.
\end{proof}


\begin{lemma}
\label{lemma-sheaf-lives-on-subscheme}
Let $f : X \to S$ be a finite type morphism of affine schemes.
Let $\mathcal{F}$ be a finite type quasi-coherent $\mathcal{O}_X$-module.
Let $x \in X$ with image $s = f(x)$ in $S$.
Set $\mathcal{F}_s = \mathcal{F}|_{X_s}$.
Then there exist a closed immersion $i : Z \to X$ of finite presentation,
and a quasi-coherent finite type $\mathcal{O}_Z$-module $\mathcal{G}$
such that $i_*\mathcal{G} = \mathcal{F}$ and
$Z_s = \text{Supp}(\mathcal{F}_s)$.
\end{lemma}

\begin{proof}
Say the morphism $f : X \to S$ is given by the ring map
$A \to B$ and that $\mathcal{F}$ is the quasi-coherent sheaf
associated to the $B$-module $M$. By
Morphisms, Lemma \href{https://stacks.math.columbia.edu/tag/01T2}{lemma locally finite type characterize (Tag 01T2)}
we know that $A \to B$ is a finite type ring map, and by
Properties, Lemma \href{https://stacks.math.columbia.edu/tag/01PB}{lemma finite type module (Tag 01PB)}
we know that $M$ is a finite $B$-module. In particular the
support of $\mathcal{F}$ is the closed subscheme of $\Spec(B)$
cut out by the annihilator
$I = \{x \in B \mid xm = 0\ \forall m \in M\}$ of $M$, see
Algebra, Lemma \href{https://stacks.math.columbia.edu/tag/00L2}{lemma support closed (Tag 00L2)}.
Let $\mathfrak q \subset B$ be the prime ideal corresponding to $x$
and let $\mathfrak p \subset A$ be the prime ideal corresponding to $s$.
Note that $X_s = \Spec(B \otimes_A \kappa(\mathfrak p))$ and
that $\mathcal{F}_s$ is the quasi-coherent sheaf associated to the
$B \otimes_A \kappa(\mathfrak p)$ module $M \otimes_A \kappa(\mathfrak p)$. By
Morphisms, Lemma \href{https://stacks.math.columbia.edu/tag/056J}{lemma support finite type (Tag 056J)}
the support of $\mathcal{F}_s$ is equal to
$V(I(B \otimes_A \kappa(\mathfrak p)))$. Since
$B \otimes_A \kappa(\mathfrak p)$ is of finite type over $\kappa(\mathfrak p)$
there exist finitely many elements $f_1, \ldots, f_m \in I$
such that
$$
I(B \otimes_A \kappa(\mathfrak p)) =
(f_1, \ldots, f_n)(B \otimes_A \kappa(\mathfrak p)).
$$
Denote $i : Z \to X$ the closed subscheme cut out by
$(f_1, \ldots, f_m)$, in a formula $Z = \Spec(B/(f_1, \ldots, f_m))$.
Since $M$ is annihilated by $I$ we can think of $M$ as an
$B/(f_1, \ldots, f_m)$-module. In other words, $\mathcal{F}$ is the
pushforward of a finite type module on $Z$.
As $Z_s = \text{Supp}(\mathcal{F}_s)$ by construction, this
proves the lemma.
\end{proof}


\begin{lemma}
\label{algebra-lemma-lift-etale}
\begin{slogan}
\'Etale ring maps lift along surjections of rings
\end{slogan}
Let $R$ be a ring and let $I \subset R$ be an ideal.
Let $R/I \to \overline{S}$ be an \'etale ring map.
Then there exists an \'etale ring map
$R \to S$ such that $\overline{S} \cong S/IS$ as $R/I$-algebras.
\end{lemma}

\begin{proof}
By Lemma \href{https://stacks.math.columbia.edu/tag/00U9}{lemma etale standard smooth (Tag 00U9)} we can write
$\overline{S} =
(R/I)[x_1, \ldots, x_n]/(\overline{f}_1, \ldots, \overline{f}_n)$
as in Definition \ref{algebra-definition-standard-smooth} with
$\overline{\Delta} =
\det(\frac{\partial \overline{f}_i}{\partial x_j})_{i, j = 1, \ldots, n}$
invertible in $\overline{S}$. Just take some lifts $f_i$ and set
$S = R[x_1, \ldots, x_n, x_{n+1}]/(f_1, \ldots, f_n, x_{n + 1}\Delta - 1)$
where $\Delta = \det(\frac{\partial f_i}{\partial x_j})_{i, j = 1, \ldots, n}$
as in Example \href{https://stacks.math.columbia.edu/tag/00T8}{example make standard smooth (Tag 00T8)}.
This proves the lemma.
\end{proof}


\begin{lemma}
\label{lemma-devissage-finite-presentation}
Assumptions and notation as in
Lemma \ref{lemma-elementary-devissage}.
If $f$ is locally of finite presentation
then $\pi$ is of finite presentation.
In this case the following are equivalent
\begin{enumerate}
\item $\mathcal{F}$ is an $\mathcal{O}_X$-module of finite presentation
in a neighbourhood of $x$,
\item $\mathcal{G}$ is an $\mathcal{O}_{Z'}$-module of finite presentation
in a neighbourhood of $z'$, and
\item $\pi_*\mathcal{G}$ is an $\mathcal{O}_{Y'}$-module of
finite presentation in a neighbourhood of $y'$.
\end{enumerate}
Still assuming $f$ locally of finite presentation the following are
equivalent to each other
\begin{enumerate}
\item[(a)] $\mathcal{F}_x$ is an $\mathcal{O}_{X, x}$-module of finite
presentation,
\item[(b)] $\mathcal{G}_{z'}$ is an $\mathcal{O}_{Z', z'}$-module of
finite presentation, and
\item[(c)] $(\pi_*\mathcal{G})_{y'}$ is an $\mathcal{O}_{Y', y'}$-module
of finite presentation.
\end{enumerate}
\end{lemma}

\begin{proof}
Assume $f$ locally of finite presentation. Then $Z' \to S$ is locally
of finite presentation as a composition of such, see
Morphisms, Lemma \href{https://stacks.math.columbia.edu/tag/01TR}{lemma composition finite presentation (Tag 01TR)}.
Note that $Y' \to S$ is also locally of finite presentation as a composition
of a smooth and an \'etale morphism. Hence
Morphisms, Lemma \href{https://stacks.math.columbia.edu/tag/02FV}{lemma finite presentation permanence (Tag 02FV)}
implies $\pi$ is locally of finite presentation.
Since $\pi$ is finite we conclude that it is also separated and
quasi-compact, hence $\pi$ is actually of finite presentation.

\medskip\noindent
To prove the equivalence of (1), (2), and (3) we also consider:
(4) $g^*\mathcal{F}$ is a $\mathcal{O}_{X'}$-module of finite presentation
in a neighbourhood of $x'$. The pullback of a module of finite presentation
is of finite presentation, see
Modules, Lemma \href{https://stacks.math.columbia.edu/tag/01BQ}{lemma pullback finite presentation (Tag 01BQ)}.
Hence (1) $\Rightarrow$ (4).
The \'etale morphism $g$ is open, see
Morphisms, Lemma \href{https://stacks.math.columbia.edu/tag/03WT}{lemma etale open (Tag 03WT)}.
Hence for any open neighbourhood $U' \subset X'$ of $x'$, the image
$g(U')$ is an open neighbourhood of $x$ and the map
$\{U' \to g(U')\}$ is an \'etale covering. Thus (4) $\Rightarrow$ (1) by
Descent, Lemma \href{https://stacks.math.columbia.edu/tag/05B0}{lemma finite presentation descends (Tag 05B0)}.
Using
Descent, Lemma \href{https://stacks.math.columbia.edu/tag/05B4}{lemma finite finitely presented module (Tag 05B4)}
and some easy topological arguments (see
More on Morphisms,
Lemma \ref{more-morphisms-lemma-finite-morphism-single-point-in-fibre})
we see that
(4) $\Leftrightarrow$ (2) $\Leftrightarrow$ (3).

\medskip\noindent
To prove the equivalence of (a), (b), (c) consider the ring maps
$$
\mathcal{O}_{X, x} \to
\mathcal{O}_{X', x'} \to
\mathcal{O}_{Z', z'} \leftarrow
\mathcal{O}_{Y', y'}
$$
The first ring map is faithfully flat. Hence
$\mathcal{F}_x$ is of finite presentation over $\mathcal{O}_{X, x}$
if and only if $g^*\mathcal{F}_{x'}$ is of finite presentation over
$\mathcal{O}_{X', x'}$, see
Algebra, Lemma \href{https://stacks.math.columbia.edu/tag/03C4}{lemma descend properties modules (Tag 03C4)}.
The second ring map is surjective (hence finite) and
finitely presented by assumption, hence
$g^*\mathcal{F}_{x'}$ is of finite presentation over $\mathcal{O}_{X', x'}$
if and only if $\mathcal{G}_{z'}$ is of finite presentation over
$\mathcal{O}_{Z', z'}$, see
Algebra, Lemma \href{https://stacks.math.columbia.edu/tag/0564}{lemma finite finitely presented extension (Tag 0564)}.
Because $\pi$ is finite, of finite presentation, and
$\pi^{-1}(\{y'\}) = \{x'\}$ the ring homomorphism
$\mathcal{O}_{Y', y'} \leftarrow \mathcal{O}_{Z', z'}$ is finite
and of finite presentation, see
More on Morphisms,
Lemma \ref{more-morphisms-lemma-finite-morphism-single-point-in-fibre}.
Hence $\mathcal{G}_{z'}$ is of finite presentation over $\mathcal{O}_{Z', z'}$
if and only if $\pi_*\mathcal{G}_{y'}$ is of finite presentation over
$\mathcal{O}_{Y', y'}$, see
Algebra, Lemma \href{https://stacks.math.columbia.edu/tag/0564}{lemma finite finitely presented extension (Tag 0564)}.
\end{proof}


\begin{lemma}
\label{lemma-devissage-flat}
Assumptions and notation as in
Lemma \ref{lemma-elementary-devissage}.
The following are equivalent
\begin{enumerate}
\item $\mathcal{F}$ is flat over $S$ in a neighbourhood of $x$,
\item $\mathcal{G}$ is flat over $S'$ in a neighbourhood of $z'$, and
\item $\pi_*\mathcal{G}$ is flat over $S'$ in a neighbourhood of $y'$.
\end{enumerate}
The following are equivalent also
\begin{enumerate}
\item[(a)] $\mathcal{F}_x$ is flat over $\mathcal{O}_{S, s}$,
\item[(b)] $\mathcal{G}_{z'}$ is flat over $\mathcal{O}_{S', s'}$, and
\item[(c)] $(\pi_*\mathcal{G})_{y'}$ is flat over $\mathcal{O}_{S', s'}$.
\end{enumerate}
\end{lemma}

\begin{proof}
To prove the equivalence of (1), (2), and (3) we also consider:
(4) $g^*\mathcal{F}$ is flat over $S$ in a neighbourhood of $x'$.
We will use
Lemma \href{https://stacks.math.columbia.edu/tag/05B9}{lemma etale flat up down (Tag 05B9)}
to equate flatness over $S$ and $S'$ without further mention.
The \'etale morphism $g$ is flat and open, see
Morphisms, Lemma \href{https://stacks.math.columbia.edu/tag/03WT}{lemma etale open (Tag 03WT)}.
Hence for any open neighbourhood $U' \subset X'$ of $x'$, the image
$g(U')$ is an open neighbourhood of $x$ and the map
$U' \to g(U')$ is surjective and flat.
Thus (4) $\Leftrightarrow$ (1) by
Morphisms, Lemma \href{https://stacks.math.columbia.edu/tag/02JZ}{lemma flat permanence (Tag 02JZ)}.
Note that
$$
\Gamma(X', g^*\mathcal{F}) =
\Gamma(Z', \mathcal{G}) =
\Gamma(Y', \pi_*\mathcal{G})
$$
Hence the flatness of $g^*\mathcal{F}$, $\mathcal{G}$ and $\pi_*\mathcal{G}$
over $S'$ are all equivalent (this uses that $X'$, $Z'$, $Y'$, and
$S'$ are all affine). Some omitted topological arguments (compare
More on Morphisms,
Lemma \ref{more-morphisms-lemma-finite-morphism-single-point-in-fibre})
regarding affine neighbourhoods now show that
(4) $\Leftrightarrow$ (2) $\Leftrightarrow$ (3).

\medskip\noindent
To prove the equivalence of (a), (b), (c) consider the commutative diagram
of local ring maps
$$
\xymatrix{
\mathcal{O}_{X', x'} \ar[r]_\iota &
\mathcal{O}_{Z', z'} &
\mathcal{O}_{Y', y'} \ar[l]^\alpha &
\mathcal{O}_{S', s'} \ar[l]^\beta \\
\mathcal{O}_{X, x} \ar[u]^\gamma & & &
\mathcal{O}_{S, s} \ar[lll]_\varphi \ar[u]_\epsilon
}
$$
We will use
Lemma \href{https://stacks.math.columbia.edu/tag/05BA}{lemma etale flat up down local ring (Tag 05BA)}
to equate flatness over $\mathcal{O}_{S, s}$ and $\mathcal{O}_{S', s'}$
without further mention.
The map $\gamma$ is faithfully flat. Hence
$\mathcal{F}_x$ is flat over $\mathcal{O}_{S, s}$
if and only if $g^*\mathcal{F}_{x'}$ is flat over
$\mathcal{O}_{S', s'}$, see
Algebra, Lemma \href{https://stacks.math.columbia.edu/tag/0584}{lemma flatness descends more general (Tag 0584)}.
As $\mathcal{O}_{S', s'}$-modules the modules
$g^*\mathcal{F}_{x'}$, $\mathcal{G}_{z'}$, and
$\pi_*\mathcal{G}_{y'}$ are all isomorphic, see
More on Morphisms,
Lemma \ref{more-morphisms-lemma-finite-morphism-single-point-in-fibre}.
This finishes the proof.
\end{proof}


\begin{lemma}
\label{more-morphisms-lemma-local-local-structure-finite-type}
\begin{slogan}
A morphism of finite type is, in \'etale neighbourhoods, finite over a
smooth morphism.
\end{slogan}
Let $f : X \to S$ be a morphism. Let $x \in X$ and set $s = f(x)$.
Assume that $f$ is locally of finite type and that $n = \dim_x(X_s)$.
Then there exists a commutative diagram
$$
\xymatrix{
X \ar[dd] & X' \ar[l]^g \ar[d]^\pi & x \ar@{|->}[dd] &
x' \ar@{|->}[l] \ar@{|->}[d] \\
& Y' \ar[d]^h & & y' \ar@{|->}[d] \\
S & S' \ar[l]_e & s & s' \ar@{|->}[l]
}
$$
and a point $x' \in X'$ with $g(x') = x$ such that with $y' = \pi(x')$,
$s' = h(y')$ we have
\begin{enumerate}
\item $h : Y' \to S'$ is smooth of relative dimension $n$,
\item all fibres of $Y' \to S'$ are geometrically integral,
\item $g : (X', x') \to (X, x)$ is an elementary \'etale neighbourhood,
\item $\pi$ is finite, and $\pi^{-1}(\{y'\}) = \{x'\}$,
\item $\kappa(y')$ is a purely transcendental extension of $\kappa(s')$, and
\item $e : (S', s') \to (S, s)$ is an elementary \'etale neighbourhood.
\end{enumerate}
Moreover, if $f$ is locally of finite presentation, then $\pi$ is
of finite presentation.
\end{lemma}

\begin{proof}
The question is local on $S$, hence we may replace $S$ by an affine open
neighbourhood of $s$. Next, we apply
Lemma \href{https://stacks.math.columbia.edu/tag/052E}{lemma local structure finite type (Tag 052E)}
to get a commutative diagram
$$
\xymatrix{
X \ar[dd] & X' \ar[l]^g \ar[d]^\pi & x \ar@{|->}[dd] &
x' \ar@{|->}[l] \ar@{|->}[d] \\
& Y \ar[d]^h & & y \ar@{|->}[d] \\
S \ar@{=}[r] & S & s & s \ar@{=}[l]
}
$$
where $h$ is smooth of relative dimension $n$ and $\kappa(y)$ is
a purely transcendental extension of $\kappa(s)$. Since the question is
local on $X$ also, we may replace $Y$ by an affine neighbourhood of $y$
(and $X'$ by the inverse image of this under $\pi$). As $S$ is affine
this guarantees that $Y \to S$ is quasi-compact, separated and smooth,
in particular of finite presentation.
Let $T$ be the connected component of $Y_s$ containing $y$.
As $Y_s$ is Noetherian we see that $T$ is open.
We also see that $T$ is geometrically connected over $\kappa(s)$ by
Varieties,
Lemma \href{https://stacks.math.columbia.edu/tag/04KV}{lemma geometrically connected if connected and point (Tag 04KV)}.
Since $T$ is also smooth over $\kappa(s)$ it is geometrically normal, see
Varieties, Lemma \href{https://stacks.math.columbia.edu/tag/056T}{lemma smooth geometrically normal (Tag 056T)}.
We conclude that $T$ is geometrically irreducible over $\kappa(s)$ (as a
connected Noetherian normal scheme is irreducible, see
Properties, Lemma \href{https://stacks.math.columbia.edu/tag/033M}{lemma normal Noetherian (Tag 033M)}).
Finally, note that the smooth morphism $h$ is normal by
Lemma \href{https://stacks.math.columbia.edu/tag/056W}{lemma smooth normal (Tag 056W)}.
At this point we have verified all assumption of
Lemma \href{https://stacks.math.columbia.edu/tag/057J}{lemma normal morphism irreducible (Tag 057J)}
hold for the morphism $h : Y \to S$ and open $T \subset Y_s$.
As a result of applying
Lemma \href{https://stacks.math.columbia.edu/tag/057J}{lemma normal morphism irreducible (Tag 057J)}
we obtain $e : S' \to S$, $s' \in S'$, $Y'$ as
in the commutative diagram
$$
\xymatrix{
X \ar[dd] & X' \ar[l]^g \ar[d]^\pi & X' \times_Y Y' \ar[l] \ar[d] &
x \ar@{|->}[dd] & x' \ar@{|->}[l] \ar@{|->}[d] &
(x', s') \ar@{|->}[l] \ar@{|->}[d] \\
& Y \ar[d]^h & Y' \ar[d] \ar[l] & & y \ar@{|->}[d] &
(y, s') \ar@{|->}[l] \ar@{|->}[d] \\
S \ar@{=}[r] & S & S' \ar[l]_e & s & s \ar@{=}[l] & s' \ar@{|->}[l]
}
$$
where $e : (S', s') \to (S, s)$ is an elementary \'etale neighbourhood,
and where $Y' \subset Y_{S'}$ is an open neighbourhood all of whose fibres
over $S'$ are geometrically irreducible, such that $Y'_{s'} = T$ via
the identification $Y_s = Y_{S', s'}$. Let $(y, s') \in Y'$ be the point
corresponding to $y \in T$; this is also the unique point of $Y \times_S S'$
lying over $y$ with residue field equal to $\kappa(y)$ which maps to $s'$
in $S'$. Similarly, let $(x', s') \in X' \times_Y Y' \subset X' \times_S S'$
be the unique point over $x'$ with residue field equal to $\kappa(x')$
lying over $s'$. Then the outer part of this diagram is a solution to the
problem posed in the lemma. Some minor details omitted.
\end{proof}


\begin{lemma}
\label{more-morphisms-lemma-local-local-structure-finite-type-affine}
Assumption and notation as in
Lemma \ref{more-morphisms-lemma-local-local-structure-finite-type}.
In addition to properties (1) -- (6) we may also arrange it so that
\begin{enumerate}
\item[(7)] $S'$, $Y'$, $X'$ are affine.
\end{enumerate}
\end{lemma}

\begin{proof}
Note that if $Y'$ is affine, then $X'$ is affine as $\pi$ is finite.
Choose an affine open neighbourhood $U' \subset S'$ of $s'$.
Choose an affine open neighbourhood $V' \subset h^{-1}(U')$ of $y'$.
Let $W' = h(V')$. This is an open neighbourhood of $s'$ in $S'$, see
Morphisms, Lemma \href{https://stacks.math.columbia.edu/tag/056G}{lemma smooth open (Tag 056G)},
contained in $U'$. Choose an affine open neighbourhood $U'' \subset W'$
of $s'$. Then $h^{-1}(U'') \cap V'$ is affine because it is equal to
$U'' \times_{U'} V'$. By construction
$h^{-1}(U'') \cap V' \to U''$ is a surjective smooth morphism whose
fibres are (nonempty) open subschemes of geometrically integral fibres
of $Y' \to S'$, and hence geometrically integral. Thus we may replace
$S'$ by $U''$ and $Y'$ by $h^{-1}(U'') \cap V'$.
\end{proof}


\begin{lemma}
\label{more-morphisms-lemma-finite-morphism-single-point-in-fibre}
Let $\pi : X \to Y$ be a finite morphism.
Let $x \in X$ with $y = \pi(x)$ such that $\pi^{-1}(\{y\}) = \{x\}$.
Then
\begin{enumerate}
\item For every neighbourhood $U \subset X$ of $x$ in $X$, there
exists a neighbourhood $V \subset Y$ of $y$ such that
$\pi^{-1}(V) \subset U$.
\item The ring map $\mathcal{O}_{Y, y} \to \mathcal{O}_{X, x}$
is finite.
\item If $\pi$ is of finite presentation, then
$\mathcal{O}_{Y, y} \to \mathcal{O}_{X, x}$ is of finite presentation.
\item For any quasi-coherent $\mathcal{O}_X$-module $\mathcal{F}$
we have $\mathcal{F}_x = \pi_*\mathcal{F}_y$ as
$\mathcal{O}_{Y, y}$-modules.
\end{enumerate}
\end{lemma}

\begin{proof}
The first assertion is purely topological; use that
$\pi$ is a continuous and closed map such that $\pi^{-1}(\{y\}) = \{x\}$.
To prove the second and third parts we may assume
$X = \Spec(B)$ and $Y = \Spec(A)$. Then
$A \to B$ is a finite ring map and $y$ corresponds to a prime
$\mathfrak p$ of $A$ such that there exists a unique prime $\mathfrak q$ of
$B$ lying over $\mathfrak p$. Then
$B_{\mathfrak q} = B_{\mathfrak p}$, see
Algebra, Lemma \href{https://stacks.math.columbia.edu/tag/00EA}{lemma unique prime over localize below (Tag 00EA)}.
In other words, the map $A_{\mathfrak p} \to B_{\mathfrak q}$
is equal to the map $A_{\mathfrak p} \to B_{\mathfrak p}$ you get
from localizing $A \to B$ at $\mathfrak p$.
Thus (2) and (3) follow from simple properties of localization
(some details omitted). For the final statement, suppose that
$\mathcal{F} = \widetilde M$ for some $B$-module $M$.
Then $\mathcal{F} = M_{\mathfrak q}$ and
$\pi_*\mathcal{F}_y = M_{\mathfrak p}$. By the above these
localizations agree. Alternatively you can use part (1) and
the definition of stalks to see that $\mathcal{F}_x = \pi_*\mathcal{F}_y$
directly.
\end{proof}


\begin{lemma}
\label{lemma-induction-step}
Let $(R, \mathfrak m)$ be a local ring. Let $R \to S$ be a finitely presented
flat ring map with geometrically integral fibres. Write
$\mathfrak p = \mathfrak mS$. Let $\mathfrak q \subset S$ be a prime ideal
lying over $\mathfrak m$. Let $N$ be a finite $S$-module.
There exist $r \geq 0$ and an $S$-module map
$$
\alpha : S^{\oplus r} \longrightarrow N
$$
such that
$\alpha : \kappa(\mathfrak p)^{\oplus r} \to N \otimes_S \kappa(\mathfrak p)$
is an isomorphism. For any such $\alpha$ the following are equivalent:
\begin{enumerate}
\item $N_{\mathfrak q}$ is $R$-flat,
\item $\alpha$ is $R$-universally injective and
$\Coker(\alpha)_{\mathfrak q}$ is $R$-flat,
\item $\alpha$ is injective and
$\Coker(\alpha)_{\mathfrak q}$ is $R$-flat,
\item $\alpha_{\mathfrak p}$ is an isomorphism and
$\Coker(\alpha)_{\mathfrak q}$ is $R$-flat, and
\item $\alpha_{\mathfrak q}$ is injective and
$\Coker(\alpha)_{\mathfrak q}$ is $R$-flat.
\end{enumerate}
\end{lemma}

\begin{proof}
To obtain $\alpha$ set
$r = \dim_{\kappa(\mathfrak p)} N \otimes_S \kappa(\mathfrak p)$ and pick
$x_1, \ldots, x_r \in N$ which form a basis of
$N \otimes_S \kappa(\mathfrak p)$. Define
$\alpha(s_1, \ldots, s_r) = \sum s_i x_i$. This proves the existence.

\medskip\noindent
Fix an $\alpha$. The most interesting implication is
(1) $\Rightarrow$ (2) which we prove first. Assume (1).
Because $S/\mathfrak mS$ is a domain with fraction field $\kappa(\mathfrak p)$
we see that
$(S/\mathfrak mS)^{\oplus r} \to
N_{\mathfrak p}/\mathfrak mN_{\mathfrak p} = N \otimes_S \kappa(\mathfrak p)$
is injective. Hence by
Lemmas \ref{lemma-universally-injective-local} and
\ref{lemma-fibres-irreducible-flat-projective-nonnoetherian}.
the map $S^{\oplus r} \to N_{\mathfrak p}$ is $R$-universally injective.
It follows that $S^{\oplus r} \to N$ is $R$-universally injective, see
Algebra, Lemma \href{https://stacks.math.columbia.edu/tag/05CJ}{lemma universally injective permanence (Tag 05CJ)}.
Then also the localization $\alpha_{\mathfrak q}$ is $R$-universally
injective, see
Algebra, Lemma \href{https://stacks.math.columbia.edu/tag/05CM}{lemma universally injective localize (Tag 05CM)}.
We conclude that $\Coker(\alpha)_{\mathfrak q}$ is $R$-flat by
Algebra, Lemma \href{https://stacks.math.columbia.edu/tag/058P}{lemma ui flat domain (Tag 058P)}.

\medskip\noindent
The implication (2) $\Rightarrow$ (3) is immediate. If (3) holds, then
$\alpha_{\mathfrak p}$ is injective as a localization of an injective
module map. By Nakayama's lemma
(Algebra, Lemma \href{https://stacks.math.columbia.edu/tag/00DV}{lemma NAK (Tag 00DV)})
$\alpha_{\mathfrak p}$ is surjective too. Hence (3) $\Rightarrow$ (4).
If (4) holds, then $\alpha_{\mathfrak p}$ is an isomorphism, so
$\alpha$ is injective as $S_{\mathfrak q} \to S_{\mathfrak p}$ is injective.
Namely, elements of $S \setminus \mathfrak p$ are nonzerodivisors on $S$
by a combination of
Lemmas \ref{lemma-invert-universally-injective} and
\ref{lemma-fibres-irreducible-flat-projective-nonnoetherian}.
Hence (4) $\Rightarrow$ (5). Finally, if (5) holds, then
$N_{\mathfrak q}$ is $R$-flat as an extension of flat modules, see
Algebra, Lemma \href{https://stacks.math.columbia.edu/tag/00HM}{lemma flat ses (Tag 00HM)}.
Hence (5) $\Rightarrow$ (1) and the proof is finished.
\end{proof}


\begin{lemma}
\label{lemma-universally-injective-local}
Let $(R, \mathfrak m)$ be a local ring. Let $u : M \to N$ be an $R$-module map.
If $M$ is a projective $R$-module, $N$ is a flat $R$-module, and
$\overline{u} : M/\mathfrak mM \to N/\mathfrak mN$ is injective
then $u$ is universally injective.
\end{lemma}

\begin{proof}
By
Algebra, Theorem \href{https://stacks.math.columbia.edu/tag/0593}{theorem projective free over local ring (Tag 0593)}
the module $M$ is free. If we show the result holds for every finitely
generated direct summand of $M$, then the lemma follows. Hence we may
assume that $M$ is finite free. Write $N = \colim_i N_i$ as
a directed colimit of finite free modules, see
Algebra, Theorem \href{https://stacks.math.columbia.edu/tag/058G}{theorem lazard (Tag 058G)}.
Note that $u : M \to N$ factors through $N_i$ for some $i$ (as $M$ is finite
free). Denote $u_i : M \to N_i$ the corresponding $R$-module map.
As $\overline{u}$ is injective we see that
$\overline{u_i} : M/\mathfrak mM \to N_i/\mathfrak mN_i$ is
injective and remains injective on composing with the maps
$N_i/\mathfrak mN_i \to N_{i'}/\mathfrak mN_{i'}$ for all $i' \geq i$.
As $M$ and $N_{i'}$ are finite free over the local ring $R$ this implies
that $M \to N_{i'}$ is a split injection for all $i' \geq i$. Hence
for any $R$-module $Q$ we see that $M \otimes_R Q \to N_{i'} \otimes_R Q$
is injective for all $i' \geq i$. As $- \otimes_R Q$ commutes with
colimits we conclude that $M \otimes_R Q \to N_{i'} \otimes_R Q$
is injective as desired.
\end{proof}


\begin{lemma}
\label{lemma-invert-universally-injective}
Assumption and notation as in
Lemma \ref{lemma-homothety-spectrum}.
Assume moreover that $N$ is projective as an $R$-module.
Then each $s \in \Sigma$ defines a
universally injective $R$-module map $s : N \to N$, and the
map $N \to \Sigma^{-1}N$ is $R$-universally injective.
\end{lemma}

\begin{proof}
Pick $s \in \Sigma$. By
Lemma \ref{lemma-universally-injective-local}
the map $s : N \to N$ is universally injective.
The map $N \to \Sigma^{-1}N$ is universally injective as the directed
colimit of the maps $s : N \to N$.
\end{proof}


\begin{lemma}
\label{lemma-fibres-irreducible-flat-projective}
Let $R$ be a ring.
Let $R \to S$ be a ring map.
Assume
\begin{enumerate}
\item $R$ is Noetherian,
\item $R \to S$ is of finite type and flat, and
\item every fibre ring $S \otimes_R \kappa(\mathfrak p)$ is
geometrically integral over $\kappa(\mathfrak p)$.
\end{enumerate}
Then $S$ is projective as an $R$-module.
\end{lemma}

\begin{proof}
Consider the set
$$
\{I \subset R \mid S/IS\text{ not projective as }R/I\text{-module}\}
$$
We have to show this set is empty. To get a contradiction assume it is
nonempty. Then it contains a maximal element $I$.
Let $J = \sqrt{I}$ be its radical. If $I \not = J$, then
$S/JS$ is projective as a $R/J$-module, and $S/IS$ is flat over $R/I$
and $J/I$ is a nilpotent ideal in $R/I$. Applying
Algebra, Lemma \href{https://stacks.math.columbia.edu/tag/05CG}{lemma lift projective (Tag 05CG)}
we see that $S/IS$ is a projective $R/I$-module, which is a contradiction.
Hence we may assume that $I$ is a radical ideal. In other words we
are reduced to proving the lemma in case $R$ is a reduced ring and
$S/IS$ is a projective $R/I$-module for every nonzero ideal $I$
of $R$.

\medskip\noindent
Assume $R$ is a reduced ring and $S/IS$ is a projective $R/I$-module
for every nonzero ideal $I$ of $R$. By generic flatness,
Algebra, Lemma \href{https://stacks.math.columbia.edu/tag/051R}{lemma generic flatness Noetherian (Tag 051R)}
(applied to a localization $R_g$ which is a domain) or the more general
Algebra, Lemma \href{https://stacks.math.columbia.edu/tag/051Z}{lemma generic flatness reduced (Tag 051Z)}
there exists a nonzero $f \in R$ such that $S_f$ is free as an
$R_f$-module. Denote $R^\wedge = \lim R/(f^n)$ the $(f)$-adic completion
of $R$. Note that the ring map
$$
R \longrightarrow R_f \times R^\wedge
$$
is a faithfully flat ring map, see
Algebra, Lemma \href{https://stacks.math.columbia.edu/tag/00MB}{lemma completion flat (Tag 00MB)}.
Hence by faithfully flat descent of projectivity, see
Algebra, Theorem \ref{algebra-theorem-ffdescent-projectivity}
it suffices to prove that $S \otimes_R R^\wedge$ is a projective
$R^\wedge$-module. To see this we will use the criterion of
Lemma \ref{lemma-flat-pure-over-complete-projective}.
First of all, note that $S/fS = (S \otimes_R R^\wedge)/f(S \otimes_R R^\wedge)$
is a projective $R/(f)$-module and that $S \otimes_R R^\wedge$ is flat
and of finite type over $R^\wedge$ as a base change of such.
Next, suppose that $\mathfrak p^\wedge$ is a prime ideal
of $R^\wedge$. Let $\mathfrak p \subset R$ be the corresponding prime
of $R$. As $R \to S$ has geometrically integral fibre rings, the
same is true for the fibre rings of any base change. Hence
$\mathfrak q^\wedge = \mathfrak p^\wedge(S \otimes_R R^\wedge)$,
is a prime ideals lying over $\mathfrak p^\wedge$
and it is the unique associated prime of
$S \otimes_R \kappa(\mathfrak p^\wedge)$. Thus we win if
$f(S \otimes_R R^\wedge) + \mathfrak q^\wedge \not = S \otimes_R R^\wedge$.
This is true because $\mathfrak p^\wedge + fR^\wedge \not = R^\wedge$ as
$f$ lies in the Jacobson radical of the $f$-adically complete ring $R^\wedge$
and because $R^\wedge \to S \otimes_R R^\wedge$ is surjective on spectra
as its fibres are nonempty (irreducible spaces are nonempty).
\end{proof}


\begin{lemma}
\label{lemma-fibres-irreducible-flat-projective-nonnoetherian}
Let $R$ be a ring. Let $R \to S$ be a ring map.
Assume
\begin{enumerate}
\item $R \to S$ is of finite presentation and flat, and
\item every fibre ring $S \otimes_R \kappa(\mathfrak p)$ is
geometrically integral over $\kappa(\mathfrak p)$.
\end{enumerate}
Then $S$ is projective as an $R$-module.
\end{lemma}

\begin{proof}
We can find a cocartesian diagram of rings
$$
\xymatrix{
S_0 \ar[r] & S \\
R_0 \ar[u] \ar[r] & R \ar[u]
}
$$
such that $R_0$ is of finite type over $\mathbf{Z}$, the map
$R_0 \to S_0$ is of finite type and flat with geometrically integral
fibres, see
More on Morphisms,
Lemmas \href{https://stacks.math.columbia.edu/tag/05FE}{lemma Noetherian approximation flat (Tag 05FE)},
\href{https://stacks.math.columbia.edu/tag/05FG}{lemma Noetherian approximation geometrically reduced (Tag 05FG)},
\href{https://stacks.math.columbia.edu/tag/05FH}{lemma Noetherian approximation geometrically irreducible (Tag 05FH)},
and \href{https://stacks.math.columbia.edu/tag/05FL}{lemma Noetherian approximation combine (Tag 05FL)}.
By
Lemma \ref{lemma-fibres-irreducible-flat-projective}
we see that $S_0$ is a projective $R_0$-module. Hence $S = S_0 \otimes_{R_0} R$
is a projective $R$-module, see
Algebra, Lemma \href{https://stacks.math.columbia.edu/tag/05A3}{lemma ascend properties modules (Tag 05A3)}.
\end{proof}


\begin{lemma}
\label{lemma-flat-pure-over-complete-projective}
Let $R$ be a ring.
Let $I \subset R$ be an ideal.
Let $R \to S$ be a ring map, and $N$ an $S$-module.
Assume
\begin{enumerate}
\item $R$ is Noetherian and $I$-adically complete,
\item $R \to S$ is of finite type,
\item $N$ is a finite $S$-module,
\item $N$ is flat over $R$,
\item $N/IN$ is projective as a $R/I$-module, and
\item for any prime $\mathfrak q \subset S$ which is an associated prime of
$N \otimes_R \kappa(\mathfrak p)$ where $\mathfrak p = R \cap \mathfrak q$
we have $IS + \mathfrak q \not = S$.
\end{enumerate}
Then $N$ is projective as an $R$-module.
\end{lemma}

\begin{proof}
By
Lemma \ref{lemma-universally-injective-to-completion-flat}
the map $N \to N^\wedge$ is universally injective.
By
Lemma \ref{lemma-lift-ML}
the module $N^\wedge$ is Mittag-Leffler.
By
Algebra, Lemma \ref{algebra-lemma-pure-submodule-ML}
we conclude that $N$ is Mittag-Leffler.
Hence $N$ is countably generated, flat and Mittag-Leffler as an $R$-module,
whence projective by
Algebra, Lemma \ref{algebra-lemma-countgen-projective}.
\end{proof}


\begin{lemma}
\label{lemma-universally-injective-to-completion-flat}
Let $R$ be a ring.
Let $I \subset R$ be an ideal.
Let $R \to S$ be a ring map, and $N$ an $S$-module.
Assume
\begin{enumerate}
\item $R$ is a Noetherian ring,
\item $S$ is a Noetherian ring,
\item $N$ is a finite $S$-module,
\item $N$ is flat over $R$, and
\item for any prime $\mathfrak q \subset S$ which is an associated prime of
$N \otimes_R \kappa(\mathfrak p)$ where $\mathfrak p = R \cap \mathfrak q$
we have $IS + \mathfrak q \not = S$.
\end{enumerate}
Then the map $N \to N^\wedge$ of $N$ into the $I$-adic completion of $N$
is universally injective as a map of $R$-modules.
\end{lemma}

\begin{proof}
This follows from
Lemma \ref{lemma-universally-injective-to-completion}
because
Algebra, Lemma
\href{https://stacks.math.columbia.edu/tag/05C2}{lemma bourbaki fibres (Tag 05C2)} and
Remark \href{https://stacks.math.columbia.edu/tag/05E0}{remark bourbaki (Tag 05E0)}
guarantee that the set of associated primes of tensor products
$N \otimes_R Q$ are contained in the set of associated primes of
the modules $N \otimes_R \kappa(\mathfrak p)$.
\end{proof}


\begin{lemma}
\label{lemma-universally-injective-to-completion}
Let $R$ be a ring.
Let $I \subset R$ be an ideal.
Let $R \to S$ be a ring map, and $N$ an $S$-module.
Assume
\begin{enumerate}
\item $R$ is a Noetherian ring,
\item $S$ is a Noetherian ring,
\item $N$ is a finite $S$-module, and
\item for any finite $R$-module $Q$, any
$\mathfrak q \in \text{Ass}_S(Q \otimes_R N)$
satisfies $IS + \mathfrak q \not = S$.
\end{enumerate}
Then the map $N \to N^\wedge$ of $N$ into the $I$-adic completion of $N$
is universally injective as a map of $R$-modules.
\end{lemma}

\begin{proof}
We have to show that for any finite $R$-module $Q$ the map
$Q \otimes_R N \to Q \otimes_R N^\wedge$ is injective, see
Algebra, Theorem \href{https://stacks.math.columbia.edu/tag/058K}{theorem universally exact criteria (Tag 058K)}.
As there is a canonical map $Q \otimes_R N^\wedge \to (Q \otimes_R N)^\wedge$
it suffices to prove that the canonical map
$Q \otimes_R N \to (Q \otimes_R N)^\wedge$ is injective.
Hence we may replace $N$ by $Q \otimes_R N$ and it suffices to prove the
injectivity for the map $N \to N^\wedge$.

\medskip\noindent
Let $K = \Ker(N \to N^\wedge)$. It suffices to show that
$K_{\mathfrak q} = 0$ for $\mathfrak q \in \text{Ass}(N)$ as $N$ is a
submodule of $\prod_{\mathfrak q \in \text{Ass}(N)} N_{\mathfrak q}$, see
Algebra, Lemma \href{https://stacks.math.columbia.edu/tag/0311}{lemma zero at ass zero (Tag 0311)}.
Pick $\mathfrak q \in \text{Ass}(N)$. By the last assumption we see that
there exists a prime $\mathfrak q' \supset IS + \mathfrak q$.
Since $K_{\mathfrak q}$ is a localization of $K_{\mathfrak q'}$
it suffices to prove the vanishing of $K_{\mathfrak q'}$.
Note that $K = \bigcap I^nN$, hence
$K_{\mathfrak q'} \subset \bigcap I^nN_{\mathfrak q'}$.
Hence $K_{\mathfrak q'} = 0$ by
Algebra, Lemma \href{https://stacks.math.columbia.edu/tag/00IP}{lemma intersect powers ideal module zero (Tag 00IP)}.
\end{proof}


\begin{lemma}
\label{lemma-lift-ML}
Let $R$ be a ring. Let $I \subset R$ be an ideal.
Let $M$ be an $R$-module.
Assume
\begin{enumerate}
\item $R$ is Noetherian and $I$-adically complete,
\item $M$ is flat over $R$, and
\item $M/IM$ is a projective $R/I$-module.
\end{enumerate}
Then the $I$-adic completion $M^\wedge$ is a flat Mittag-Leffler
$R$-module.
\end{lemma}

\begin{proof}
Choose a surjection $F \to M$ where $F$ is a free $R$-module. By
Algebra, Lemma \ref{algebra-lemma-split-completed-sequence}
the module $M^\wedge$ is a direct summand of the module $F^\wedge$.
Hence it suffices to prove the lemma for $F$.
In this case the lemma follows from
Lemma \ref{lemma-completed-direct-sum-ML}.
\end{proof}


\begin{lemma}
\label{lemma-completed-direct-sum-ML}
Let $R$ be a ring. Let $I \subset R$ be an ideal. Let $A$ be a set.
Assume $R$ is Noetherian and complete with respect to $I$. The completion
$(\bigoplus\nolimits_{\alpha \in A} R)^\wedge$
is flat and Mittag-Leffler.
\end{lemma}

\begin{proof}
By
More on Algebra, Lemma
\ref{more-algebra-lemma-ui-completion-direct-sum-into-product}
the map $(\bigoplus\nolimits_{\alpha \in A} R)^\wedge
\to \prod_{\alpha \in A} R$ is universally injective.
Thus, by
Algebra, Lemmas \href{https://stacks.math.columbia.edu/tag/058P}{lemma ui flat domain (Tag 058P)} and
\ref{algebra-lemma-pure-submodule-ML}
it suffices to show that $\prod_{\alpha \in A} R$ is flat and Mittag-Leffler.
By
Algebra, Proposition \href{https://stacks.math.columbia.edu/tag/05CZ}{proposition characterize coherent (Tag 05CZ)}
(and
Algebra, Lemma \href{https://stacks.math.columbia.edu/tag/05CY}{lemma Noetherian coherent (Tag 05CY)})
we see that $\prod_{\alpha \in A} R$ is flat.
Thus we conclude because a product of copies of $R$ is Mittag-Leffler, see
Algebra, Lemma \href{https://stacks.math.columbia.edu/tag/05D0}{lemma product over Noetherian ring (Tag 05D0)}.
\end{proof}


\begin{lemma}
\label{algebra-lemma-split-completed-sequence}
Let $R$ be a ring. Let $I \subset R$ be an ideal.
Let $0 \to K \to P \to M \to 0$ be a short exact sequence of
$R$-modules. If $M$ is flat over $R$ and $M/IM$ is a projective
$R/I$-module, then the sequence of $I$-adic completions
$$
0 \to K^\wedge \to P^\wedge \to M^\wedge \to 0
$$
is a split exact sequence.
\end{lemma}

\begin{proof}
As $M$ is flat, each of the sequences
$$
0 \to K/I^nK \to P/I^nP \to M/I^nM \to 0
$$
is short exact, see
Lemma \href{https://stacks.math.columbia.edu/tag/00HL}{lemma flat tor zero (Tag 00HL)}
and the sequence $0 \to K^\wedge \to P^\wedge \to M^\wedge \to 0$
is a short exact sequence, see
Lemma \href{https://stacks.math.columbia.edu/tag/0315}{lemma completion generalities (Tag 0315)}.
It suffices to show that we can find splittings
$s_n : M/I^nM \to P/I^nP$ such that $s_{n + 1} \bmod I^n = s_n$.
We will construct these $s_n$ by induction on $n$.
Pick any splitting $s_1$, which exists as $M/IM$ is a projective $R/I$-module.
Assume given $s_n$ for some $n > 0$. Set
$P_{n + 1} = \{x \in P \mid x \bmod I^nP \in \Im(s_n)\}$.
The map $\pi : P_{n + 1}/I^{n + 1}P_{n + 1} \to M/I^{n + 1}M$ is surjective
(details omitted). As $M/I^{n + 1}M$ is projective as a $R/I^{n + 1}$-module
by
Lemma \href{https://stacks.math.columbia.edu/tag/05CG}{lemma lift projective (Tag 05CG)}
we may choose a section
$t : M/I^{n + 1}M \to P_{n + 1}/I^{n + 1}P_{n + 1}$
of $\pi$. Setting $s_{n + 1}$ equal to the composition of
$t$ with the canonical map $P_{n + 1}/I^{n + 1}P_{n + 1} \to P/I^{n + 1}P$
works.
\end{proof}


\begin{lemma}
\label{more-algebra-lemma-ui-completion-direct-sum-into-product}
Let $R$ be a ring.
Let $I \subset R$ be an ideal.
Let $A$ be a set.
Assume $R$ is Noetherian and complete with respect to $I$. There is a
canonical map
$$
\left(\bigoplus\nolimits_{\alpha \in A} R\right)^\wedge
\longrightarrow
\prod\nolimits_{\alpha \in A} R
$$
from the $I$-adic completion of the direct sum into the product
which is universally injective.
\end{lemma}

\begin{proof}
By definition an element $x$ of the left hand side is $x = (x_n)$ where
$x_n = (x_{n, \alpha}) \in \bigoplus\nolimits_{\alpha \in A} R/I^n$
such that $x_{n, \alpha} = x_{n + 1, \alpha} \bmod I^n$.
As $R = R^\wedge$ we see that for any $\alpha$ there exists a $y_\alpha \in R$
such that $x_{n, \alpha} = y_\alpha \bmod I^n$. Note that for each $n$ there
are only finitely many $\alpha$ such that the elements $x_{n, \alpha}$ are
nonzero. Conversely, given $(y_\alpha) \in \prod_\alpha R$ such that for each
$n$ there are only finitely many $\alpha$ such that $y_{\alpha} \bmod I^n$
is nonzero, then this defines an element of the left hand side.
Hence we can think of an element of the left hand side as infinite
``convergent sums'' $\sum_\alpha y_\alpha$ with $y_\alpha \in R$
such that for each $n$ there are only finitely many $y_\alpha$
which are nonzero modulo $I^n$. The displayed map maps this element
to the element to $(y_\alpha)$ in the product.
In particular the map is injective.

\medskip\noindent
Let $Q$ be a finite $R$-module. We have to show that the map
$$
Q \otimes_R \left(\bigoplus\nolimits_{\alpha \in A} R\right)^\wedge
\longrightarrow
Q \otimes_R \left(\prod\nolimits_{\alpha \in A} R\right)
$$
is injective, see
Algebra, Theorem \href{https://stacks.math.columbia.edu/tag/058K}{theorem universally exact criteria (Tag 058K)}.
Choose a presentation $R^{\oplus k} \to R^{\oplus m} \to Q \to 0$
and denote $q_1, \ldots, q_m \in Q$ the corresponding generators for $Q$.
By Artin-Rees
(Algebra, Lemma \href{https://stacks.math.columbia.edu/tag/00IN}{lemma Artin Rees (Tag 00IN)})
there exists a constant $c$ such that
$\Im(R^{\oplus k} \to R^{\oplus m}) \cap (I^N)^{\oplus m}
\subset \Im((I^{N - c})^{\oplus k} \to R^{\oplus m})$.
Let us contemplate the diagram
$$
\xymatrix{
\bigoplus_{l = 1}^k \left(\bigoplus\nolimits_{\alpha \in A} R\right)^\wedge
\ar[r] \ar[d] &
\bigoplus_{j = 1}^m \left(\bigoplus\nolimits_{\alpha \in A} R\right)^\wedge
\ar[r] \ar[d] &
Q \otimes_R \left(\bigoplus\nolimits_{\alpha \in A} R\right)^\wedge
\ar[r] \ar[d] &
0 \\
\bigoplus_{l = 1}^k \left(\prod\nolimits_{\alpha \in A} R\right)
\ar[r] &
\bigoplus_{j = 1}^m \left(\prod\nolimits_{\alpha \in A} R\right)
\ar[r] &
Q \otimes_R \left(\prod\nolimits_{\alpha \in A} R\right)
\ar[r] &
0
}
$$
with exact rows. Pick an element $\sum_j \sum_\alpha y_{j, \alpha}$ of
$\bigoplus_{j = 1, \ldots, m}
\left(\bigoplus\nolimits_{\alpha \in A} R\right)^\wedge$.
If this element maps to zero in the module
$Q \otimes_R \left(\prod\nolimits_{\alpha \in A} R\right)$,
then we see in particular that
$\sum_j q_j \otimes y_{j, \alpha} = 0$ in $Q$ for each $\alpha$.
Thus we can find an element
$(z_{1, \alpha}, \ldots, z_{k, \alpha}) \in \bigoplus_{l = 1, \ldots, k} R$
which maps to
$(y_{1, \alpha}, \ldots, y_{m, \alpha}) \in \bigoplus_{j = 1, \ldots, m} R$.
Moreover, if $y_{j, \alpha} \in I^{N_\alpha}$ for $j = 1, \ldots, m$, then
we may assume that $z_{l, \alpha} \in I^{N_\alpha - c}$ for
$l = 1, \ldots, k$.
Hence the sum $\sum_l \sum_\alpha z_{l, \alpha}$ is ``convergent'' and
defines an element of
$\bigoplus_{l = 1, \ldots, k}
\left(\bigoplus\nolimits_{\alpha \in A} R\right)^\wedge$
which maps to the element $\sum_j \sum_\alpha y_{j, \alpha}$ we started
out with. Thus the right vertical arrow is injective and we win.
\end{proof}


\begin{lemma}
\label{lemma-weak-bourbaki-pre}
Let $R \to S$ be a ring map of finite presentation.
Let $N$ be a finitely presented $S$-module
which is flat as an $R$-module. Let $M$ be an $R$-module.
Let $\mathfrak q$ be a prime of $S$ lying over $\mathfrak p \subset R$.
Then
$$
\mathfrak q \in \text{WeakAss}_S(M \otimes_R N)
\Leftrightarrow
\Big(
\mathfrak p \in \text{WeakAss}_R(M)
\text{ and }
\overline{\mathfrak q} \in \text{Ass}_{\overline{S}}(\overline{N})
\Big)
$$
Here $\overline{S} = S \otimes_R \kappa(\mathfrak p)$,
$\overline{\mathfrak q} = \mathfrak q \overline{S}$, and
$\overline{N} = N \otimes_R \kappa(\mathfrak p)$.
\end{lemma}

\begin{proof}
Pick $g \in S$ as in Lemma \ref{lemma-weak-bourbaki-pre-pre}.
Apply Proposition \ref{proposition-finite-presentation-flat-at-point}
to the morphism of schemes $\Spec(S_g) \to \Spec(R)$, the quasi-coherent
module associated to $N_g$, and the points
corresponding to the primes $\mathfrak qS_g$ and $\mathfrak p$. Translating
into algebra we obtain a commutative diagram of rings
$$
\xymatrix{
S \ar[r] & S_g \ar[r] & S' \\
& R \ar[lu] \ar[u] \ar[r] & R' \ar[u]
}
\quad\quad
\xymatrix{
\mathfrak q \ar@{-}[r] \ar@{-}[rd] &
\mathfrak qS_g \ar@{-}[d] \ar@{-}[r] & \mathfrak q' \ar@{-}[d] \\
& \mathfrak p \ar@{-}[r] & \mathfrak p'
}
$$
endowed with primes as shown, the horizontal arrows are \'etale,
and $N \otimes_S S'$ is projective as an $R'$-module. Set
$N' = N \otimes_S S'$, $M' = M \otimes_R R'$,
$\overline{S}' = S' \otimes_{R'} \kappa(\mathfrak q')$,
$\overline{\mathfrak q}' = \mathfrak q' \overline{S}'$,
and
$$
\overline{N}' = N' \otimes_{R'} \kappa(\mathfrak p') =
\overline{N} \otimes_{\overline{S}} \overline{S}'
$$
By Lemma \ref{lemma-etale-weak-assassin-up-down} we have
\begin{align*}
\text{WeakAss}_{S'}(M' \otimes_{R'} N') & =
(\Spec(S') \to \Spec(S))^{-1}\text{WeakAss}_S(M \otimes_R N) \\
\text{WeakAss}_{R'}(M') & =
(\Spec(R') \to \Spec(R))^{-1}\text{WeakAss}_R(M) \\
\text{Ass}_{\overline{S}'}(\overline{N}') & =
(\Spec(\overline{S}') \to \Spec(\overline{S}))^{-1}
\text{Ass}_{\overline{S}}(\overline{N})
\end{align*}
Use Algebra, Lemma \href{https://stacks.math.columbia.edu/tag/058A}{lemma ass weakly ass (Tag 058A)}
for $\overline{N}$ and $\overline{N}'$. In particular we have
\begin{align*}
\mathfrak q \in \text{WeakAss}_S(M \otimes_R N)
& \Leftrightarrow
\mathfrak q' \in \text{WeakAss}_{S'}(M' \otimes_{R'} N') \\
\mathfrak p \in \text{WeakAss}_R(M)
& \Leftrightarrow
\mathfrak p' \in \text{WeakAss}_{R'}(M') \\
\overline{\mathfrak q} \in \text{Ass}_{\overline{S}}(\overline{N})
& \Leftrightarrow
\overline{\mathfrak q}' \in \text{WeakAss}_{\overline{S}'}(\overline{N}')
\end{align*}
Our careful choice of $g$ and the formula for
$\text{Ass}_{\overline{S}'}(\overline{N}')$ above shows that
\begin{equation}
\label{equation-key-observation}
\text{if }\mathfrak r' \in \text{Ass}_{\overline{S}'}(\overline{N}')
\text{ lies over }\mathfrak r \subset \overline{S}\text{ then }
\mathfrak r \subset \overline{\mathfrak q}
\end{equation}
This will be a key observation later in the proof. We will use
the characterization of weakly associated primes given in
Algebra, Lemma \href{https://stacks.math.columbia.edu/tag/0566}{lemma weakly ass local (Tag 0566)} without further mention.

\medskip\noindent
Suppose that
$\overline{\mathfrak q} \not \in \text{Ass}_{\overline{S}}(\overline{N})$.
Then
$\overline{\mathfrak q}' \not \in \text{Ass}_{\overline{S}'}(\overline{N}')$.
By
Algebra, Lemmas \href{https://stacks.math.columbia.edu/tag/00LD}{lemma ass zero divisors (Tag 00LD)},
\href{https://stacks.math.columbia.edu/tag/00LC}{lemma finite ass (Tag 00LC)}, and
\href{https://stacks.math.columbia.edu/tag/00DS}{lemma silly (Tag 00DS)}
there exists an element $\overline{a}' \in \overline{\mathfrak q}'$
which is not a zerodivisor on $\overline{N}'$.
After replacing $\overline{a}'$ by $\lambda \overline{a}'$ for some nonzero
$\lambda \in \kappa(\mathfrak p)$ we can find
$a' \in \mathfrak q'$ mapping to $\overline{a}'$. By
Lemma \ref{lemma-invert-universally-injective}
the map $a' : N'_{\mathfrak p'} \to N'_{\mathfrak p'}$ is
$R'_{\mathfrak p'}$-universally injective. In particular
we see that $a' : M' \otimes_{R'} N' \to M' \otimes_{R'} N'$ is
injective after localizing at $\mathfrak p'$ and hence after
localizing at $\mathfrak q'$. Clearly this implies that
$\mathfrak q' \not \in \text{WeakAss}_{S'}(M' \otimes_{R'} N')$.
We conclude that $\mathfrak q \in \text{WeakAss}_S(M \otimes_R N)$ implies
$\overline{\mathfrak q} \in \text{Ass}_{\overline{S}}(\overline{N})$.

\medskip\noindent
Assume $\mathfrak q \in \text{WeakAss}_S(M \otimes_R N)$. We want
to show $\mathfrak p \in \text{WeakAss}_S(M)$.
Let $z \in M \otimes_R N$ be an element such that $\mathfrak q$
is minimal over $J = \text{Ann}_S(z)$.
Let $f_i \in \mathfrak p$, $i \in I$ be a set of generators of the
ideal $\mathfrak p$. Since $\mathfrak q$ lies over $\mathfrak p$, for every $i$
we can choose an $n_i \geq 1$ and $g_i \in S$, $g_i \not \in \mathfrak q$
with $g_i f_i^{n_i} \in J$, i.e., $g_i f_i^{n_i} z = 0$.
Let $z' \in (M' \otimes_{R'} N')_{\mathfrak p'}$ be the image of $z$.
Observe that $z'$ is nonzero because $z$ has nonzero image in
$(M \otimes_R N)_\mathfrak q$ and because $S_\mathfrak q \to S'_{\mathfrak q'}$
is faithfully flat. We claim that $f_i^{n_i} z' = 0$.

\medskip\noindent
Proof of the claim: Let $g'_i \in S'$ be the image of $g_i$.
By the key observation (\ref{equation-key-observation})
we find that the image $\overline{g}'_i \in \overline{S}'$
is not contained in $\mathfrak r'$ for any
$\mathfrak r' \in \text{Ass}_{\overline{S}'}(\overline{N})$.
Hence by Lemma \ref{lemma-invert-universally-injective}
we see that $g'_i : N'_{\mathfrak p'} \to N'_{\mathfrak p'}$ is
$R'_{\mathfrak p'}$-universally injective. In particular
we see that $g'_i : M' \otimes_{R'} N' \to M' \otimes_{R'} N'$ is
injective after localizating at $\mathfrak p'$. The claim
follows because $g_i f_i^{n_i} z' = 0$.

\medskip\noindent
Our claim shows that the annihilator of $z'$ in $R'_{\mathfrak p'}$
contains the elements $f_i^{n_i}$. As $R \to R'$ is \'etale we have
$\mathfrak p'R'_{\mathfrak p'} = \mathfrak pR'_{\mathfrak p'}$
by Algebra, Lemma \href{https://stacks.math.columbia.edu/tag/00U4}{lemma etale at prime (Tag 00U4)}.
Hence the annihilator of $z'$ in $R'_{\mathfrak p'}$ has radical equal to
$\mathfrak p' R_{\mathfrak p'}$ (here we use $z'$ is not zero).
On the other hand
$$
z' \in (M' \otimes_{R'} N')_{\mathfrak p'} =
M'_{\mathfrak p'} \otimes_{R'_{\mathfrak p'}} N'_{\mathfrak p'}
$$
The module $N'_{\mathfrak p'}$ is projective over the local ring
$R'_{\mathfrak p'}$ and hence free
(Algebra, Theorem \href{https://stacks.math.columbia.edu/tag/0593}{theorem projective free over local ring (Tag 0593)}).
Thus we can find a finite free direct summand $F' \subset N'_{\mathfrak p'}$
such that $z' \in M'_{\mathfrak p'} \otimes_{R'_{\mathfrak p'}} F'$.
If $F'$ has rank $n$, then we deduce that
$\mathfrak p' R'_{\mathfrak p'} \in
\text{WeakAss}_{R'_{\mathfrak p'}}({M'_{\mathfrak p'}}^{\oplus n})$.
This implies
$\mathfrak p'R'_{\mathfrak p'} \in \text{WeakAss}(M'_{\mathfrak p'})$
for example by Algebra, Lemma \href{https://stacks.math.columbia.edu/tag/0548}{lemma weakly ass (Tag 0548)}.
Then $\mathfrak p' \in \text{WeakAss}_{R'}(M')$
which in turn gives $\mathfrak p \in \text{WeakAss}_R(M)$.
This finishes the proof of the implication
``$\Rightarrow$'' of the equivalence of the lemma.

\medskip\noindent
Assume that $\mathfrak p \in \text{WeakAss}_R(M)$ and
$\overline{\mathfrak q} \in \text{Ass}_{\overline{S}}(\overline{N})$.
We want to show that $\mathfrak q$ is weakly associated to $M \otimes_R N$.
Note that $\overline{\mathfrak q}'$ is a maximal element of
$\text{Ass}_{\overline{S}'}(\overline{N}')$.
This is a consequence of (\ref{equation-key-observation})
and the fact that there are no inclusions among the primes
of $\overline{S}'$ lying over $\overline{\mathfrak q}$
(as fibres of \'etale morphisms are discrete
Morphisms, Lemma \href{https://stacks.math.columbia.edu/tag/02GL}{lemma etale over field (Tag 02GL)}).
Thus, after replacing $R, S, \mathfrak p, \mathfrak q, M, N$ by
$R', S', \mathfrak p', \mathfrak q', M', N'$
we may assume, in addition to the assumptions of the lemma, that
\begin{enumerate}
\item $\mathfrak p \in \text{WeakAss}_R(M)$,
\item $\overline{\mathfrak q} \in \text{Ass}_{\overline{S}}(\overline{N})$,
\item $N$ is projective as an $R$-module, and
\item $\overline{\mathfrak q}$ is maximal in
$\text{Ass}_{\overline{S}}(\overline{N})$.
\end{enumerate}
There is one more reduction, namely, we may replace
$R, S, M, N$ by their localizations at $\mathfrak p$.
This leads to one more condition, namely,
\begin{enumerate}
\item[(5)] $R$ is a local ring with maximal ideal $\mathfrak p$.
\end{enumerate}
We will finish by showing that (1) -- (5) imply
$\mathfrak q \in \text{WeakAss}(M \otimes_R N)$.

\medskip\noindent
Since $R$ is local and $\mathfrak p \in \text{WeakAss}_R(M)$
we can pick a $y \in M$ whose annihilator $I$ has radical
equal to $\mathfrak p$.
Write $\overline{\mathfrak q} = (\overline{g}_1, \ldots, \overline{g}_n)$
for some $\overline{g}_i \in \overline{S}$. Choose $g_i \in S$
mapping to $\overline{g}_i$.
Then $\mathfrak q = \mathfrak pS + g_1S + \ldots + g_nS$.
Consider the map
$$
\Psi : N/IN \longrightarrow (N/IN)^{\oplus n}, \quad
z \longmapsto (g_1z, \ldots, g_nz).
$$
This is a homomorphism of projective $R/I$-modules.
The local ring $R/I$ is auto-associated
(More on Algebra, Definition \ref{more-algebra-definition-auto-ass})
as $\mathfrak p/I$ is locally nilpotent.
The map $\Psi \otimes \kappa(\mathfrak p)$ is not injective, because
$\overline{\mathfrak q} \in \text{Ass}_{\overline{S}}(\overline{N})$.
Hence More on Algebra, Lemma \ref{more-algebra-lemma-P-fPD-zero}
implies $\Psi$ is not injective. Pick $z \in N/IN$
nonzero in the kernel of $\Psi$. The annihilator $J = \text{Ann}_S(z)$
contains $IS$ and $g_i$ by construction. Thus
$\sqrt{J} \subset S$ contains $\mathfrak q$.
Let $\mathfrak s \subset S$ be a prime minimal over $J$.
Then $\mathfrak q \subset \mathfrak s$,
$\mathfrak s$ lies over $\mathfrak p$, and
$\mathfrak s \in \text{WeakAss}_S(N/IN)$.
The last fact by definition of weakly associated primes.
Apply the ``$\Rightarrow$'' part of the lemma (which we've already proven)
to the ring map $R \to S$ and the modules $R/I$ and $N$
to conclude that
$\overline{\mathfrak s} \in \text{Ass}_{\overline{S}}(\overline{N})$.
Since $\overline{\mathfrak q} \subset \overline{\mathfrak s}$
the maximality of $\overline{\mathfrak q}$, see condition (4) above,
implies that $\overline{\mathfrak q} = \overline{\mathfrak s}$.
This shows that $\mathfrak q = \mathfrak s$ and we conlude
what we want.
\end{proof}


\begin{lemma}
\label{lemma-weak-bourbaki-pre-pre}
Let $R \to S$ be a ring map of finite presentation. Let $N$ be a
finitely presented $S$-module. Let $\mathfrak q \subset S$ be a prime ideal
lying over $\mathfrak p \subset R$. Set
$\overline{S} = S \otimes_R \kappa(\mathfrak p)$,
$\overline{\mathfrak q} = \mathfrak q \overline{S}$, and
$\overline{N} = N \otimes_R \kappa(\mathfrak p)$. Then
we can find a $g \in S$ with
$g \not \in \mathfrak q$ such that
$\overline{g} \in \mathfrak r$ for all
$\mathfrak r \in \text{Ass}_{\overline{S}}(\overline{N})$
such that $\mathfrak r \not \subset \overline{\mathfrak q}$.
\end{lemma}

\begin{proof}
Namely, if $\text{Ass}_{\overline{S}}(\overline{N}) =
\{\mathfrak r_1, \ldots, \mathfrak r_n\}$
(finiteness by Algebra, Lemma \href{https://stacks.math.columbia.edu/tag/00LC}{lemma finite ass (Tag 00LC)}),
then after renumbering we may assume that
$$
\mathfrak r_1 \subset \overline{\mathfrak q},
\ldots,
\mathfrak r_r \subset \overline{\mathfrak q}, \quad
\mathfrak r_{r + 1} \not \subset \overline{\mathfrak q},
\ldots,
\mathfrak r_n \not \subset \overline{\mathfrak q}
$$
Since $\overline{\mathfrak q}$ is a prime ideal we see that the product
$\mathfrak r_{r + 1} \ldots \mathfrak r_n$ is not contained in
$\overline{\mathfrak q}$ and hence we can pick an element
$a$ of $\overline{S}$ contained in
$\mathfrak r_{r + 1}, \ldots, \mathfrak r_n$ but not in
$\overline{\mathfrak q}$.
If there exists $g \in S$ mapping to $a$, then $g$
works. In general we can find a nonzero element
$\lambda \in \kappa(\mathfrak p)$
such that $\lambda a$ is the image of a $g \in S$.
\end{proof}


\begin{lemma}
\label{lemma-bourbaki-finite-type-general-base-at-point}
Let $S$ be a scheme.
Let $f : X \to S$ be locally of finite type.
Let $x \in X$ with image $s \in S$.
Let $\mathcal{F}$ be a finite type quasi-coherent sheaf on $X$.
Let $\mathcal{G}$ be a quasi-coherent sheaf on $S$.
If $\mathcal{F}$ is flat at $x$ over $S$, then
$$
x \in \text{WeakAss}_X(\mathcal{F} \otimes_{\mathcal{O}_X} f^*\mathcal{G})
\Leftrightarrow
s \in \text{WeakAss}_S(\mathcal{G})
\text{ and }
x \in \text{Ass}_{X_s}(\mathcal{F}_s).
$$
\end{lemma}

\begin{proof}
In this paragraph we reduce to $f$ being of finite presentation.
The question is local on $X$ and $S$, hence we may assume $X$ and $S$
are affine. Write $X = \Spec(B)$, $S = \Spec(A)$ and write
$B = A[x_1, \ldots, x_n]/I$. In other words we obtain a closed immersion
$i : X \to \mathbf{A}^n_S$ over $S$. Denote $t = i(x) \in \mathbf{A}^n_S$.
Note that $i_*\mathcal{F}$ is a finite type quasi-coherent sheaf on
$\mathbf{A}^n_S$ which is flat at $t$ over $S$ and note that
$$
i_*(\mathcal{F} \otimes_{\mathcal{O}_X} f^*\mathcal{G}) =
i_*\mathcal{F} \otimes_{\mathcal{O}_{\mathbf{A}^n_S}} p^*\mathcal{G}
$$
where $p : \mathbf{A}^n_S \to S$ is the projection. Note that
$t$ is a weakly associated point of
$i_*(\mathcal{F} \otimes_{\mathcal{O}_X} f^*\mathcal{G})$
if and only if $x$ is a weakly associated point of
$\mathcal{F} \otimes_{\mathcal{O}_X} f^*\mathcal{G}$, see
Divisors, Lemma \href{https://stacks.math.columbia.edu/tag/05EZ}{lemma weakly associated finite (Tag 05EZ)}.
Similarly $x \in \text{Ass}_{X_s}(\mathcal{F}_s)$ if and only
if $t \in \text{Ass}_{\mathbf{A}^n_s}((i_*\mathcal{F})_s)$ (see
Algebra, Lemma \href{https://stacks.math.columbia.edu/tag/05BY}{lemma ass quotient ring (Tag 05BY)}).
Hence it suffices to prove the lemma in case $X = \mathbf{A}^n_S$.
Thus we may assume that $X \to S$ is of finite presentation.

\medskip\noindent
In this paragraph we reduce to $\mathcal{F}$ being of finite presentation
and flat over $S$.
Choose an elementary \'etale neighbourhood $e : (S', s') \to (S, s)$
and an open $V \subset X \times_S \Spec(\mathcal{O}_{S', s'})$
as in Proposition \ref{proposition-finite-type-flat-at-point}.
Let $x' \in X' = X \times_S S'$ be the unique point mapping to $x$
and $s'$. Then it suffices to prove the statement for
$X' \to S'$, $x'$, $s'$, $(X' \to X)^*\mathcal{F}$, and $e^*\mathcal{G}$, see
Lemma \ref{lemma-etale-weak-assassin-up-down}.
Let $v \in V$ the unique point mapping to $x'$
and let $s' \in \Spec(\mathcal{O}_{S', s'})$ be the closed point.
Then $\mathcal{O}_{V, v} = \mathcal{O}_{X', x'}$
and $\mathcal{O}_{\Spec(\mathcal{O}_{S', s'}), s'} =
\mathcal{O}_{S', s'}$ and similarly for the stalks of pullbacks of
$\mathcal{F}$ and $\mathcal{G}$.
Also $V_{s'} \subset X'_{s'}$ is an open subscheme.
Since the condition of being a weakly associated point
depend only on the stalk of the sheaf, we may
replace
$X' \to S'$, $x'$, $s'$, $(X' \to X)^*\mathcal{F}$, and $e^*\mathcal{G}$
by
$V \to \Spec(\mathcal{O}_{S', s'})$, $v$, $s'$, $(V \to X)^*\mathcal{F}$,
and $(\Spec(\mathcal{O}_{S', s'}) \to S)^*\mathcal{G}$.
Thus we may assume that $f$ is of finite presentation and
$\mathcal{F}$ of finite presentation and flat over $S$.

\medskip\noindent
Assume $f$ is of finite presentation and
$\mathcal{F}$ of finite presentation and flat over $S$.
After shrinking $X$ and $S$ to affine neighbourhoods
of $x$ and $s$, this case is handled by
Lemma \ref{lemma-weak-bourbaki-pre}.
\end{proof}


\begin{lemma}
\label{lemma-bourbaki-finite-type-general-base}
Let $f : X \to S$ be a morphism which is locally of finite type.
Let $\mathcal{F}$ be a finite type quasi-coherent sheaf on $X$
which is flat over $S$. Let $\mathcal{G}$ be a quasi-coherent sheaf on $S$.
Then we have
$$
\text{WeakAss}_X(\mathcal{F} \otimes_{\mathcal{O}_X} f^*\mathcal{G}) =
\bigcup\nolimits_{s \in \text{WeakAss}_S(\mathcal{G})}
\text{Ass}_{X_s}(\mathcal{F}_s)
$$
\end{lemma}

\begin{proof}
Immediate consequence of
Lemma \ref{lemma-bourbaki-finite-type-general-base-at-point}.
\end{proof}


\begin{proposition}
\label{proposition-finite-presentation-flat-at-point}
Let $f : X \to S$ be a morphism of schemes.
Let $\mathcal{F}$ be a quasi-coherent sheaf on $X$.
Let $x \in X$ with image $s \in S$.
Assume that
\begin{enumerate}
\item $f$ is locally of finite presentation,
\item $\mathcal{F}$ is of finite presentation, and
\item $\mathcal{F}$ is flat at $x$ over $S$.
\end{enumerate}
Then there exists a commutative diagram of pointed schemes
$$
\xymatrix{
(X, x) \ar[d] & (X', x') \ar[l]^g \ar[d] \\
(S, s) & (S', s') \ar[l]
}
$$
whose horizontal arrows are elementary \'etale neighbourhoods
such that $X'$, $S'$ are affine and such that
$\Gamma(X', g^*\mathcal{F})$ is a projective
$\Gamma(S', \mathcal{O}_{S'})$-module.
\end{proposition}

\begin{proof}
By openness of flatness, see
More on Morphisms, Theorem \href{https://stacks.math.columbia.edu/tag/0399}{theorem openness flatness (Tag 0399)}
we may replace $X$ by an open neighbourhood of $x$ and assume that
$\mathcal{F}$ is flat over $S$. Next, we apply
Proposition \ref{proposition-existence-complete-at-x}
to find a diagram as in the statement of the proposition such
that $g^*\mathcal{F}/X'/S'$ has a complete d\'evissage over $s'$.
(In particular $S'$ and $X'$ are affine.) By
Morphisms, Lemma \href{https://stacks.math.columbia.edu/tag/02JZ}{lemma flat permanence (Tag 02JZ)}
we see that $g^*\mathcal{F}$ is flat over $S$ and by
Lemma \href{https://stacks.math.columbia.edu/tag/05B9}{lemma etale flat up down (Tag 05B9)}
we see that it is flat over $S'$. Via
Remark \ref{remark-same-notion}
we deduce that
$$
\Gamma(X', g^*\mathcal{F})/
\Gamma(X', \mathcal{O}_{X'})/
\Gamma(S', \mathcal{O}_{S'})
$$
has a complete d\'evissage over the prime of $\Gamma(S', \mathcal{O}_{S'})$
corresponding to $s'$. Thus
Lemma \ref{lemma-complete-devissage-flat-finitely-presented-module}
implies that the result of the proposition holds after replacing
$S'$ by a standard open neighbourhood of $s'$.
\end{proof}


\begin{proposition}
\label{proposition-existence-complete-at-x}
Let $S$ be a scheme.
Let $X$ be locally of finite type over $S$.
Let $x \in X$ be a point with image $s \in S$.
There exists a commutative diagram
$$
\xymatrix{
(X, x) \ar[d] & (X', x') \ar[l]^g \ar[d] \\
(S, s) & (S', s') \ar[l]
}
$$
of pointed schemes such that the horizontal
arrows are elementary \'etale neighbourhoods
and such that $g^*\mathcal{F}/X'/S'$ has a complete
d\'evissage at $x$.
\end{proposition}

\begin{proof}
We prove this by induction on the integer
$d = \dim_x(\text{Supp}(\mathcal{F}_s))$.
By
Lemma \ref{lemma-elementary-devissage-variant}
there exists a diagram
$$
\xymatrix{
(X, x) \ar[d] & (X', x') \ar[l]^g \ar[d] \\
(S, s) & (S', s') \ar[l]
}
$$
of pointed schemes such that the horizontal
arrows are elementary \'etale neighbourhoods
and such that $g^*\mathcal{F}/X'/S'$ has a one step d\'evissage at $x'$.
The local nature of the problem implies that we may replace
$(X, x) \to (S, s)$ by $(X', x') \to (S', s')$. Thus after doing so
we may assume that there exists a one step d\'evissage
$(Z_1, Y_1, i_1, \pi_1, \mathcal{G}_1)$ of $\mathcal{F}/X/S$ at $x$.

\medskip\noindent
We apply
Lemma \ref{lemma-existence-alpha}
to find a map
$$
\alpha_1 :
\mathcal{O}_{Y_1}^{\oplus r_1}
\longrightarrow
\pi_{1, *}\mathcal{G}_1
$$
which induces an isomorphism of vector spaces over $\kappa(\xi_1)$
where $\xi_1 \in Y_1$ is the unique generic point of the fibre of
$Y_1$ over $s$. Moreover
$\dim_{y_1}(\text{Supp}(\Coker(\alpha_1)_s)) < d$.
It may happen that the stalk of $\Coker(\alpha_1)_s$
at $y_1$ is zero. In this case we may shrink $Y_1$ by
Lemma \ref{lemma-shrink} (\ref{item-shrink-on-Y})
and assume that $\Coker(\alpha_1) = 0$ so we obtain a
complete d\'evissage of length zero.

\medskip\noindent
Assume now that the stalk of $\Coker(\alpha_1)_s$
at $y_1$ is not zero. In this case, by induction, there exists a
commutative diagram
\begin{equation}
\label{equation-overcome-this}
\vcenter{
\xymatrix{
(Y_1, y_1) \ar[d] & (Y'_1, y'_1) \ar[l]^h \ar[d] \\
(S, s) & (S', s') \ar[l]
}
}
\end{equation}
of pointed schemes such that the horizontal
arrows are elementary \'etale neighbourhoods
and such that $h^*\Coker(\alpha_1)/Y'_1/S'$ has a complete
d\'evissage
$$
(Z_k, Y_k, i_k, \pi_k, \mathcal{G}_k, \alpha_k, z_k, y_k)_{k = 2, \ldots, n}
$$
at $y'_1$. (In particular $i_2 : Z_2 \to Y'_1$ is a closed immersion into
$Y'_2$.) At this point we apply
Lemma \ref{lemma-elementary-etale-neighbourhood}
to $S, X, \mathcal{F}, x, s$, the system
$(Z_1, Y_1, i_1, \pi_1, \mathcal{G}_1)$ and
diagram (\ref{equation-overcome-this}). We obtain a diagram
$$
\xymatrix{
& & (X'', x'') \ar[lld] \ar[d] & (Z''_1, z''_1) \ar[l] \ar[lld] \ar[d] \\
(X, x) \ar[d] & (Z_1, z_1) \ar[l] \ar[d] &
(S'', s'') \ar[lld] & (Y''_1, y''_1) \ar[lld] \ar[l] \\
(S, s) & (Y_1, y_1) \ar[l]
}
$$
with all the properties as listed in the referenced lemma.
In particular $Y''_1 \subset Y'_1 \times_{S'} S''$. Set
$X_1 = Y'_1 \times_{S'} S''$ and let $\mathcal{F}_1$ denote the
pullback of $\Coker(\alpha_1)$. By
Lemma \ref{lemma-base-change-complete-at-x}
the system
\begin{equation}
\label{equation-shrink-this}
(Z_k \times_{S'} S'',
Y_k \times_{S'} S'', i''_k, \pi''_k, \mathcal{G}''_k,
\alpha''_k, z''_k, y''_k)_{k = 2, \ldots, n}
\end{equation}
is a complete d\'evissage of $\mathcal{F}_1$
to $X_1$. Again, the nature of the problem allows
us to replace $(X, x) \to (S, s)$ by $(X'', x'') \to (S'', s'')$.
In this we see that we may assume:
\begin{enumerate}
\item[(a)] There exists a one step d\'evissage
$(Z_1, Y_1, i_1, \pi_1, \mathcal{G}_1)$ of $\mathcal{F}/X/S$ at $x$,
\item[(b)] there exists an $\alpha_1 : \mathcal{O}_{Y_1}^{\oplus r_1}
\to \pi_{1, *}\mathcal{G}_1$ such that $\alpha \otimes \kappa(\xi_1)$
is an isomorphism,
\item[(c)] $Y_1 \subset X_1$ is open, $y_1 = x_1$, and
$\mathcal{F}_1|_{Y_1} \cong \Coker(\alpha_1)$, and
\item[(d)] there exists a complete d\'evissage
$(Z_k, Y_k, i_k, \pi_k, \mathcal{G}_k, \alpha_k, z_k, y_k)_{k = 2, \ldots, n}$
of $\mathcal{F}_1/X_1/S$ at $x_1$.
\end{enumerate}
To finish the proof all we have to do is shrink the one step d\'evissage
and the complete d\'evissage such that they fit together to a complete
d\'evissage. (We suggest the reader do this on their own using
Lemmas \ref{lemma-shrink} and
\ref{lemma-shrink-complete}
instead of reading the proof that follows.) Since $Y_1 \subset X_1$
is an open neighbourhood of $x_1$ we may apply
Lemma \ref{lemma-shrink-complete} (\ref{item-shrink-on-X-complete})
to find a standard shrinking $S', X'_1, Z'_2, Y'_2, \ldots, Y'_n$
of the datum (d) so that $X'_1 \subset Y_1$. Note that $X'_1$ is also
a standard open of the affine scheme $Y_1$. Next, we shrink the datum
(a) as follows: first we shrink the base $S$ to $S'$, see
Lemma \ref{lemma-shrink} (\ref{item-shrink-base}) and then
we shrink the result to $S''$, $X''$, $Z''_1$, $Y''_1$ using
Lemma \ref{lemma-shrink} (\ref{item-shrink-on-Y})
such that eventually $Y''_1 = X'_1 \times_S S''$ and $S'' \subset S'$.
Then we see that
$$
Z''_1, Y''_1, Z'_2 \times_{S'} S'', Y'_2 \times_{S'} S'', \ldots,
Y'_n \times_{S'} S''
$$
gives the complete d\'evissage we were looking for.
\end{proof}


\begin{lemma}
\label{lemma-existence-alpha}
Let $S$, $X$, $\mathcal{F}$, $s$ be as in
Definition \ref{definition-one-step-devissage}.
Let $(Z, Y, i, \pi, \mathcal{G})$ be a one step d\'evissage
of $\mathcal{F}/X/S$ over $s$.
Let $\xi \in Y_s$ be the (unique) generic point.
Then there exists an integer $r > 0$ and an $\mathcal{O}_Y$-module map
$$
\alpha : \mathcal{O}_Y^{\oplus r} \longrightarrow \pi_*\mathcal{G}
$$
such that
$$
\alpha :
\kappa(\xi)^{\oplus r}
\longrightarrow
(\pi_*\mathcal{G})_\xi \otimes_{\mathcal{O}_{Y, \xi}} \kappa(\xi)
$$
is an isomorphism. Moreover, in this case we have
$$
\dim(\text{Supp}(\Coker(\alpha)_s)) < \dim(\text{Supp}(\mathcal{F}_s)).
$$
\end{lemma}

\begin{proof}
By assumption the schemes $S$ and $Y$ are affine.
Write $S = \Spec(A)$ and $Y = \Spec(B)$.
As $\pi$ is finite the $\mathcal{O}_Y$-module $\pi_*\mathcal{G}$
is a finite type quasi-coherent $\mathcal{O}_Y$-module.
Hence $\pi_*\mathcal{G} = \widetilde{N}$ for some finite $B$-module $N$.
Let $\mathfrak p \subset B$ be the prime ideal corresponding to $\xi$.
To obtain $\alpha$ set
$r = \dim_{\kappa(\mathfrak p)} N \otimes_B \kappa(\mathfrak p)$
and pick $x_1, \ldots, x_r \in N$ which form a basis of
$N \otimes_B \kappa(\mathfrak p)$. Take $\alpha : B^{\oplus r} \to N$
to be the map given by the formula $\alpha(b_1, \ldots, b_r) = \sum b_ix_i$.
It is clear that
$\alpha : \kappa(\mathfrak p)^{\oplus r} \to N \otimes_B \kappa(\mathfrak p)$
is an isomorphism as desired. Finally, suppose $\alpha$ is any map with this
property. Then $N' = \Coker(\alpha)$ is a finite $B$-module
such that $N' \otimes \kappa(\mathfrak p) = 0$. By Nakayama's lemma
(Algebra, Lemma \href{https://stacks.math.columbia.edu/tag/00DV}{lemma NAK (Tag 00DV)})
we see that $N'_{\mathfrak p} = 0$. Since the fibre $Y_s$ is
geometrically irreducible of dimension $n$ with generic point $\xi$
and since we have just seen that $\xi$ is not in the support of
$\Coker(\alpha)$ the last assertion of the lemma holds.
\end{proof}


\begin{lemma}
\label{lemma-elementary-etale-neighbourhood}
Let $S$, $X$, $\mathcal{F}$, $x$, $s$ be as in
Definition \ref{definition-one-step-devissage-at-x}.
Let $(Z, Y, i, \pi, \mathcal{G}, z, y)$ be a one step d\'evissage
of $\mathcal{F}/X/S$ at $x$. Let
$$
\xymatrix{
(Y, y) \ar[d] & (Y', y') \ar[l] \ar[d] \\
(S, s) & (S', s') \ar[l]
}
$$
be a commutative diagram of pointed schemes such that the horizontal
arrows are elementary \'etale neighbourhoods. Then there exists
a commutative diagram
$$
\xymatrix{
& & (X'', x'') \ar[lld] \ar[d] & (Z'', z'') \ar[l] \ar[lld] \ar[d] \\
(X, x) \ar[d] & (Z, z) \ar[l] \ar[d] &
(S'', s'') \ar[lld] & (Y'', y'') \ar[lld] \ar[l] \\
(S, s) & (Y, y) \ar[l]
}
$$
of pointed schemes with the following properties:
\begin{enumerate}
\item $(S'', s'') \to (S', s')$ is an elementary \'etale neighbourhood and
the morphism $S'' \to S$ is the composition $S'' \to S' \to S$,
\item $Y''$ is an open subscheme of $Y' \times_{S'} S''$,
\item $Z'' = Z \times_Y Y''$,
\item $(X'', x'') \to (X, x)$ is an elementary \'etale neighbourhood, and
\item $(Z'', Y'', i'', \pi'', \mathcal{G}'', z'', y'')$ is a one step
d\'evissage at $x''$ of the sheaf $\mathcal{F}''$.
\end{enumerate}
Here $\mathcal{F}''$ (resp.\ $\mathcal{G}''$) is the pullback of
$\mathcal{F}$ (resp.\ $\mathcal{G}$) via the morphism $X'' \to X$
(resp.\ $Z'' \to Z$) and $i'' : Z'' \to X''$ and $\pi'' : Z'' \to Y''$
are as in the diagram.
\end{lemma}

\begin{proof}
Let $(S'', s'') \to (S', s')$ be any elementary \'etale neighbourhood
with $S''$ affine. Let $Y'' \subset Y' \times_{S'} S''$ be any affine
open neighbourhood containing the point $y'' = (y', s'')$. Then we
obtain an affine $(Z'', z'')$ by (3). Moreover $Z_{S''} \to X_{S''}$
is a closed immersion and $Z'' \to Z_{S''}$ is an \'etale
morphism. Hence
Lemma \href{https://stacks.math.columbia.edu/tag/057R}{lemma lift etale (Tag 057R)}
applies and we can find an \'etale morphism $X'' \to X_{S'}$ of affines
such that $Z'' \cong X'' \times_{X_{S'}} Z_{S'}$. Denote $i'' : Z'' \to X''$
the corresponding closed immersion. Setting $x'' = i''(z'')$ we obtain a
commutative diagram as in the lemma.
Properties (1), (2), (3), and (4) hold by construction.
Thus it suffices to show that (5) holds for a suitable choice of
$(S'', s'') \to (S', s')$ and $Y''$.

\medskip\noindent
We first list those properties which hold for any choice of
$(S'', s'') \to (S', s')$ and $Y''$ as in the first paragraph.
As we have $Z'' = X'' \times_X Z$ by construction we see that
$i''_*\mathcal{G}'' = \mathcal{F}''$ (with notation as in the
statement of the lemma), see
Cohomology of Schemes, Lemma \href{https://stacks.math.columbia.edu/tag/02KG}{lemma affine base change (Tag 02KG)}.
Set $n = \dim(\text{Supp}(\mathcal{F}_s)) = \dim_x(\text{Supp}(\mathcal{F}_s))$.
The morphism $Y'' \to S''$ is smooth of relative dimension $n$
(because $Y' \to S'$ is smooth of relative dimension $n$
as the composition $Y' \to Y_{S'} \to S'$ of an \'etale and
smooth morphism of relative dimension $n$ and because base change
preserves smooth morphisms of relative dimension $n$).
We have $\kappa(y'') = \kappa(y)$ and $\kappa(s) = \kappa(s'')$
hence $\kappa(y'')$ is a purely transcendental extension of $\kappa(s'')$.
The morphism of fibres $X''_{s''} \to X_s$ is an \'etale morphism of affine
schemes over $\kappa(s) = \kappa(s'')$ mapping the point $x''$ to the
point $x$ and pulling back $\mathcal{F}_s$ to $\mathcal{F}''_{s''}$.
Hence
$$
\dim(\text{Supp}(\mathcal{F}''_{s''})) =
\dim(\text{Supp}(\mathcal{F}_s)) = n =
\dim_x(\text{Supp}(\mathcal{F}_s)) =
\dim_{x''}(\text{Supp}(\mathcal{F}''_{s''}))
$$
because dimension is invariant under \'etale localization, see
Descent, Lemma \href{https://stacks.math.columbia.edu/tag/04N4}{lemma dimension at point local (Tag 04N4)}.
As $\pi'' : Z'' \to Y''$ is the base change of $\pi$ we see that
$\pi''$ is finite and as $\kappa(y) = \kappa(y'')$ we see that
$\pi^{-1}(\{y''\}) = \{z''\}$.

\medskip\noindent
At this point we have verified all the conditions of
Definition \ref{definition-one-step-devissage}
except we have not verified that $Y'' \to S''$ has geometrically
irreducible fibres. Of course in general this is not going to be
true, and it is at this point that we will use that
$\kappa(s) \subset \kappa(y)$ is purely transcendental. Namely,
let $T \subset Y'_{s'}$ be the irreducible component of
$Y'_{s'}$ containing $y' = (y, s')$. Note that $T$ is an open subscheme
of $Y'_{s'}$ as this is a smooth scheme over $\kappa(s')$. By
Varieties,
Lemma \href{https://stacks.math.columbia.edu/tag/04KV}{lemma geometrically connected if connected and point (Tag 04KV)}
we see that $T$ is geometrically connected because $\kappa(s') = \kappa(s)$
is algebraically closed in $\kappa(y') = \kappa(y)$.
As $T$ is smooth we see that $T$ is geometrically irreducible. Hence
More on Morphisms,
Lemma \href{https://stacks.math.columbia.edu/tag/057J}{lemma normal morphism irreducible (Tag 057J)}
applies and we can find an elementary \'etale morphism
$(S'', s'') \to (S', s')$ and an affine open $Y'' \subset Y'_{S''}$
such that all fibres of $Y'' \to S''$ are geometrically irreducible
and such that $T = Y''_{s''}$. After shrinking (first $Y''$ and then $S''$)
we may assume that both $Y''$ and $S''$ are affine.
This finishes the proof of the lemma.
\end{proof}


\begin{lemma}
\label{lemma-shrink}
With assumption and notation as in
Definition \ref{definition-shrink}
we have:
\begin{enumerate}
\item
\label{item-shrink-base}
If $S' \subset S$ is a standard open neighbourhood of $s$, then
setting $X' = X_{S'}$, $Z' = Z_{S'}$ and $Y' = Y_{S'}$ we obtain a
standard shrinking.
\item
\label{item-shrink-on-Y}
Let $W \subset Y$ be a standard open neighbourhood of $y$.
Then there exists a standard shrinking with $Y' = W \times_S S'$.
\item
\label{item-shrink-on-X}
Let $U \subset X$ be an open neighbourhood of $x$.
Then there exists a standard shrinking with $X' \subset U$.
\end{enumerate}
\end{lemma}

\begin{proof}
Part (1) is immediate from
Lemma \ref{lemma-base-change-one-step-at-x}
and the fact that the inverse image of a standard open under a morphism
of affine schemes is a standard open, see
Algebra, Lemma \href{https://stacks.math.columbia.edu/tag/00E2}{lemma spec functorial (Tag 00E2)}.

\medskip\noindent
Let $W \subset Y$ as in (2). Because $Y \to S$ is smooth it is open, see
Morphisms, Lemma \href{https://stacks.math.columbia.edu/tag/056G}{lemma smooth open (Tag 056G)}.
Hence we can find a standard open neighbourhood $S'$ of $s$
contained in the image of $W$. Then the fibres of $W_{S'} \to S'$
are nonempty open subschemes of the fibres of $Y \to S$ over $S'$
and hence geometrically irreducible too. Setting $Y' = W_{S'}$
and $Z' = \pi^{-1}(Y')$ we see that $Z' \subset Z$ is a standard open
neighbourhood of $z$. Let $\overline{h} \in \Gamma(Z, \mathcal{O}_Z)$
be a function such that $Z' = D(\overline{h})$. As $i : Z \to X$
is a closed immersion, we can find a function $h \in \Gamma(X, \mathcal{O}_X)$
such that $i^\sharp(h) = \overline{h}$. Take $X' = D(h) \subset X$.
In this way we obtain a standard shrinking as in (2).

\medskip\noindent
Let $U \subset X$ be as in (3). We may after shrinking $U$ assume that
$U$ is a standard open. By
More on Morphisms,
Lemma \ref{more-morphisms-lemma-finite-morphism-single-point-in-fibre}
there exists a standard open $W \subset Y$ neighbourhood of $y$ such
that $\pi^{-1}(W) \subset i^{-1}(U)$. Apply (2) to get a standard
shrinking $X', S', Z', Y'$ with $Y' = W_{S'}$. Since
$Z' \subset \pi^{-1}(W) \subset i^{-1}(U)$ we may replace $X'$ by
$X' \cap U$ (still a standard open as $U$ is also standard open)
without violating any of the conditions defining a standard shrinking.
Hence we win.
\end{proof}


\begin{lemma}
\label{lemma-shrink-complete}
With assumption and notation as in
Definition \ref{definition-shrink-complete}
we have:
\begin{enumerate}
\item
\label{item-shrink-base-complete}
If $S' \subset S$ is a standard open neighbourhood of $s$, then
setting $X' = X_{S'}$, $Z'_k = Z_{S'}$ and $Y'_k = Y_{S'}$ we obtain a
standard shrinking.
\item
\label{item-shrink-on-Y-complete}
Let $W \subset Y_n$ be a standard open neighbourhood of $y$.
Then there exists a standard shrinking with $Y'_n = W \times_S S'$.
\item
\label{item-shrink-on-X-complete}
Let $U \subset X$ be an open neighbourhood of $x$.
Then there exists a standard shrinking with $X' \subset U$.
\end{enumerate}
\end{lemma}

\begin{proof}
Part (1) is immediate from
Lemmas \ref{lemma-base-change-complete-at-x} and
\ref{lemma-shrink}.

\medskip\noindent
Proof of (2). For convenience denote $X = Y_0$. We apply
Lemma \ref{lemma-shrink} (\ref{item-shrink-on-Y})
to find a standard shrinking
$S', Y'_{n - 1}, Z'_n, Y'_n$
of the one step d\'evissage of $\Coker(\alpha_{n - 1})/Y_{n - 1}/S$
at $y_{n - 1}$ with $Y'_n = W \times_S S'$. We may repeat this procedure
and find a standard shrinking
$S'', Y''_{n - 2}, Z''_{n - 1}, Y''_{n - 1}$
of the one step d\'evissage of $\Coker(\alpha_{n - 2})/Y_{n - 2}/S$
at $y_{n - 2}$ with $Y''_{n - 1} = Y'_{n - 1} \times_S S''$.
We may continue in this manner until we obtain
$S^{(n)}, Y^{(n)}_0, Z^{(n)}_1, Y^{(n)}_1$.
At this point it is clear that we obtain our desired standard shrinking
by taking $S^{(n)}$, $X^{(n)}$, $Z_k^{(n - k)} \times_S S^{(n)}$, and
$Y_k^{(n - k)} \times_S S^{(n)}$ with the desired property.

\medskip\noindent
Proof of (3). We use induction on the length of the complete
d\'evissage. First we apply
Lemma \ref{lemma-shrink} (\ref{item-shrink-on-X})
to find a standard shrinking
$S', X', Z'_1, Y'_1$
of the one step d\'evissage of $\mathcal{F}/X/S$ at $x$
with $X' \subset U$. If $n = 1$, then we are done.
If $n > 1$, then by induction we can find a standard shrinking
$S''$, $Y''_1$, $Z''_k$, and $Y''_k$ of the complete d\'evissage
$(Z_k, Y_k, i_k, \pi_k, \mathcal{G}_k, \alpha_k, z_k, y_k)_{k = 2, \ldots, n}$
of $\Coker(\alpha_1)/Y_1/S$ at $x$ such that
$Y''_1 \subset Y'_1$. Using
Lemma \ref{lemma-shrink} (\ref{item-shrink-on-Y})
we can find $S''' \subset S'$, $X''' \subset X'$, $Z'''_1$ and
$Y'''_1 = Y''_1 \times_S S'''$ which is a standard shrinking.
The solution to our problem is to take
$$
S''', X''', Z'''_1, Y'''_1, Z''_2 \times_S S''',
Y''_2 \times_S S''', \ldots, Z''_n \times_S S''', Y''_n \times_S S'''
$$
This ends the proof of the lemma.
\end{proof}


\begin{lemma}
\label{lemma-base-change-complete}
Let $S$, $X$, $\mathcal{F}$, $s$ be as in
Definition \ref{definition-complete-devissage}.
Let $(S', s') \to (S, s)$ be any morphism of pointed schemes.
Let $(Z_k, Y_k, i_k, \pi_k, \mathcal{G}_k, \alpha_k)_{k = 1, \ldots, n}$
be a complete d\'evissage of $\mathcal{F}/X/S$ over $s$.
Given this data let $X', Z'_k, Y'_k, i'_k, \pi'_k$ be the base
changes of $X, Z_k, Y_k, i_k, \pi_k$ via $S' \to S$.
Let $\mathcal{F}'$ be the pullback of $\mathcal{F}$ to $X'$
and let $\mathcal{G}'_k$ be the pullback of $\mathcal{G}_k$ to $Z'_k$.
Let $\alpha'_k$ be the pullback of $\alpha_k$ to $Y'_k$.
If $S'$ is affine, then
$(Z'_k, Y'_k, i'_k, \pi'_k, \mathcal{G}'_k, \alpha'_k)_{k = 1, \ldots, n}$
is a complete d\'evissage of $\mathcal{F}'/X'/S'$ over $s'$.
\end{lemma}

\begin{proof}
By
Lemma \ref{lemma-base-change-one-step}
we know that the base change of a one step d\'evissage is a one step
d\'evissage. Hence it suffices to prove that formation of
$\Coker(\alpha_k)$ commutes with base change and that
condition (2) of
Definition \ref{definition-complete-devissage}
is preserved by base change. The first is true as
$\pi'_{k, *}\mathcal{G}'_k$ is the pullback of
$\pi_{k, *}\mathcal{G}_k$ (by
Cohomology of Schemes, Lemma \href{https://stacks.math.columbia.edu/tag/02KG}{lemma affine base change (Tag 02KG)})
and because $\otimes$ is right exact. The second because
by the same token we have
$$
(\pi_{k, *}\mathcal{G}_k)_{\xi_k}
\otimes_{\mathcal{O}_{Y_k, \xi_k}} \kappa(\xi_k)
\otimes_{\kappa(\xi_k)} \kappa(\xi'_k)
\cong
(\pi'_{k, *}\mathcal{G}'_k)_{\xi'_k}
\otimes_{\mathcal{O}_{Y'_k, \xi'_k}} \kappa(\xi'_k)
$$
with obvious notation.
\end{proof}


\begin{lemma}
\label{lemma-base-change-complete-at-x}
Let $S$, $X$, $\mathcal{F}$, $x$, $s$ be as in
Definition \ref{definition-complete-devissage-at-x}.
Let $(S', s') \to (S, s)$ be a morphism of pointed schemes
which induces an isomorphism $\kappa(s) = \kappa(s')$. Let
$(Z_k, Y_k, i_k, \pi_k, \mathcal{G}_k, \alpha_k, z_k, y_k)_{k = 1, \ldots, n}$
be a complete d\'evissage of $\mathcal{F}/X/S$ at $x$.
Let
$(Z'_k, Y'_k, i'_k, \pi'_k, \mathcal{G}'_k, \alpha'_k)_{k = 1, \ldots, n}$
be as constructed in
Lemma \ref{lemma-base-change-complete}
and let $x' \in X'$ (resp.\ $z'_k \in Z'$, $y'_k \in Y'$) be the
unique point mapping to both $x \in X$ (resp.\ $z_k \in Z_k$, $y_k \in Y_k$)
and $s' \in S'$.
If $S'$ is affine, then
$(Z'_k, Y'_k, i'_k, \pi'_k, \mathcal{G}'_k, \alpha'_k,
z'_k, y'_k)_{k = 1, \ldots, n}$
is a complete d\'evissage of $\mathcal{F}'/X'/S'$ at $x'$.
\end{lemma}

\begin{proof}
Combine
Lemma \ref{lemma-base-change-complete}
and
Lemma \ref{lemma-base-change-one-step-at-x}.
\end{proof}


\begin{lemma}
\label{lemma-base-change-one-step}
Let $S$, $X$, $\mathcal{F}$, $s$ be as in
Definition \ref{definition-one-step-devissage}.
Let $(Z, Y, i, \pi, \mathcal{G})$ be a one step d\'evissage
of $\mathcal{F}/X/S$ over $s$.
Let $(S', s') \to (S, s)$ be any morphism of pointed schemes.
Given this data let $X', Z', Y', i', \pi'$ be the base
changes of $X, Z, Y, i, \pi$ via $S' \to S$.
Let $\mathcal{F}'$ be the pullback of $\mathcal{F}$ to $X'$
and let $\mathcal{G}'$ be the pullback of $\mathcal{G}$ to $Z'$.
If $S'$ is affine, then $(Z', Y', i', \pi', \mathcal{G}')$
is a one step d\'evissage of $\mathcal{F}'/X'/S'$ over $s'$.
\end{lemma}

\begin{proof}
Fibre products of affines are affine, see
Schemes, Lemma \href{https://stacks.math.columbia.edu/tag/01JQ}{lemma fibre product affines (Tag 01JQ)}.
Base change preserves
closed immersions,
morphisms of finite presentation,
finite morphisms,
smooth morphisms,
morphisms with geometrically irreducible fibres, and
morphisms of relative dimension $n$, see
Morphisms, Lemmas \href{https://stacks.math.columbia.edu/tag/01QR}{lemma base change closed immersion (Tag 01QR)},
\href{https://stacks.math.columbia.edu/tag/01TS}{lemma base change finite presentation (Tag 01TS)},
\href{https://stacks.math.columbia.edu/tag/01WL}{lemma base change finite (Tag 01WL)},
\href{https://stacks.math.columbia.edu/tag/01VB}{lemma base change smooth (Tag 01VB)},
\href{https://stacks.math.columbia.edu/tag/02NK}{lemma base change relative dimension d (Tag 02NK)}, and
More on Morphisms, Lemma
\href{https://stacks.math.columbia.edu/tag/0555}{lemma base change fibres geometrically irreducible (Tag 0555)}.
We have $i'_*\mathcal{G}' \cong \mathcal{F}'$ because pushforward
along the finite morphism $i$ commutes with base change, see
Cohomology of Schemes, Lemma \href{https://stacks.math.columbia.edu/tag/02KG}{lemma affine base change (Tag 02KG)}.
We have
$\dim(\text{Supp}(\mathcal{F}_s)) = \dim(\text{Supp}(\mathcal{F}'_{s'}))$
by
Morphisms, Lemma \href{https://stacks.math.columbia.edu/tag/02FY}{lemma dimension fibre after base change (Tag 02FY)}
because
$$
\text{Supp}(\mathcal{F}_s) \times_s s' = \text{Supp}(\mathcal{F}'_{s'}).
$$
This proves the lemma.
\end{proof}


\begin{lemma}
\label{lemma-base-change-one-step-at-x}
Let $S$, $X$, $\mathcal{F}$, $x$, $s$ be as in
Definition \ref{definition-one-step-devissage-at-x}.
Let $(Z, Y, i, \pi, \mathcal{G}, z, y)$ be a one step d\'evissage
of $\mathcal{F}/X/S$ at $x$.
Let $(S', s') \to (S, s)$ be a morphism of pointed schemes
which induces an isomorphism $\kappa(s) = \kappa(s')$.
Let $(Z', Y', i', \pi', \mathcal{G}')$ be as constructed in
Lemma \ref{lemma-base-change-one-step}
and let $x' \in X'$ (resp.\ $z' \in Z'$, $y' \in Y'$) be the
unique point mapping to both $x \in X$ (resp.\ $z \in Z$, $y \in Y$)
and $s' \in S'$.
If $S'$ is affine, then $(Z', Y', i', \pi', \mathcal{G}', z', y')$
is a one step d\'evissage of $\mathcal{F}'/X'/S'$ at $x'$.
\end{lemma}

\begin{proof}
By
Lemma \ref{lemma-base-change-one-step}
$(Z', Y', i', \pi', \mathcal{G}')$ is a one step d\'evissage of
$\mathcal{F}'/X'/S'$ over $s'$. Properties (1) -- (4) of
Definition \ref{definition-one-step-devissage-at-x}
hold for $(Z', Y', i', \pi', \mathcal{G}', z', y')$
as the assumption that $\kappa(s) = \kappa(s')$ insures that the fibres
$X'_{s'}$, $Z'_{s'}$, and $Y'_{s'}$ are isomorphic to
$X_s$, $Z_s$, and $Y_s$.
\end{proof}


\begin{lemma}
\label{lemma-complete-devissage-flat-finitely-presented-module}
Let $R \to S$ be a ring map of finite presentation.
Let $N$ be a finitely presented $S$-module flat over $R$.
Let $\mathfrak r \subset R$ be a prime ideal.
Assume there exists a complete d\'evissage of $N/S/R$ over $\mathfrak r$.
Then there exists an $f \in R$, $f \not \in \mathfrak r$
such that
$$
N_f \cong B_1^{\oplus r_1} \oplus \ldots \oplus B_n^{\oplus r_n}
$$
as $R$-modules where each $B_i$ is a smooth $R_f$-algebra with geometrically
irreducible fibres. Moreover, $N_f$ is projective as an $R_f$-module.
\end{lemma}

\begin{proof}
Let $(A_i, B_i, M_i, \alpha_i)_{i = 1, \ldots, n}$ be the given
complete d\'evissage. We prove the lemma by induction on $n$.
Note that the assertions of the lemma are entirely about the structure
of $N$ as an $R$-module. Hence we may replace $N$ by $M_1$, and we
may think of $M_1$ as a $B_1$-module. See
Remark \ref{remark-finite-presentation}
in order to see why $M_1$ is of finite presentation as a $B_1$-module. By
Lemma \ref{lemma-induction-step-fp}
we may, after replacing $R$ by $R_f$ for some
$f \in R$, $f \not \in \mathfrak r$, assume
the map $\alpha_1 : B_1^{\oplus r_1} \to M_1$ is $R$-universally injective.
Since $M_1$ and $B_1^{\oplus r_1}$ are $R$-flat and finitely presented as
$B_1$-modules we see that $\Coker(\alpha_1)$ is $R$-flat
(Algebra, Lemma \href{https://stacks.math.columbia.edu/tag/058P}{lemma ui flat domain (Tag 058P)})
and finitely presented as a $B_1$-module. Note that
$(A_i, B_i, M_i, \alpha_i)_{i = 2, \ldots, n}$ is a complete
d\'evissage of $\Coker(\alpha_1)$. Hence the induction hypothesis
implies that, after replacing
$R$ by $R_f$ for some $f \in R$, $f \not \in \mathfrak r$,
we may assume that $\Coker(\alpha_1)$ has a decomposition
as in the lemma and is projective. In particular
$M_1 = B_1^{\oplus r_1} \oplus \Coker(\alpha_1)$.
This proves the statement regarding the decomposition.
The statement on projectivity follows as $B_1$ is projective as
an $R$-module by
Lemma \ref{lemma-fibres-irreducible-flat-projective-nonnoetherian}.
\end{proof}


\begin{lemma}
\label{lemma-induction-step-fp}
Let $R$ be a ring. Let $R \to S$ be a finitely presented
flat ring map with geometrically integral fibres. Let
$\mathfrak q \subset S$ be a prime ideal lying over the prime
$\mathfrak r \subset R$. Set $\mathfrak p = \mathfrak r S$.
Let $N$ be a finitely presented $S$-module.
There exists $r \geq 0$ and an $S$-module map
$$
\alpha : S^{\oplus r} \longrightarrow N
$$
such that
$\alpha : \kappa(\mathfrak p)^{\oplus r} \to N \otimes_S \kappa(\mathfrak p)$
is an isomorphism. For any such $\alpha$ the following are equivalent:
\begin{enumerate}
\item $N_{\mathfrak q}$ is $R$-flat,
\item there exists an $f \in R$, $f \not \in \mathfrak r$ such that
$\alpha_f : S_f^{\oplus r} \to N_f$ is $R_f$-universally injective and
a $g \in S$, $g \not \in \mathfrak q$ such that $\Coker(\alpha)_g$
is $R$-flat,
\item $\alpha_{\mathfrak r}$ is $R_{\mathfrak r}$-universally injective and
$\Coker(\alpha)_{\mathfrak q}$ is $R$-flat
\item $\alpha_{\mathfrak r}$ is injective and
$\Coker(\alpha)_{\mathfrak q}$ is $R$-flat,
\item $\alpha_{\mathfrak p}$ is an isomorphism and
$\Coker(\alpha)_{\mathfrak q}$ is $R$-flat, and
\item $\alpha_{\mathfrak q}$ is injective and
$\Coker(\alpha)_{\mathfrak q}$ is $R$-flat.
\end{enumerate}
\end{lemma}

\begin{proof}
To obtain $\alpha$ set
$r = \dim_{\kappa(\mathfrak p)} N \otimes_S \kappa(\mathfrak p)$ and pick
$x_1, \ldots, x_r \in N$ which form a basis of
$N \otimes_S \kappa(\mathfrak p)$. Define
$\alpha(s_1, \ldots, s_r) = \sum s_i x_i$. This proves the existence.

\medskip\noindent
Fix a choice of $\alpha$.
We may apply
Lemma \ref{lemma-induction-step}
to the map
$\alpha_{\mathfrak r} : S_{\mathfrak r}^{\oplus r} \to N_{\mathfrak r}$.
Hence we see that (1), (3), (4), (5), and (6) are all equivalent.
Since it is also clear that (2) implies (3) we see that all we have to
do is show that (1) implies (2).

\medskip\noindent
Assume (1). By openness of flatness, see
Algebra, Theorem \href{https://stacks.math.columbia.edu/tag/00RC}{theorem openness flatness (Tag 00RC)},
the set
$$
U_1 = \{\mathfrak q' \subset S \mid N_{\mathfrak q'}\text{ is flat over }R\}
$$
is open in $\Spec(S)$. It contains $\mathfrak q$ by assumption
and hence $\mathfrak p$. Because $S^{\oplus r}$ and $N$ are finitely presented
$S$-modules the set
$$
U_2 = \{\mathfrak q' \subset S \mid
\alpha_{\mathfrak q'}\text{ is an isomorphism}\}
$$
is open in $\Spec(S)$, see
Algebra, Lemma \href{https://stacks.math.columbia.edu/tag/05GF}{lemma map between finitely presented (Tag 05GF)}.
It contains $\mathfrak p$ by (5). As $R \to S$
is finitely presented and flat the map
$\Phi : \Spec(S) \to \Spec(R)$ is open, see
Algebra, Proposition \href{https://stacks.math.columbia.edu/tag/00I1}{proposition fppf open (Tag 00I1)}.
For any prime $\mathfrak r' \in \Phi(U_1 \cap U_2)$ we see that
there exists a prime $\mathfrak q'$ lying over $\mathfrak r'$ such that
$N_{\mathfrak q'}$ is flat and such that $\alpha_{\mathfrak q'}$ is
an isomorphism, which implies that $\alpha \otimes \kappa(\mathfrak p')$
is an isomorphism where $\mathfrak p' = \mathfrak r' S$. Thus
$\alpha_{\mathfrak r'}$ is $R_{\mathfrak r'}$-universally injective
by the implication (1) $\Rightarrow$ (3).
Hence if we pick $f \in R$, $f \not \in \mathfrak r$ such that
$D(f) \subset \Phi(U_1 \cap U_2)$ then we conclude that
$\alpha_f$ is $R_f$-universally injective, see
Algebra, Lemma \href{https://stacks.math.columbia.edu/tag/05CL}{lemma universally injective check stalks (Tag 05CL)}.
The same reasoning also shows that for any
$\mathfrak q' \in U_1 \cap \Phi^{-1}(\Phi(U_1 \cap U_2))$
the module $\Coker(\alpha)_{\mathfrak q'}$ is $R$-flat.
Note that $\mathfrak q \in U_1 \cap \Phi^{-1}(\Phi(U_1 \cap U_2))$.
Hence we can find a $g \in S$, $g \not \in \mathfrak q$ such
that $D(g) \subset U_1 \cap \Phi^{-1}(\Phi(U_1 \cap U_2))$
and we win.
\end{proof}


\begin{lemma}
\label{lemma-base-change-universally-flat-local}
Let $R \to S$ be a ring map.
Let $N$ be an $S$-module.
Let $S \to S'$ be a ring map.
Assume
\begin{enumerate}
\item $R \to S$ is a local homomorphism of local rings
\item $S$ is essentially of finite presentation over $R$,
\item $N$ is of finite presentation over $S$,
\item $N$ is flat over $R$,
\item $S \to S'$ is flat, and
\item the image of $\Spec(S') \to \Spec(S)$ contains
all primes $\mathfrak q$ of $S$ lying over $\mathfrak m_R$
such that $\mathfrak q$ is an associated prime of $N/\mathfrak m_R N$.
\end{enumerate}
Then $N \to N \otimes_S S'$ is $R$-universally injective.
\end{lemma}

\begin{proof}
Set $N' = N \otimes_R S'$. Consider the commutative diagram
$$
\xymatrix{
N \ar[d] \ar[r] & N' \ar[d] \\
\Sigma^{-1}N \ar[r] & \Sigma^{-1}N'
}
$$
where $\Sigma \subset S$ is the set of elements which are not a
zerodivisor on $N/\mathfrak m_R N$. If we can show that the map
$N \to \Sigma^{-1}N'$ is universally injective, then $N \to N'$
is too (see
Algebra, Lemma \href{https://stacks.math.columbia.edu/tag/05CJ}{lemma universally injective permanence (Tag 05CJ)}).

\medskip\noindent
By
Lemma \ref{lemma-homothety-spectrum}
the ring $\Sigma^{-1}S$ is a semi-local ring whose maximal ideals
correspond to associated primes of $N/\mathfrak m_R N$.
Hence the image of
$\Spec(\Sigma^{-1}S') \to \Spec(\Sigma^{-1}S)$
contains all these maximal ideals by assumption. By
Algebra, Lemma \href{https://stacks.math.columbia.edu/tag/00HQ}{lemma ff rings (Tag 00HQ)}
the ring map $\Sigma^{-1}S \to \Sigma^{-1}S'$ is faithfully flat.
Hence $\Sigma^{-1}N \to \Sigma^{-1}N'$, which is the map
$$
N \otimes_S \Sigma^{-1}S \longrightarrow N \otimes_S \Sigma^{-1}S'
$$
is universally injective, see
Algebra, Lemmas \href{https://stacks.math.columbia.edu/tag/05CK}{lemma faithfully flat universally injective (Tag 05CK)} and
\href{https://stacks.math.columbia.edu/tag/05CH}{lemma universally injective tensor (Tag 05CH)}.
Finally, we apply
Lemma \ref{lemma-homothety-universally-injective}
to see that $N \to \Sigma^{-1}N$ is universally injective.
As the composition of universally injective module maps is universally
injective (see
Algebra, Lemma \href{https://stacks.math.columbia.edu/tag/05CI}{lemma composition universally injective (Tag 05CI)})
we conclude that $N \to \Sigma^{-1}N'$ is universally injective and we win.
\end{proof}


\begin{lemma}
\label{lemma-base-change-universally-flat}
Let $R \to S$ be a ring map.
Let $N$ be an $S$-module.
Let $S \to S'$ be a ring map.
Assume
\begin{enumerate}
\item $R \to S$ is of finite presentation and $N$ is of finite presentation
over $S$,
\item $N$ is flat over $R$,
\item $S \to S'$ is flat, and
\item the image of $\Spec(S') \to \Spec(S)$ contains
all primes $\mathfrak q$ such that $\mathfrak q$ is an associated prime
of $N \otimes_R \kappa(\mathfrak p)$ where $\mathfrak p$ is the inverse
image of $\mathfrak q$ in $R$.
\end{enumerate}
Then $N \to N \otimes_S S'$ is $R$-universally injective.
\end{lemma}

\begin{proof}
By
Algebra, Lemma \href{https://stacks.math.columbia.edu/tag/05CL}{lemma universally injective check stalks (Tag 05CL)}
it suffices to show that $N_{\mathfrak q} \to (N \otimes_R S')_{\mathfrak q}$
is a $R_{\mathfrak p}$-universally injective for any prime $\mathfrak q$
of $S$ lying over $\mathfrak p$ in $R$. Thus we may apply
Lemma \ref{lemma-base-change-universally-flat-local}
to the ring maps
$R_{\mathfrak p} \to S_{\mathfrak q} \to S'_{\mathfrak q}$
and the module $N_{\mathfrak q}$.
\end{proof}


\begin{lemma}
\label{lemma-homothety-spectrum}
Let $R \to S$ be a ring map.
Let $N$ be a $S$-module.
Assume
\begin{enumerate}
\item $R$ is a local ring with maximal ideal $\mathfrak m$,
\item $\overline{S} = S/\mathfrak m S$ is Noetherian, and
\item $\overline{N} = N/\mathfrak m_R N$ is a finite $\overline{S}$-module.
\end{enumerate}
Let $\Sigma \subset S$ be the multiplicative subset of elements which are not
a zerodivisor on $\overline{N}$. Then $\Sigma^{-1}S$ is a semi-local ring
whose spectrum consists of primes $\mathfrak q \subset S$ contained in an
element of $\text{Ass}_S(\overline{N})$. Moreover, any maximal
ideal of $\Sigma^{-1}S$ corresponds to an associated prime of
$\overline{N}$ over $\overline{S}$.
\end{lemma}

\begin{proof}
Note that
$\text{Ass}_S(\overline{N}) = \text{Ass}_{\overline{S}}(\overline{N})$, see
Algebra, Lemma \href{https://stacks.math.columbia.edu/tag/05BY}{lemma ass quotient ring (Tag 05BY)}.
This is a finite set by
Algebra, Lemma \href{https://stacks.math.columbia.edu/tag/00LC}{lemma finite ass (Tag 00LC)}.
Say $\{\mathfrak q_1, \ldots, \mathfrak q_r\} = \text{Ass}_S(\overline{N})$.
We have $\Sigma = S \setminus (\bigcup \mathfrak q_i)$ by
Algebra, Lemma \href{https://stacks.math.columbia.edu/tag/00LD}{lemma ass zero divisors (Tag 00LD)}.
By the description of $\Spec(\Sigma^{-1}S)$ in
Algebra, Lemma \href{https://stacks.math.columbia.edu/tag/00E3}{lemma spec localization (Tag 00E3)}
and by
Algebra, Lemma \href{https://stacks.math.columbia.edu/tag/00DS}{lemma silly (Tag 00DS)}
we see that the primes of $\Sigma^{-1}S$ correspond to the primes of
$S$ contained in one of the $\mathfrak q_i$.
Hence the maximal ideals of $\Sigma^{-1}S$ correspond one-to-one with the
maximal (w.r.t.\ inclusion) elements of the set
$\{\mathfrak q_1, \ldots, \mathfrak q_r\}$. This proves the lemma.
\end{proof}


\begin{lemma}
\label{lemma-homothety-universally-injective}
Assumption and notation as in
Lemma \ref{lemma-homothety-spectrum}.
Assume moreover that
\begin{enumerate}
\item $S$ is local and $R \to S$ is a local homomorphism,
\item $S$ is essentially of finite presentation over $R$,
\item $N$ is finitely presented over $S$, and
\item $N$ is flat over $R$.
\end{enumerate}
Then each $s \in \Sigma$ defines a
universally injective $R$-module map $s : N \to N$, and the
map $N \to \Sigma^{-1}N$ is $R$-universally injective.
\end{lemma}

\begin{proof}
By
Algebra, Lemma \href{https://stacks.math.columbia.edu/tag/046Y}{lemma mod injective general (Tag 046Y)}
the sequence $0 \to N \to N \to N/sN \to 0$ is exact and
$N/sN$ is flat over $R$. This implies that $s : N \to N$
is universally injective, see
Algebra, Lemma \href{https://stacks.math.columbia.edu/tag/00HL}{lemma flat tor zero (Tag 00HL)}.
The map $N \to \Sigma^{-1}N$ is universally injective as the directed
colimit of the maps $s : N \to N$.
\end{proof}


\begin{lemma}
\label{lemma-etale-weak-assassin-up-down}
Let $h : U \to S$ be an \'etale morphism of schemes.
Let $\mathcal{G}$ be a quasi-coherent $\mathcal{O}_S$-module.
Let $u \in U$ be a point with image $s \in S$. Then
$$
u \in \text{WeakAss}(h^*\mathcal{G})
\Leftrightarrow
s \in \text{WeakAss}(\mathcal{G})
$$
\end{lemma}

\begin{proof}
After replacing $S$ and $U$ by affine neighbourhoods of $s$ and $u$
we may assume that $g$ is a standard \'etale morphism of affines, see
Morphisms, Lemma \href{https://stacks.math.columbia.edu/tag/02GT}{lemma etale locally standard etale (Tag 02GT)}.
Thus we may assume $S = \Spec(A)$ and
$X = \Spec(A[x, 1/g]/(f))$, where $f$ is monic and $f'$
is invertible in $A[x, 1/g]$.
Note that $A[x, 1/g]/(f) = (A[x]/(f))_g$ is also the localization
of the finite free $A$-algebra $A[x]/(f)$. Hence we may think of
$U$ as an open subscheme of the scheme $T = \Spec(A[x]/(f))$
which is finite locally free over $S$. This reduces us to
Lemma \ref{lemma-finite-flat-weak-assassin-up-down}
above.
\end{proof}


\begin{lemma}
\label{lemma-finite-flat-weak-assassin-up-down}
Let $g : T \to S$ be a finite flat morphism of schemes.
Let $\mathcal{G}$ be a quasi-coherent $\mathcal{O}_S$-module.
Let $t \in T$ be a point with image $s \in S$. Then
$$
t \in \text{WeakAss}(g^*\mathcal{G})
\Leftrightarrow
s \in \text{WeakAss}(\mathcal{G})
$$
\end{lemma}

\begin{proof}
The implication ``$\Leftarrow$'' follows immediately from
Divisors, Lemma \href{https://stacks.math.columbia.edu/tag/05F0}{lemma weakly ass pullback (Tag 05F0)}.
Assume $t \in \text{WeakAss}(g^*\mathcal{G})$.
Let $\Spec(A) \subset S$ be an affine open neighbourhood of $s$.
Let $\mathcal{G}$ be the quasi-coherent sheaf associated to the $A$-module $M$.
Let $\mathfrak p \subset A$ be the prime ideal corresponding to $s$.
As $g$ is finite flat we have $g^{-1}(\Spec(A)) = \Spec(B)$
for some finite flat $A$-algebra $B$. Note that
$g^*\mathcal{G}$ is the quasi-coherent $\mathcal{O}_{\Spec(B)}$-module
associated to the $B$-module $M \otimes_A B$ and $g_*g^*\mathcal{G}$ is the
quasi-coherent $\mathcal{O}_{\Spec(A)}$-module associated to the
$A$-module $M \otimes_A B$. By
Algebra, Lemma \href{https://stacks.math.columbia.edu/tag/00NZ}{lemma finite flat local (Tag 00NZ)}
we have $B_{\mathfrak p} \cong A_{\mathfrak p}^{\oplus n}$
for some integer $n \geq 0$. Note that $n \geq 1$ as we assumed there
exists at least one point of $T$ lying over $s$. Hence we see by
looking at stalks that
$$
s \in \text{WeakAss}(\mathcal{G})
\Leftrightarrow
s \in \text{WeakAss}(g_*g^*\mathcal{G})
$$
Now the assumption that $t \in \text{WeakAss}(g^*\mathcal{G})$
implies that $s \in \text{WeakAss}(g_*g^*\mathcal{G})$ by
Divisors, Lemma \href{https://stacks.math.columbia.edu/tag/05EZ}{lemma weakly associated finite (Tag 05EZ)}
and hence by the above $s \in  \text{WeakAss}(\mathcal{G})$.
\end{proof}


\begin{lemma}
\label{more-algebra-lemma-P-fPD-zero}
Let $R$ be a ring. The following are equivalent
\begin{enumerate}
\item $R$ has property (P) of
Lemma \ref{more-algebra-lemma-auto-ass-implies-P},
\item any injective map of projective $R$-modules is
universally injective,
\item if $u : N \to M$ is injective and $N$, $M$ are finite projective
$R$-modules then $\Coker(u)$ is a finite projective $R$-module,
\item if $N \subset M$ and $N$, $M$ are finite projective as $R$-modules, then
$N$ is a direct summand of $M$, and
\item any injective map $R \to R^{\oplus n}$ is a split injection.
\end{enumerate}
\end{lemma}

\begin{proof}
The implication (1) $\Rightarrow$ (2) is
Lemma \ref{more-algebra-lemma-P-universally-injective}.
It is clear that (3) and (4) are equivalent.
We have (2) $\Rightarrow$ (3), (4) by
Algebra, Lemma \href{https://stacks.math.columbia.edu/tag/058L}{lemma universally exact split (Tag 058L)}.
Part (5) is a special case of (4).
Assume (5). Let $I = (a_1, \ldots, a_n)$ be a proper finitely generated
ideal of $R$. As $I \not = R$ we see that
$R \to R^{\oplus n}$, $x \mapsto (xa_1, \ldots, xa_n)$
is not a split injection. Hence it has a nonzero kernel and we conclude
that $\text{Ann}_R(I) \not = 0$. Thus (1) holds.
\end{proof}


\begin{lemma}
\label{more-algebra-lemma-P-universally-injective}
Let $R$ be a ring having property (P) of
Lemma \ref{more-algebra-lemma-auto-ass-implies-P}.
Let $u : N \to M$ be a homomorphism of projective $R$-modules.
Then $u$ is universally injective if and only if $u$ is injective.
\end{lemma}

\begin{proof}
Assume $u$ is injective. Our goal is to show $u$ is universally injective.
First we choose a module $Q$ such that $N \oplus Q$ is free. On considering
the map $N \oplus Q \to M \oplus Q$ we see that it suffices to prove
the lemma in case $N$ is free. In this case $N$ is a directed colimit of
finite free $R$-modules. Thus we reduce to the case that $N$ is a finite
free $R$-module, say $N = R^{\oplus n}$. We prove the lemma by induction
on $n$. The case $n = 0$ is trivial.

\medskip\noindent
Let $u : R^{\oplus n} \to M$ be an injective module map with $M$ projective.
Choose an $R$-module $Q$ such that $M \oplus Q$ is free. After replacing
$u$ by the composition $R^{\oplus n} \to M \to M \oplus Q$ we see that
we may assume that $M$ is free. Then we can find a direct summand
$R^{\oplus m} \subset M$ such that $u(R^{\oplus n}) \subset R^{\oplus m}$.
Hence we may assume that $M = R^{\oplus m}$.
In this case $u$ is given by a matrix $A = (a_{ij})$ so that
$u(x_1, \ldots, x_n) = (\sum x_i a_{i1}, \ldots, \sum x_i a_{im})$.
As $u$ is injective, in particular
$u(x, 0, \ldots, 0) = (xa_{11}, xa_{12}, \ldots, xa_{1m}) \not = 0$ if
$x \not = 0$, and as $R$ has property (P) we see that
$a_{11}R + a_{12}R + \ldots + a_{1m}R = R$. Hence see that
$R(a_{11}, \ldots, a_{1m}) \subset R^{\oplus m}$ is a direct summand
of $R^{\oplus m}$, in particular $R^{\oplus m}/R(a_{11}, \ldots, a_{1m})$
is a projective $R$-module. We get a commutative diagram
$$
\xymatrix{
0 \ar[r] &
R \ar[rr] \ar[d]^1 & & R^{\oplus n} \ar[r] \ar[d]^u &
R^{\oplus n - 1} \ar[r] \ar[d] & 0 \\
0 \ar[r] & R \ar[rr]^{(a_{11}, \ldots, a_{1m})} & &
R^{\oplus m} \ar[r] & R^{\oplus m}/R(a_{11}, \ldots, a_{1m}) \ar[r] & 0
}
$$
with split exact rows. Thus the right vertical arrow is injective
and we may apply the induction hypothesis to conclude that
the right vertical arrow is universally injective. It follows that the
middle vertical arrow is universally injective.
\end{proof}


\begin{lemma}
\label{more-algebra-lemma-auto-ass-implies-P}
An auto-associated ring $R$ has the following property: (P)
Every proper finitely generated ideal $I \subset R$ has a nonzero
annihilator.
\end{lemma}

\begin{proof}
By assumption there exists a nonzero element $x \in R$ such that for every
$f \in \mathfrak m$ we have $f^n x = 0$. Say $I = (f_1, \ldots, f_r)$.
Then $x$ is in the kernel of $R \to \bigoplus R_{f_i}$. Hence we see
that there exists a nonzero $y \in R$ such that $f_i y = 0$ for all $i$, see
Algebra, Lemma \href{https://stacks.math.columbia.edu/tag/0565}{lemma when injective covering (Tag 0565)}.
As $y \in \text{Ann}_R(I)$ we win.
\end{proof}


\begin{lemma}
\label{lemma-finite-presentation-flat-along-fibre}
Let $f : X \to S$ be a morphism of schemes.
Let $\mathcal{F}$ be a quasi-coherent sheaf on $X$.
Let $s \in S$.
Assume that
\begin{enumerate}
\item $f$ is of finite presentation,
\item $\mathcal{F}$ is of finite presentation, and
\item $\mathcal{F}$ is flat over $S$ at every point of the fibre $X_s$.
\end{enumerate}
Then there exists an elementary \'etale neighbourhood
$(S', s') \to (S, s)$ and a commutative diagram of schemes
$$
\xymatrix{
X \ar[d] & X' \ar[l]^g \ar[d] \\
S & S' \ar[l]
}
$$
such that $g$ is \'etale, $X_s \subset g(X')$, the schemes
$X'$, $S'$ are affine, and such that
$\Gamma(X', g^*\mathcal{F})$ is a projective
$\Gamma(S', \mathcal{O}_{S'})$-module.
\end{lemma}

\begin{proof}
For every point $x \in X_s$ we can use
Proposition \ref{proposition-finite-presentation-flat-at-point}
to find a commutative diagram
$$
\xymatrix{
(X, x) \ar[d] & (Y_x, y_x) \ar[l]^{g_x} \ar[d] \\
(S, s) & (S_x, s_x) \ar[l]
}
$$
whose horizontal arrows are elementary \'etale neighbourhoods
such that $Y_x$, $S_x$ are affine and such that
$\Gamma(Y_x, g_x^*\mathcal{F})$ is a projective
$\Gamma(S_x, \mathcal{O}_{S_x})$-module. In particular
$g_x(Y_x) \cap X_s$ is an open neighbourhood of $x$ in $X_s$.
Because $X_s$ is quasi-compact we can find a finite number of points
$x_1, \ldots, x_n \in X_s$ such that $X_s$ is the union of
the $g_{x_i}(Y_{x_i}) \cap X_s$. Choose an elementary \'etale neighbourhood
$(S' , s') \to (S, s)$ which dominates each of the neighbourhoods
$(S_{x_i}, s_{x_i})$, see
More on Morphisms,
Lemma \href{https://stacks.math.columbia.edu/tag/057B}{lemma elementary etale neighbourhoods (Tag 057B)}.
We may also assume that $S'$ is affine.
Set $X' = \coprod Y_{x_i} \times_{S_{x_i}} S'$ and endow it with the
obvious morphism $g : X' \to X$.
By construction $g(X')$ contains $X_s$ and
$$
\Gamma(X', g^*\mathcal{F})
=
\bigoplus \Gamma(Y_{x_i}, g_{x_i}^*\mathcal{F})
\otimes_{\Gamma(S_{x_i}, \mathcal{O}_{S_{x_i}})}
\Gamma(S', \mathcal{O}_{S'}).
$$
This is a projective $\Gamma(S', \mathcal{O}_{S'})$-module, see
Algebra, Lemma \href{https://stacks.math.columbia.edu/tag/05A3}{lemma ascend properties modules (Tag 05A3)}.
\end{proof}


\begin{lemma}
\label{lemma-criterion}
Let $f : X \to S$ be a morphism of schemes of finite type.
Let $\mathcal{F}$ be a quasi-coherent $\mathcal{O}_X$-module
of finite type. Let $s \in S$. Let $(S', s') \to (S, s)$ be an
elementary \'etale neighbourhood and let
$$
\xymatrix{
X \ar[d] & X' \ar[l]^g \ar[d] \\
S & S' \ar[l]
}
$$
be a commutative diagram of morphisms of schemes. Assume
\begin{enumerate}
\item $\mathcal{F}$ is flat over $S$ at all points of $X_s$,
\item $X' \to S'$ is of finite type,
\item $g^*\mathcal{F}$ is pure along $X'_{s'}$,
\item $g : X' \to X$ is \'etale, and
\item $g(X')$ contains $\text{Ass}_{X_s}(\mathcal{F}_s)$.
\end{enumerate}
In this situation $\mathcal{F}$ is pure along $X_s$ if and only
if the image of $X' \to X \times_S S'$ contains the points of
$\text{Ass}_{X \times_S S'/S'}(\mathcal{F} \times_S S')$
lying over points in $S'$ which specialize to $s'$.
\end{lemma}

\begin{proof}
Since the morphism $S' \to S$ is \'etale, we see that if $\mathcal{F}$
is pure along $X_s$, then $\mathcal{F} \times_S S'$ is pure along
$X_s$, see
Lemma \ref{lemma-quasi-finite-base-change}.
Since purity satisfies flat descent, see
Lemma \ref{lemma-flat-descend-pure},
we see that if $\mathcal{F} \times_S S'$ is pure along $X_{s'}$, then
$\mathcal{F}$ is pure along $X_s$. Hence we may replace $S$ by $S'$
and assume that $S = S'$ so that $g : X' \to X$ is an \'etale morphism
between schemes of finite type over $S$. Moreover, we may replace
$S$ by $\Spec(\mathcal{O}_{S, s})$ and assume that $S$ is local.

\medskip\noindent
First, assume that $\mathcal{F}$ is pure along $X_s$.
In this case every point of $\text{Ass}_{X/S}(\mathcal{F})$
specializes to a point of $X_s$ by purity. Hence by
Lemma \ref{lemma-associated-point-specializes}
we see that every point of $\text{Ass}_{X/S}(\mathcal{F})$
specializes to a point of $\text{Ass}_{X_s}(\mathcal{F}_s)$.
Thus every point of $\text{Ass}_{X/S}(\mathcal{F})$ is in the
image of $g$ (as the image is open and contains
$\text{Ass}_{X_s}(\mathcal{F}_s)$).

\medskip\noindent
Conversely, assume that $g(X')$ contains $\text{Ass}_{X/S}(\mathcal{F})$.
Let $S^h = \Spec(\mathcal{O}_{S, s}^h)$ be the henselization
of $S$ at $s$. Denote $g^h : (X')^h \to X^h$ the base change of $g$
by $S^h \to S$, and denote $\mathcal{F}^h$ the pullback of $\mathcal{F}$
to $X^h$. By
Divisors, Lemma \href{https://stacks.math.columbia.edu/tag/05DC}{lemma base change relative assassin (Tag 05DC)} and
Remark \href{https://stacks.math.columbia.edu/tag/05KL}{remark base change relative assassin (Tag 05KL)}
the relative assassin $\text{Ass}_{X^h/S^h}(\mathcal{F}^h)$
is the inverse image of $\text{Ass}_{X/S}(\mathcal{F})$ via the projection
$X^h \to X$. As we have assumed that $g(X')$ contains
$\text{Ass}_{X/S}(\mathcal{F})$ we conclude that the base change
$g^h((X')^h) = g(X') \times_S S^h$ contains
$\text{Ass}_{X^h/S^h}(\mathcal{F}^h)$. In this way
we reduce to the case where $S$ is the spectrum of a henselian local ring.
Let $x \in \text{Ass}_{X/S}(\mathcal{F})$. To finish the proof of the
lemma we have to show that $x$ specializes to a point of $X_s$, see
criterion (4) for purity in discussion following
Definition \ref{definition-pure}.
By assumption there exists a $x' \in X'$ such that $g(x') = x$.
As $g : X' \to X$ is \'etale, we see that
$x' \in \text{Ass}_{X'/S}(g^*\mathcal{F})$, see
Lemma \ref{lemma-etale-weak-assassin-up-down} (applied to
the morphism of fibres $X'_w \to X_w$ where $w \in S$ is the image of $x'$).
Since $g^*\mathcal{F}$ is pure along $X'_s$ we see that $x' \leadsto y$
for some $y \in X'_s$. Hence $x = g(x') \leadsto g(y)$ and
$g(y) \in X_s$ as desired.
\end{proof}


\begin{lemma}
\label{lemma-finite-type-flat-pure-along-fibre-is-universal}
Let $f : X \to S$ be a morphism of schemes.
Let $\mathcal{F}$ be a quasi-coherent $\mathcal{O}_X$-module.
Let $s \in S$.
Assume
\begin{enumerate}
\item $f$ is of finite type,
\item $\mathcal{F}$ is of finite type,
\item $\mathcal{F}$ is flat over $S$ at all points of $X_s$, and
\item $\mathcal{F}$ is pure along $X_s$.
\end{enumerate}
Then $\mathcal{F}$ is universally pure along $X_s$.
\end{lemma}

\begin{proof}
We first make a preliminary remark. Suppose that $(S', s') \to (S, s)$
is an elementary \'etale neighbourhood. Denote $\mathcal{F}'$ the
pullback of $\mathcal{F}$ to $X' = X \times_S S'$.
By the discussion following
Definition \ref{definition-pure}
we see that $\mathcal{F}'$ is pure along $X'_{s'}$. Moreover, $\mathcal{F}'$
is flat over $S'$ along $X'_{s'}$. Then it suffices to prove
that $\mathcal{F}'$ is universally pure along $X'_{s'}$. Namely, given
any morphism $(T, t) \to (S, s)$ of pointed schemes
the fibre product $(T', t') = (T \times_S S', (t, s'))$ is flat over $(T, t)$
and hence if $\mathcal{F}_{T'}$ is pure along $X_{t'}$ then
$\mathcal{F}_T$ is pure along $X_t$ by
Lemma \ref{lemma-flat-descend-pure}.
Thus during the proof we may always replace $(s, S)$ by an elementary
\'etale neighbourhood.
We may also replace $S$ by $\Spec(\mathcal{O}_{S, s})$
due to the local nature of the problem.

\medskip\noindent
Choose an elementary \'etale neighbourhood $(S', s') \to (S, s)$ and
a commutative diagram
$$
\xymatrix{
X \ar[d] & X' \ar[l]^g \ar[d] \\
S & \Spec(\mathcal{O}_{S', s'}) \ar[l]
}
$$
such that $X' \to X \times_S \Spec(\mathcal{O}_{S', s'})$
is \'etale, $X_s = g((X')_{s'})$, the scheme $X'$ is affine,
and such that $\Gamma(X', g^*\mathcal{F})$ is a free
$\mathcal{O}_{S', s'}$-module, see
Lemma \ref{lemma-finite-type-flat-along-fibre-free-variant}.
Note that $X' \to \Spec(\mathcal{O}_{S', s'})$ is of finite type
(as a quasi-compact morphism which is the composition of an \'etale morphism
and the base change of a finite type morphism).
By our preliminary remarks in the first paragraph of the proof
we may replace $S$ by $\Spec(\mathcal{O}_{S', s'})$. Hence
we may assume there exists a commutative diagram
$$
\xymatrix{
X \ar[dr] & & X' \ar[ll]^g \ar[ld] \\
& S &
}
$$
of schemes of finite type over $S$, where $g$ is \'etale, $X_s \subset g(X')$,
with $S$ local with closed point $s$, with $X'$ affine, and with
$\Gamma(X', g^*\mathcal{F})$ a free $\Gamma(S, \mathcal{O}_S)$-module.
Note that in this case $g^*\mathcal{F}$ is universally pure over $S$, see
Lemma \ref{lemma-affine-locally-projective-pure}.

\medskip\noindent
In this situation we apply
Lemma \ref{lemma-criterion}
to deduce that $\text{Ass}_{X/S}(\mathcal{F}) \subset g(X')$
from our assumption that $\mathcal{F}$ is pure along $X_s$
and flat over $S$ along $X_s$. By
Divisors, Lemma \href{https://stacks.math.columbia.edu/tag/05DC}{lemma base change relative assassin (Tag 05DC)} and
Remark \href{https://stacks.math.columbia.edu/tag/05KL}{remark base change relative assassin (Tag 05KL)}
we see that for any morphism of pointed schemes
$(T, t) \to (S, s)$ we have
$$
\text{Ass}_{X_T/T}(\mathcal{F}_T) \subset
(X_T \to X)^{-1}(\text{Ass}_{X/S}(\mathcal{F})) \subset
g(X') \times_S T = g_T(X'_T).
$$
Hence by
Lemma \ref{lemma-criterion}
applied to the base change of our displayed diagram to $(T, t)$
we conclude that $\mathcal{F}_T$ is pure along $X_t$ as desired.
\end{proof}


\begin{lemma}
\label{lemma-affine-locally-projective-pure}
Let $f : X \to S$ be a finite type, affine morphism of schemes.
Let $\mathcal{F}$ be a finite type quasi-coherent $\mathcal{O}_X$-module
such that $f_*\mathcal{F}$ is locally projective on $S$, see
Properties, Definition \ref{properties-definition-locally-projective}.
Then $\mathcal{F}$ is universally pure over $S$.
\end{lemma}

\begin{proof}
After reducing to the case where $S$ is the spectrum of a henselian
local ring this follows from
Lemma \ref{lemma-explain-why-pure}.
\end{proof}


\begin{lemma}
\label{lemma-explain-why-pure}
Let $R$ be a local ring with maximal ideal $\mathfrak m$.
Let $R \to S$ be a ring map. Let $N$ be an $S$-module.
Assume
\begin{enumerate}
\item $N$ is projective as an $R$-module, and
\item $S/\mathfrak mS$ is Noetherian and $N/\mathfrak mN$ is a finite
$S/\mathfrak mS$-module.
\end{enumerate}
Then for any prime $\mathfrak q \subset S$ which is an associated prime of
$N \otimes_R \kappa(\mathfrak p)$ where $\mathfrak p = R \cap \mathfrak q$
we have $\mathfrak q + \mathfrak m S \not = S$.
\end{lemma}

\begin{proof}
Note that the hypotheses of
Lemmas \ref{lemma-homothety-spectrum} and
\ref{lemma-invert-universally-injective}
are satisfied. We will use the conclusions of these lemmas without further
mention. Let $\Sigma \subset S$ be the multiplicative set of elements
which are not zerodivisors on $N/\mathfrak mN$. The map
$N \to \Sigma^{-1}N$ is $R$-universally injective. Hence we see that
any $\mathfrak q \subset S$ which is an associated prime of
$N \otimes_R \kappa(\mathfrak p)$ is also an associated prime of
$\Sigma^{-1}N \otimes_R \kappa(\mathfrak p)$. Clearly this implies that
$\mathfrak q$ corresponds to a prime of $\Sigma^{-1}S$.
Thus $\mathfrak q \subset \mathfrak q'$ where $\mathfrak q'$
corresponds to an associated prime of $N/\mathfrak mN$ and we win.
\end{proof}


\begin{lemma}
\label{lemma-quasi-finite-base-change}
Let $f : X \to S$ be a morphism of schemes which is of finite type.
Let $\mathcal{F}$ be a finite type quasi-coherent $\mathcal{O}_X$-module.
Let $s \in S$. Let $(S', s') \to (S, s)$ be a morphism of pointed schemes.
If $S' \to S$ is quasi-finite at $s'$ and $\mathcal{F}$ is pure along $X_s$,
then $\mathcal{F}_{S'}$ is pure along $X_{s'}$.
\end{lemma}

\begin{proof}
It $(T \to S', t' \leadsto t, \xi)$ is an impurity of
$\mathcal{F}_{S'}$ above $s'$ with $T \to S'$ quasi-finite at $t$,
then $(T \to S, t' \to t, \xi)$ is an impurity of $\mathcal{F}$
above $s$ with $T \to S$ quasi-finite at $t$, see
Morphisms, Lemma \href{https://stacks.math.columbia.edu/tag/01TL}{lemma composition quasi finite (Tag 01TL)}.
Hence the lemma follows immediately from the characterization (2)
of purity given following
Definition \ref{definition-pure}.
\end{proof}


\begin{lemma}
\label{lemma-flat-descend-pure}
Let $f : X \to S$ be a morphism of schemes which is of finite type.
Let $\mathcal{F}$ be a finite type quasi-coherent $\mathcal{O}_X$-module.
Let $s \in S$. Let $(S', s') \to (S, s)$ be a morphism of pointed schemes.
Assume $S' \to S$ is flat at $s'$.
\begin{enumerate}
\item If $\mathcal{F}_{S'}$ is pure along $X_{s'}$,
then $\mathcal{F}$ is pure along $X_s$.
\item If $\mathcal{F}_{S'}$ is universally pure along $X_{s'}$,
then $\mathcal{F}$ is universally pure along $X_s$.
\end{enumerate}
\end{lemma}

\begin{proof}
Let $(T \to S, t' \leadsto t, \xi)$ be an impurity of
$\mathcal{F}$ above $s$. Set $T_1 = T \times_S S'$, and let $t_1$
be the unique point of $T_1$ mapping to $t$ and $s'$. Since
$T_1 \to T$ is flat at $t_1$, see
Morphisms, Lemma \href{https://stacks.math.columbia.edu/tag/01U9}{lemma base change flat (Tag 01U9)},
there exists a specialization $t'_1 \leadsto t_1$ lying over
$t' \leadsto t$, see
Algebra, Section \href{https://stacks.math.columbia.edu/tag/00HU}{section going up (Tag 00HU)}.
Choose a point $\xi_1 \in X_{t'_1}$ which corresponds to a generic
point of $\Spec(\kappa(t'_1) \otimes_{\kappa(t')} \kappa(\xi))$, see
Schemes, Lemma \href{https://stacks.math.columbia.edu/tag/01JT}{lemma points fibre product (Tag 01JT)}.
By
Divisors, Lemma \href{https://stacks.math.columbia.edu/tag/05DC}{lemma base change relative assassin (Tag 05DC)}
we see that $\xi_1 \in \text{Ass}_{X_{T_1}/T_1}(\mathcal{F}_{T_1})$.
As the Zariski closure of $\{\xi_1\}$ in $X_{T_1}$ maps into the
Zariski closure of $\{\xi\}$ in $X_T$ we conclude that
this closure is disjoint from $X_{t_1}$. Hence
$(T_1 \to S', t'_1 \leadsto t_1, \xi_1)$
is an impurity of $\mathcal{F}_{S'}$ above $s'$.
In other words we have proved the contrapositive to part (2) of the
lemma. Finally, if $(T, t) \to (S, s)$ is an elementary
\'etale neighbourhood, then $(T_1, t_1) \to (S', s')$ is an
elementary \'etale neighbourhood too, and in this way we see that (1) holds.
\end{proof}


\begin{lemma}
\label{lemma-associated-point-specializes}
Let $f : X \to S$ be a morphism of schemes of finite type.
Let $\mathcal{F}$ be a quasi-coherent $\mathcal{O}_X$-module of finite type.
Let $s \in S$.
Assume that $\mathcal{F}$ is flat over $S$ at all points of $X_s$.
Let $x' \in \text{Ass}_{X/S}(\mathcal{F})$ with $f(x') = s'$
such that $s' \leadsto s$ is a specialization in $S$. If
$x'$ specializes to a point of $X_s$, then $x' \leadsto x$
with $x \in \text{Ass}_{X_s}(\mathcal{F}_s)$.
\end{lemma}

\begin{proof}
Say $x' \leadsto t$ with $t \in X_s$. Then we can find specializations
$x' \leadsto x \leadsto t$ with $x$ corresponding to a generic
point of an irreducible component of
$\overline{\{x'\}} \cap f^{-1}(\{s\})$. By assumption $\mathcal{F}$ is flat
over $S$ at $x$. By More on Morphisms, Lemma
\href{https://stacks.math.columbia.edu/tag/0GSI}{lemma associated point specializes (Tag 0GSI)}
we see that $x \in \text{Ass}_{X/S}(\mathcal{F})$ as desired.
\end{proof}


\begin{lemma}
\label{lemma-finite-type-flat-along-fibre-free-variant}
Let $f : X \to S$ be a morphism of schemes.
Let $\mathcal{F}$ be a quasi-coherent sheaf on $X$.
Let $s \in S$.
Assume that
\begin{enumerate}
\item $f$ is of finite type,
\item $\mathcal{F}$ is of finite type, and
\item $\mathcal{F}$ is flat over $S$ at all points of $X_s$.
\end{enumerate}
Then there exists an elementary \'etale neighbourhood $(S', s') \to (S, s)$
and a commutative diagram of schemes
$$
\xymatrix{
X \ar[d] & X' \ar[l]^g \ar[d] \\
S & \Spec(\mathcal{O}_{S', s'}) \ar[l]
}
$$
such that $X' \to X \times_S \Spec(\mathcal{O}_{S', s'})$
is \'etale, $X_s = g((X')_{s'})$, the scheme $X'$ is affine,
and such that $\Gamma(X', g^*\mathcal{F})$ is a free
$\mathcal{O}_{S', s'}$-module.
\end{lemma}

\begin{proof}
(The only difference with
Lemma \href{https://stacks.math.columbia.edu/tag/05L0}{lemma finite type flat along fibre free (Tag 05L0)}
is that we do not assume $f$ is of finite presentation.)
For every point $x \in X_s$ we can use
Lemma \ref{lemma-finite-type-flat-at-point-free-variant}
to find an elementary \'etale neighbourhood $(S_x , s_x) \to (S, s)$
and a commutative diagram
$$
\xymatrix{
(X, x) \ar[d] & (Y_x, y_x) \ar[l]^{g_x} \ar[d] \\
(S, s) & (\Spec(\mathcal{O}_{S_x, s_x}), s_x) \ar[l]
}
$$
such that $Y_x \to X \times_S \Spec(\mathcal{O}_{S_x, s_x})$
is \'etale, $\kappa(x) = \kappa(y_x)$, the scheme $Y_x$ is affine, and
such that $\Gamma(Y_x, g_x^*\mathcal{F})$ is a free
$\mathcal{O}_{S_x, s_x}$-module. In particular
$g_x((Y_x)_{s_x})$ is an open neighbourhood of $x$ in $X_s$.
Because $X_s$ is quasi-compact we can find a finite number of points
$x_1, \ldots, x_n \in X_s$ such that $X_s$ is the union of
the $g_{x_i}((Y_{x_i})_{s_{x_i}})$. Choose an elementary \'etale neighbourhood
$(S' , s') \to (S, s)$ which dominates each of the neighbourhoods
$(S_{x_i}, s_{x_i})$, see
More on Morphisms,
Lemma \href{https://stacks.math.columbia.edu/tag/057B}{lemma elementary etale neighbourhoods (Tag 057B)}.
Set
$$
X' = \coprod Y_{x_i} \times_{\Spec(\mathcal{O}_{S_{x_i}, s_{x_i}})}
\Spec(\mathcal{O}_{S', s'})
$$
and endow it with the obvious morphism $g : X' \to X$.
By construction $X_s = g(X'_{s'})$ and
$$
\Gamma(X', g^*\mathcal{F})
=
\bigoplus \Gamma(Y_{x_i}, g_{x_i}^*\mathcal{F})
\otimes_{\mathcal{O}_{S_{x_i}, s_{x_i}}}
\mathcal{O}_{S', s'}.
$$
This is a free $\mathcal{O}_{S', s'}$-module as a direct sum
of base changes of free modules.
\end{proof}


\begin{lemma}
\label{lemma-finite-type-flat-at-point-free-variant}
Let $f : X \to S$ be a morphism of schemes.
Let $\mathcal{F}$ be a quasi-coherent sheaf on $X$.
Let $x \in X$ with image $s \in S$.
Assume that
\begin{enumerate}
\item $f$ is locally of finite type,
\item $\mathcal{F}$ is of finite type, and
\item $\mathcal{F}$ is flat at $x$ over $S$.
\end{enumerate}
Then there exists an elementary \'etale neighbourhood $(S', s') \to (S, s)$
and a commutative diagram of pointed schemes
$$
\xymatrix{
(X, x) \ar[d] & (X', x') \ar[l]^g \ar[d] \\
(S, s) & (\Spec(\mathcal{O}_{S', s'}), s') \ar[l]
}
$$
such that $X' \to X \times_S \Spec(\mathcal{O}_{S', s'})$
is \'etale, $\kappa(x) = \kappa(x')$, the scheme $X'$ is
affine, and such that $\Gamma(X', g^*\mathcal{F})$ is a free
$\mathcal{O}_{S', s'}$-module.
\end{lemma}

\begin{proof}
(The only difference with
Lemma \ref{lemma-finite-type-flat-at-point-free}
is that we do not assume $f$ is of finite presentation.)
The problem is local on $X$ and $S$. Hence we may assume $X$ and
$S$ are affine, say $X = \Spec(B)$ and $S = \Spec(A)$.
Since $B$ is a finite type $A$-algebra we can find a surjection
$A[x_1, \ldots, x_n] \to B$. In other words, we can choose a closed
immersion $i : X \to \mathbf{A}^n_S$. Set $t = i(x)$ and
$\mathcal{G} = i_*\mathcal{F}$. Note that $\mathcal{G}_t \cong \mathcal{F}_x$
are $\mathcal{O}_{S, s}$-modules. Hence $\mathcal{G}$ is flat over $S$ at $t$.
We apply
Lemma \ref{lemma-finite-type-flat-at-point-free}
to the morphism $\mathbf{A}^n_S \to S$, the point $t$, and the
sheaf $\mathcal{G}$. Thus we can find an
elementary \'etale neighbourhood $(S', s') \to (S, s)$
and a commutative diagram of pointed schemes
$$
\xymatrix{
(\mathbf{A}^n_S, t) \ar[d] & (Y, y) \ar[l]^h \ar[d] \\
(S, s) & (\Spec(\mathcal{O}_{S', s'}), s') \ar[l]
}
$$
such that $Y \to \mathbf{A}^n_{\mathcal{O}_{S', s'}}$
is \'etale, $\kappa(t) = \kappa(y)$, the scheme $Y$ is
affine, and such that $\Gamma(Y, h^*\mathcal{G})$ is a projective
$\mathcal{O}_{S', s'}$-module. Then a solution to the original
problem is given by the closed subscheme
$X' = Y \times_{\mathbf{A}^n_S} X$ of $Y$.
\end{proof}


\begin{lemma}
\label{lemma-finite-type-flat-at-point-free}
Let $f : X \to S$ be a morphism of schemes.
Let $\mathcal{F}$ be a quasi-coherent sheaf on $X$.
Let $x \in X$ with image $s \in S$.
Assume that
\begin{enumerate}
\item $f$ is locally of finite presentation,
\item $\mathcal{F}$ is of finite type, and
\item $\mathcal{F}$ is flat at $x$ over $S$.
\end{enumerate}
Then there exists an elementary \'etale neighbourhood $(S', s') \to (S, s)$
and a commutative diagram of pointed schemes
$$
\xymatrix{
(X, x) \ar[d] & (X', x') \ar[l]^g \ar[d] \\
(S, s) & (\Spec(\mathcal{O}_{S', s'}), s') \ar[l]
}
$$
such that $X' \to X \times_S \Spec(\mathcal{O}_{S', s'})$
is \'etale, $\kappa(x) = \kappa(x')$, the scheme $X'$ is
affine of finite presentation over $\mathcal{O}_{S', s'}$,
the sheaf $g^*\mathcal{F}$ is of finite presentation over $\mathcal{O}_{X'}$,
and such that $\Gamma(X', g^*\mathcal{F})$ is a free
$\mathcal{O}_{S', s'}$-module.
\end{lemma}

\begin{proof}
To prove the lemma we may replace $(S, s)$ by any elementary \'etale
neighbourhood, and we may also replace $S$ by
$\Spec(\mathcal{O}_{S, s})$. Hence by
Proposition \ref{proposition-finite-type-flat-at-point}
we may assume that $\mathcal{F}$ is finitely presented and flat over
$S$ in a neighbourhood of $x$. In this case the result follows from
Proposition \ref{proposition-finite-presentation-flat-at-point}
because
Algebra, Theorem \href{https://stacks.math.columbia.edu/tag/0593}{theorem projective free over local ring (Tag 0593)}
assures us that projective $=$ free over a local ring.
\end{proof}


\begin{lemma}
\label{lemma-quasi-finite-impurity-elementary}
In Situation \ref{situation-pre-pure}.
If there exists an impurity $(g : T \to S, t' \leadsto t, \xi)$
of $\mathcal{F}$ above $s$ with $g$ quasi-finite at $t$, then there
exists an impurity $(g : T \to S, t' \leadsto t, \xi)$ such that
$(T, t) \to (S, s)$ is an elementary \'etale neighbourhood.
\end{lemma}

\begin{proof}
Let $(g : T \to S, t' \leadsto t, \xi)$ be an impurity of
$\mathcal{F}$ above $s$ such that $g$ is quasi-finite at $t$.
After shrinking $T$ we may assume that $g$ is locally of finite type.
Apply
More on Morphisms,
Lemma \href{https://stacks.math.columbia.edu/tag/02LK}{lemma etale makes quasi finite finite at point (Tag 02LK)}
to $T \to S$ and $t \mapsto s$. This gives us a diagram
$$
\xymatrix{
T \ar[d] & T \times_S U \ar[l] \ar[d] & V \ar[l] \ar[ld] \\
S & U \ar[l]
}
$$
where $(U, u) \to (S, s)$ is an elementary \'etale neighbourhood
and $V \subset T \times_S U$ is an open neighbourhood of $v = (t, u)$
such that $V \to U$ is finite and such that $v$ is the unique point of $V$
lying over $u$. Since the morphism $V \to T$ is \'etale
hence flat we see that there exists a specialization $v' \leadsto v$ such
that $v' \mapsto t'$. Note that $\kappa(t') \subset \kappa(v')$
is finite separable. Pick any point $\zeta \in X_{v'}$ mapping to
$\xi \in X_{t'}$. By
Divisors, Lemma \href{https://stacks.math.columbia.edu/tag/05DC}{lemma base change relative assassin (Tag 05DC)}
we see that $\zeta \in \text{Ass}_{X_V/V}(\mathcal{F}_V)$.
Moreover, the closure $\overline{\{\zeta\}}$ does not meet
the fibre $X_v$ as by assumption the closure $\overline{\{\xi\}}$
does not meet $X_t$. In other words $(V \to S, v' \leadsto v, \zeta)$
is an impurity of $\mathcal{F}$ above $S$.

\medskip\noindent
Next, let $u' \in U'$ be the image of $v'$ and let
$\theta \in X_U$ be the image of $\zeta$.
Then $\theta \mapsto u'$ and $u' \leadsto u$.
By
Divisors, Lemma \href{https://stacks.math.columbia.edu/tag/05DC}{lemma base change relative assassin (Tag 05DC)}
we see that $\theta \in \text{Ass}_{X_U/U}(\mathcal{F})$.
Moreover, as $\pi : X_V \to X_U$ is finite we see that
$\pi\big(\overline{\{\zeta\}}\big) = \overline{\{\pi(\zeta)\}}$. Since
$v$ is the unique point of $V$ lying over $u$ we see that
$X_u \cap \overline{\{\pi(\zeta)\}} = \emptyset$ because
$X_v \cap \overline{\{\zeta\}} = \emptyset$. In this way we conclude that
$(U \to S, u' \leadsto u, \theta)$ is an impurity of
$\mathcal{F}$ above $s$ and we win.
\end{proof}


\begin{lemma}
\label{lemma-impurity-on-henselization}
In Situation \ref{situation-pre-pure}.
If there exists an impurity $(S^h \to S, s' \leadsto s, \xi)$
of $\mathcal{F}$ above $s$ then there exists an impurity
$(T \to S, t' \leadsto t, \xi)$ of $\mathcal{F}$ above $s$
where $(T, t) \to (S, s)$ is an elementary \'etale neighbourhood.
\end{lemma}

\begin{proof}
We may replace $S$ by an affine neighbourhood of $s$.
Say $S = \Spec(A)$ and $s$ corresponds to the prime
$\mathfrak p \subset A$. Then
$\mathcal{O}_{S, s}^h = \colim_{(T, t)} \Gamma(T, \mathcal{O}_T)$
where the limit is over the opposite of the
cofiltered category of affine elementary \'etale neighbourhoods
$(T, t)$ of $(S, s)$, see
More on Morphisms,
Lemma \href{https://stacks.math.columbia.edu/tag/05KS}{lemma describe henselization (Tag 05KS)}
and its proof. Hence $S^h = \lim_i T_i$ and we win by
Lemma \ref{lemma-impure-limit}.
\end{proof}


\begin{lemma}
\label{lemma-impure-limit}
In Situation \ref{situation-pre-pure}.
Let $(g : T \to S, t' \leadsto t, \xi)$ be an impurity of
$\mathcal{F}$ above $s$. Assume $T = \lim_{i \in I} T_i$
is a directed limit of affine schemes over $S$. Then for
some $i$ the triple $(T_i \to S, t'_i \leadsto t_i, \xi_i)$
is an impurity of $\mathcal{F}$ above $s$.
\end{lemma}

\begin{proof}
The notation in the statement means this: Let $p_i : T \to T_i$
be the projection morphisms, let $t_i = p_i(t)$ and $t'_i = p_i(t')$.
Finally $\xi_i \in X_{T_i}$ is the image of $\xi$. By
Divisors, Lemma \href{https://stacks.math.columbia.edu/tag/05DC}{lemma base change relative assassin (Tag 05DC)}
it is true that $\xi_i$ is a point of the relative
assassin of $\mathcal{F}_{T_i}$ over $T_i$. Thus the only point is to
show that $\overline{\{\xi_i\}} \cap X_{t_i} = \emptyset$ for some $i$.

\medskip\noindent
First proof. Let $Z_i = \overline{\{\xi_i\}} \subset X_{T_i}$
and $Z = \overline{\{\xi\}} \subset X_T$
endowed with the reduced induced scheme structure.
Then $Z = \lim Z_i$ by
Limits, Lemma \href{https://stacks.math.columbia.edu/tag/0CUG}{lemma inverse limit irreducibles (Tag 0CUG)}.
Choose a field $k$ and a morphism $\Spec(k) \to T$ whose image is $t$.
Then
$$
\emptyset =
Z \times_T \Spec(k) = (\lim Z_i) \times_{(\lim T_i)} \Spec(k)
= \lim Z_i \times_{T_i} \Spec(k)
$$
because limits commute with fibred products (limits commute with limits).
Each $Z_i \times_{T_i} \Spec(k)$ is quasi-compact because $X_{T_i} \to T_i$
is of finite type and hence $Z_i \to T_i$ is of finite type.
Hence $Z_i \times_{T_i} \Spec(k)$ is empty for some $i$ by
Limits, Lemma \href{https://stacks.math.columbia.edu/tag/01Z2}{lemma limit nonempty (Tag 01Z2)}.
Since the image of the composition $\Spec(k) \to T \to T_i$ is $t_i$
we obtain what we want.

\medskip\noindent
Second proof. Set $Z = \overline{\{\xi\}}$. Apply
Limits, Lemma \href{https://stacks.math.columbia.edu/tag/05BD}{lemma separate (Tag 05BD)}
to this situation to obtain an open neighbourhood
$V \subset T$ of $t$, a commutative diagram
$$
\xymatrix{
V \ar[d] \ar[r]_a & T' \ar[d]^b \\
T \ar[r]^g & S,
}
$$
and a closed subscheme $Z' \subset X_{T'}$ such that
\begin{enumerate}
\item the morphism $b : T' \to S$ is locally of finite presentation,
\item we have $Z' \cap X_{a(t)} = \emptyset$, and
\item $Z \cap X_V$ maps into $Z'$ via the morphism $X_V \to X_{T'}$.
\end{enumerate}
We may assume $V$ is an affine open of $T$, hence by
Limits, Lemmas \href{https://stacks.math.columbia.edu/tag/01Z4}{lemma descend opens (Tag 01Z4)} and
\href{https://stacks.math.columbia.edu/tag/01Z6}{lemma limit affine (Tag 01Z6)}
we can find an $i$ and an affine open $V_i \subset T_i$ with
$V = f_i^{-1}(V_i)$. By
Limits,
Proposition \href{https://stacks.math.columbia.edu/tag/01ZC}{proposition characterize locally finite presentation (Tag 01ZC)}
after possibly increasing $i$ a bit we can find a morphism
$a_i : V_i \to T'$ such that $a = a_i \circ f_i|_V$.
The induced morphism $X_{V_i} \to X_{T'}$ maps $\xi_i$ into
$Z'$. As $Z' \cap X_{a(t)} = \emptyset$ we conclude that
$(T_i \to S, t'_i \leadsto t_i, \xi_i)$ is an impurity of
$\mathcal{F}$ above $s$.
\end{proof}


\begin{lemma}
\label{lemma-impure-finite-presentation}
In Situation \ref{situation-pre-pure}.
If there exists an impurity of $\mathcal{F}$ above $s$, then
there exists an impurity $(g : T \to S, t' \leadsto t, \xi)$
of $\mathcal{F}$ above $s$ such that $g$ is locally of finite
presentation and $t$ a closed point of the fibre of $g$ above $s$.
\end{lemma}

\begin{proof}
Let $(g : T \to S, t' \leadsto t, \xi)$ be any impurity of
$\mathcal{F}$ above $s$. We apply
Limits, Lemma \href{https://stacks.math.columbia.edu/tag/05BD}{lemma separate (Tag 05BD)}
to $t \in T$ and $Z = \overline{\{\xi\}}$ to obtain an open neighbourhood
$V \subset T$ of $t$, a commutative diagram
$$
\xymatrix{
V \ar[d] \ar[r]_a & T' \ar[d]^b \\
T \ar[r]^g & S,
}
$$
and a closed subscheme $Z' \subset X_{T'}$ such that
\begin{enumerate}
\item the morphism $b : T' \to S$ is locally of finite presentation,
\item we have $Z' \cap X_{a(t)} = \emptyset$, and
\item $Z \cap X_V$ maps into $Z'$ via the morphism $X_V \to X_{T'}$.
\end{enumerate}
As $t'$ specializes to $t$ we may replace $T$ by the open neighbourhood
$V$ of $t$. Thus we have a commutative diagram
$$
\xymatrix{
X_T \ar[d] \ar[r] &
X_{T'} \ar[d] \ar[r] &
X \ar[d] \\
T \ar[r]^a & T' \ar[r]^b & S
}
$$
where $b \circ a = g$. Let $\xi' \in X_{T'}$ denote the
image of $\xi$. By
Divisors, Lemma \href{https://stacks.math.columbia.edu/tag/05DC}{lemma base change relative assassin (Tag 05DC)}
we see that $\xi' \in \text{Ass}_{X_{T'}/T'}(\mathcal{F}_{T'})$.
Moreover, by construction the closure of $\overline{\{\xi'\}}$
is contained in the closed subset $Z'$ which avoids the fibre
$X_{a(t)}$. In this way we see that $(T' \to S, a(t') \leadsto a(t), \xi')$
is an impurity of $\mathcal{F}$ above $s$.

\medskip\noindent
Thus we may assume that $g : T \to S$ is locally of finite presentation.
Let $Z = \overline{\{\xi\}}$. By assumption $Z_t = \emptyset$. By
More on Morphisms, Lemma \href{https://stacks.math.columbia.edu/tag/054W}{lemma empty generic fibre (Tag 054W)}
this means that $Z_{t''} = \emptyset$ for $t''$ in an open subset
of $\overline{\{t\}}$. Since the fibre of
$T \to S$ over $s$ is a Jacobson scheme, see
Morphisms, Lemma \href{https://stacks.math.columbia.edu/tag/02J6}{lemma ubiquity Jacobson schemes (Tag 02J6)}
we find that there exist a closed point $t'' \in \overline{\{t\}}$ such that
$Z_{t''} = \emptyset$. Then $(g : T \to S, t' \leadsto t'', \xi)$ is the
desired impurity.
\end{proof}


\begin{lemma}
\label{lemma-pure-along-X-s}
In Situation \ref{situation-pre-pure} the following
are equivalent
\begin{enumerate}
\item there exists an impurity $(S^h \to S, s' \leadsto s, \xi)$
of $\mathcal{F}$ above $s$ where $S^h$ is the henselization of $S$ at $s$,
\item there exists an impurity $(T \to S, t' \leadsto t, \xi)$
of $\mathcal{F}$ above $s$ such that $(T, t) \to (S, s)$ is an
elementary \'etale neighbourhood, and
\item there exists an impurity $(T \to S, t' \leadsto t, \xi)$
of $\mathcal{F}$ above $s$ such that $T \to S$ is quasi-finite at $t$.
\end{enumerate}
\end{lemma}

\begin{proof}
As an \'etale morphism is locally quasi-finite it is clear that
(2) implies (3). We have seen that (3) implies (2) in
Lemma \ref{lemma-quasi-finite-impurity-elementary}.
We have seen that (1) implies (2) in
Lemma \ref{lemma-impurity-on-henselization}.
Finally, if $(T \to S, t' \leadsto t, \xi)$ is an impurity
of $\mathcal{F}$ above $s$ such that $(T, t) \to (S, s)$ is an
elementary \'etale neighbourhood, then we can choose a factorization
$S^h \to T \to S$ of the structure morphism $S^h \to S$.
Choose any point $s' \in S^h$ mapping to $t'$ and choose any
$\xi' \in X_{s'}$ mapping to $\xi \in X_{t'}$. Then
$(S^h \to S, s' \leadsto s, \xi')$ is an impurity of
$\mathcal{F}$ above $s$. We omit the details.
\end{proof}


\begin{lemma}
\label{algebra-lemma-pure-submodule-ML}
Let $0 \to M_1 \to M_2 \to M_3 \to 0$ be a
universally exact sequence of $R$-modules.  Then:
\begin{enumerate}
\item If $M_2$ is Mittag-Leffler, then $M_1$ is Mittag-Leffler.
\item If $M_1$ and $M_3$ are Mittag-Leffler, then $M_2$ is Mittag-Leffler.
\end{enumerate}
\end{lemma}

\begin{proof}
For any family $(Q_{\alpha})_{\alpha \in A}$ of $R$-modules we have a
commutative diagram
$$
\xymatrix{
0 \ar[r] & M_1 \otimes_R (\prod_{\alpha} Q_{\alpha}) \ar[r] \ar[d] & M_2
\otimes_R (\prod_{\alpha} Q_{\alpha}) \ar[r] \ar[d] & M_3 \otimes_R
(\prod_{\alpha} Q_{\alpha}) \ar[r] \ar[d] & 0 \\
0 \ar[r] & \prod_{\alpha}(M_1 \otimes Q_{\alpha}) \ar[r] & \prod_{\alpha}(M_2
\otimes Q_{\alpha}) \ar[r] & \prod_{\alpha}(M_3 \otimes Q_{\alpha})\ar[r] & 0
}
$$
with exact rows. Thus (1) and (2) follow from
Proposition \href{https://stacks.math.columbia.edu/tag/059M}{proposition ML tensor (Tag 059M)}.
\end{proof}


\begin{lemma}
\label{algebra-lemma-countgen-projective}
Let $M$ be an $R$-module.  If $M$ is flat, Mittag-Leffler, and countably
generated, then $M$ is projective.
\end{lemma}

\begin{proof}
By Lazard's theorem (Theorem \href{https://stacks.math.columbia.edu/tag/058G}{theorem lazard (Tag 058G)}), we can write $M =
\colim_{i
\in I} M_i$ for a directed system of finite free $R$-modules $(M_i, f_{ij})$
indexed by a set $I$.  By Lemma \href{https://stacks.math.columbia.edu/tag/059W}{lemma ML countable colimit (Tag 059W)}, we may assume
$I$ is countable.  Now let
$$
0 \to N_1 \to N_2 \to N_3 \to 0
$$
be an exact sequence of $R$-modules.  We must show that applying
$\Hom_R(M, -)$
preserves exactness.  Since $M_i$ is finite free,
$$
0 \to \Hom_R(M_i, N_1) \to \Hom_R(M_i, N_2) \to
\Hom_R(M_i, N_3) \to 0
$$
is exact for each $i$.  Since $M$ is Mittag-Leffler, $(\Hom_R(M_i,
N_{1}))$ is
a Mittag-Leffler inverse system.  So by Lemma \href{https://stacks.math.columbia.edu/tag/0598}{lemma ML exact sequence (Tag 0598)},
$$
0 \to \lim_{i \in I} \Hom_R(M_i, N_1) \to
\lim_{i \in I} \Hom_R(M_i, N_2) \to
\lim_{i \in I} \Hom_R(M_i, N_3) \to 0
$$
is exact.  But for any $R$-module $N$ there is a functorial isomorphism
$\Hom_R(M, N) \cong \lim_{i \in I} \Hom_R(M_i, N)$, so
$$
0 \to \Hom_R(M, N_1) \to \Hom_R(M, N_2) \to
\Hom_R(M, N_3) \to 0
$$
is exact.
\end{proof}


\begin{theorem}
\label{algebra-theorem-ffdescent-projectivity}
Let $R \to S$ be a faithfully flat ring map.  Let $M$ be an $R$-module.
 If the $S$-module $M \otimes_R S$ is projective, then $M$ is projective.
\end{theorem}

\begin{proof}
We are going to construct a Kaplansky d\'evissage of $M$ to show that it is a
direct sum of projective modules and hence projective.  By Theorem
\href{https://stacks.math.columbia.edu/tag/058Y}{theorem projective direct sum (Tag 058Y)} we can write $M \otimes_R S =
\bigoplus_{i \in I} Q_i$ as a direct sum of countably generated $S$-modules
$Q_i$.  Choose a well-ordering on $M$.  Using transfinite recursion we are going
to define an increasing family of submodules $M_{\alpha}$ of $M$, one for each
ordinal $\alpha$, such that $M_{\alpha} \otimes_R S$ is a direct sum of some
subset of the $Q_i$.

\medskip\noindent
For $\alpha = 0$ let $M_0 = 0$.  If $\alpha$ is a limit ordinal and $M_{\beta}$
has been defined for all $\beta < \alpha$, then define $M_\alpha =
\bigcup_{\beta < \alpha} M_{\beta}$.  Since each $M_{\beta} \otimes_R S$ for
$\beta < \alpha$ is a direct sum of a subset of the $Q_i$, the same will be
true of $M_{\alpha} \otimes_R S$.  If $\alpha + 1$ is a successor ordinal and
$M_{\alpha}$ has been defined, then define $M_{\alpha + 1}$ as follows.  If
$M_{\alpha} = M$, then let $M_{\alpha +1} = M$.  Otherwise choose the smallest
$x \in M$ (with respect to the fixed well-ordering) such that $x \notin
M_{\alpha}$. Since $S$ is flat over $R$, $(M/M_{\alpha}) \otimes_R S = M
\otimes_R S/M_{\alpha} \otimes_R S$, so since $M_{\alpha} \otimes_R S$ is
a direct sum of some $Q_i$, the same is true of $(M/M_{\alpha}) \otimes_R
S$.   By Lemma \ref{algebra-lemma-adapted-submodule}, we can find a countably generated
$R$-submodule $P$ of $M/M_{\alpha}$ containing the image of $x$ in
$M/M_{\alpha}$ and such that $P \otimes_R S$ (which equals $\Im(P
\otimes_R S \to M \otimes_R S)$ since $S$ is flat over $R$) is a
direct sum of some $Q_i$. Since $M \otimes_R S = \bigoplus_{i \in I} Q_i$
is projective and projectivity passes to direct summands, $P \otimes_R S$ is
also projective.  Thus by Lemma \ref{algebra-lemma-ffdescent-countable-projectivity},
$P$ is projective.  Finally we define $M_{\alpha + 1}$ to be the preimage of $P$
in $M$, so that $M_{\alpha + 1}/M_{\alpha} = P$ is countably generated and
projective.  In particular $M_{\alpha}$ is a direct summand of $M_{\alpha + 1}$
since projectivity of $M_{\alpha + 1}/M_{\alpha}$ implies the sequence $0
\to M_{\alpha} \to M_{\alpha + 1} \to
M_{\alpha + 1}/M_{\alpha} \to 0$ splits.

\medskip\noindent
Transfinite induction on $M$ (using the fact that we constructed
$M_{\alpha + 1}$ to contain the smallest $x \in M$ not contained in
$M_{\alpha}$) shows that
each $x \in M$ is contained in some $M_{\alpha}$.  Thus, there is some large
enough ordinal $S$ satisfying: for each $x \in M$ there is $\alpha \in S$ such
that $x \in M_{\alpha}$.  This means $(M_{\alpha})_{\alpha \in S}$ satisfies
property (1) of a Kaplansky d\'evissage of $M$.  The other properties are clear
by construction.  We conclude
$M = \bigoplus_{\alpha + 1 \in S} M_{\alpha + 1}/M_{\alpha}$.
Since each $M_{\alpha + 1}/M_{\alpha}$ is projective
by construction, $M$ is projective.
\end{proof}


\begin{lemma}
\label{lemma-flattening-module-map}
Let $A$ be a ring. Let $u : M \to N$ be a surjective map of $A$-modules.
If $M$ is projective as an $A$-module, then there exists an ideal
$I \subset A$ such that for any ring map $\varphi : A \to B$
the following are equivalent
\begin{enumerate}
\item $u \otimes 1 : M \otimes_A B \to N \otimes_A B$ is an
isomorphism, and
\item $\varphi(I) = 0$.
\end{enumerate}
\end{lemma}

\begin{proof}
As $M$ is projective we can find a projective $A$-module $C$
such that $F = M \oplus C$ is a free $A$-module.
By replacing $u$ by $u \oplus 1 : F = M \oplus C \to N \oplus C$
we see that we may assume $M$ is free. In this case let $I$ be
the ideal of $A$ generated by coefficients of all the elements of
$\Ker(u)$ with respect to some (fixed) basis of $M$.
The reason this works is that, since $u$ is surjective and
$\otimes_A B$ is right exact, $\Ker(u \otimes 1)$ is
the image of $\Ker(u) \otimes_A B$ in $M \otimes_A B$.
\end{proof}


\begin{lemma}
\label{lemma-universally-separating}
Let $S$ be a scheme.
Let $g : X' \to X$ be a flat morphism of schemes over $S$
with $X$ locally of finite type over $S$.
Let $\mathcal{F}$ be a finite type quasi-coherent $\mathcal{O}_X$-module
which is flat over $S$. If $\text{Ass}_{X/S}(\mathcal{F}) \subset g(X')$
then the canonical map
$$
\mathcal{F} \longrightarrow g_*g^*\mathcal{F}
$$
is injective, and remains injective after any base change.
\end{lemma}

\begin{proof}
The final assertion means that $\mathcal{F}_T \to (g_T)_*g_T^*\mathcal{F}_T$
is injective for any morphism $T \to S$. The assumption
$\text{Ass}_{X/S}(\mathcal{F}) \subset g(X')$ is preserved by base change, see
Divisors, Lemma \href{https://stacks.math.columbia.edu/tag/05DC}{lemma base change relative assassin (Tag 05DC)} and
Remark \href{https://stacks.math.columbia.edu/tag/05KL}{remark base change relative assassin (Tag 05KL)}.
The same holds for the assumption of flatness and finite type.
Hence it suffices to prove the injectivity of the displayed arrow.
Let $\mathcal{K} = \Ker(\mathcal{F} \to g_*g^*\mathcal{F})$.
Our goal is to prove that $\mathcal{K} = 0$.
In order to do this it suffices to prove that
$\text{WeakAss}_X(\mathcal{K}) = \emptyset$, see
Divisors, Lemma \href{https://stacks.math.columbia.edu/tag/05AP}{lemma weakly ass zero (Tag 05AP)}.
We have
$\text{WeakAss}_X(\mathcal{K}) \subset \text{WeakAss}_X(\mathcal{F})$, see
Divisors, Lemma \href{https://stacks.math.columbia.edu/tag/05AN}{lemma ses weakly ass (Tag 05AN)}.
As $\mathcal{F}$ is flat we see from
Lemma \ref{lemma-bourbaki-finite-type-general-base}
that $\text{WeakAss}_X(\mathcal{F}) \subset \text{Ass}_{X/S}(\mathcal{F})$.
By assumption any point $x$ of $\text{Ass}_{X/S}(\mathcal{F})$
is the image of some $x' \in X'$. Since $g$ is flat the
local ring map $\mathcal{O}_{X, x} \to \mathcal{O}_{X', x'}$
is faithfully flat, hence the map
$$
\mathcal{F}_x
\longrightarrow
g^*\mathcal{F}_{x'} =
\mathcal{F}_x \otimes_{\mathcal{O}_{X, x}} \mathcal{O}_{X', x'}
$$
is injective (see
Algebra, Lemma \href{https://stacks.math.columbia.edu/tag/05CK}{lemma faithfully flat universally injective (Tag 05CK)}).
This implies that $\mathcal{K}_x = 0$ as desired.
\end{proof}


\begin{lemma}
\label{lemma-iso-sheaf}
In Situation \ref{situation-iso}.
\begin{enumerate}
\item Each of the functors $F_{iso}$, $F_{inj}$, $F_{surj}$, $F_{zero}$
satisfies the sheaf property for the fpqc topology.
\item If $f$ is quasi-compact and $\mathcal{G}$ is of finite type,
then $F_{surj}$ is limit preserving.
\item If $f$ is quasi-compact and $\mathcal{F}$ of finite type, then
$F_{zero}$ is limit preserving.
\item If $f$ is quasi-compact, $\mathcal{F}$ is of finite type, and
$\mathcal{G}$ is of finite presentation, then $F_{iso}$ is limit preserving.
\end{enumerate}
\end{lemma}

\begin{proof}
Let $\{T_i \to T\}_{i \in I}$ be an fpqc covering of schemes over $S$.
Set $X_i = X_{T_i} = X \times_S T_i$ and $u_i = u_{T_i}$.
Note that $\{X_i \to X_T\}_{i \in I}$ is an fpqc covering of $X_T$, see
Topologies, Lemma \href{https://stacks.math.columbia.edu/tag/022D}{lemma fpqc (Tag 022D)}.
In particular, for every $x \in X_T$ there exists an $i \in I$
and an $x_i \in X_i$ mapping to $x$. Since
$\mathcal{O}_{X_T, x} \to \mathcal{O}_{X_i, x_i}$ is flat, hence
faithfully flat (see
Algebra, Lemma \href{https://stacks.math.columbia.edu/tag/00HR}{lemma local flat ff (Tag 00HR)})
we conclude that $(u_i)_{x_i}$ is injective, surjective, bijective, or zero
if and only if $(u_T)_x$ is injective, surjective, bijective, or zero.
Whence part (1) of the lemma.

\medskip\noindent
Proof of (2). Assume $f$ quasi-compact and $\mathcal{G}$ of finite type.
Let $T = \lim_{i \in I} T_i$ be a directed limit of affine $S$-schemes
and assume that $u_T$ is surjective.
Set $X_i = X_{T_i} = X \times_S T_i$ and
$u_i = u_{T_i} : \mathcal{F}_i = \mathcal{F}_{T_i}
\to \mathcal{G}_i = \mathcal{G}_{T_i}$.
To prove part (2) we have to show that $u_i$ is surjective for some $i$.
Pick $i_0 \in I$ and replace $I$ by $\{i \mid i \geq i_0\}$.
Since $f$ is quasi-compact the scheme $X_{i_0}$ is quasi-compact.
Hence we may choose affine opens $W_1, \ldots, W_m \subset X$
and an affine open covering
$X_{i_0} = U_{1, i_0} \cup \ldots \cup U_{m, i_0}$ such that
$U_{j, i_0}$ maps into $W_j$ under the projection morphism $X_{i_0} \to X$.
For any $i \in I$ let $U_{j, i}$ be the inverse image of $U_{j, i_0}$.
Setting $U_j = \lim_i U_{j, i}$ we see that $X_T = U_1 \cup \ldots \cup U_m$
is an affine open covering of $X_T$. Now it suffices to show, for a given
$j \in \{1, \ldots, m\}$ that $u_i|_{U_{j, i}}$ is surjective for some
$i = i(j) \in I$. Using
Properties, Lemma \href{https://stacks.math.columbia.edu/tag/01PB}{lemma finite type module (Tag 01PB)}
this translates into the following algebra problem:
Let $A$ be a ring and let $u : M \to N$ be an $A$-module map.
Suppose that $R = \colim_{i \in I} R_i$ is a directed colimit
of $A$-algebras. If $N$ is a finite $A$-module and if
$u \otimes 1 : M \otimes_A R \to N \otimes_A R$ is surjective, then
for some $i$ the map
$u \otimes 1 : M \otimes_A R_i \to N \otimes_A R_i$ is surjective.
This is
Algebra, Lemma \href{https://stacks.math.columbia.edu/tag/05LI}{lemma module map property in colimit (Tag 05LI)} part (2).

\medskip\noindent
Proof of (3). Exactly the same arguments as given in the proof of (2)
reduces this to the following algebra problem:
Let $A$ be a ring and let $u : M \to N$ be an $A$-module map.
Suppose that $R = \colim_{i \in I} R_i$ is a directed colimit
of $A$-algebras. If $M$ is a finite $A$-module and if
$u \otimes 1 : M \otimes_A R \to N \otimes_A R$ is zero, then
for some $i$ the map
$u \otimes 1 : M \otimes_A R_i \to N \otimes_A R_i$ is zero.
This is
Algebra, Lemma \href{https://stacks.math.columbia.edu/tag/05LI}{lemma module map property in colimit (Tag 05LI)} part (1).

\medskip\noindent
Proof of (4).
Assume $f$ quasi-compact and $\mathcal{F}, \mathcal{G}$ of finite presentation.
Arguing in exactly the same manner as in the previous paragraph
(using in addition also
Properties, Lemma \href{https://stacks.math.columbia.edu/tag/01PC}{lemma finite presentation module (Tag 01PC)})
part (3) translates into the following algebra statement:
Let $A$ be a ring and let $u : M \to N$ be an $A$-module map.
Suppose that $R = \colim_{i \in I} R_i$ is a directed colimit
of $A$-algebras. Assume $M$ is a finite $A$-module, $N$ is a finitely
presented $A$-module, and
$u \otimes 1 : M \otimes_A R \to N \otimes_A R$ is an isomorphism.
Then for some $i$ the map
$u \otimes 1 : M \otimes_A R_i \to N \otimes_A R_i$ is an isomorphism.
This is
Algebra, Lemma \href{https://stacks.math.columbia.edu/tag/05LI}{lemma module map property in colimit (Tag 05LI)} part (3).
\end{proof}



\noindent
We intentionally restrict ourselves here to morphisms which are
of finite type and not just morphisms which are locally of
finite type, see
Remark \href{https://stacks.math.columbia.edu/tag/05J5}{remark discuss finite type (Tag 05J5)}
for a discussion. In the situation of the definition
Lemma \ref{lemma-pure-along-X-s}
tells us that the following are equivalent
\begin{enumerate}
\item $\mathcal{F}$ is pure along $X_s$,
\item there is no impurity $(g : T \to S, t' \leadsto t, \xi)$ with $g$
quasi-finite at $t$,
\item there does not exist any impurity of the form
$(S^h \to S, s' \leadsto s, \xi)$, where $S^h$ is the henselization
of $S$ at $s$.
\end{enumerate}
If we denote $X^h = X \times_S S^h$ and $\mathcal{F}^h$ the pullback
of $\mathcal{F}$ to $X^h$, then we can formulate the last condition
in the following more positive way:
\begin{enumerate}
\item[(4)] All points of $\text{Ass}_{X^h/S^h}(\mathcal{F}^h)$ specialize
to points of $X_s$.
\end{enumerate}
In particular, it is clear that $\mathcal{F}$ is pure along $X_s$
if and only if the pullback of $\mathcal{F}$ to
$X \times_S \Spec(\mathcal{O}_{S, s})$ is pure along $X_s$.



\begin{lemma}
\label{algebra-lemma-adapted-submodule}
Let $R \to S$ be a ring map, and let $M$ be an $R$-module.  Suppose $M
\otimes_R S = \bigoplus_{i \in I} Q_i$ is a direct sum of countably generated
$S$-modules $Q_i$.  If $N$ is a countably generated submodule of $M$, then
there is a countably generated submodule $N'$ of $M$ such that $N' \supset N$
and $\Im(N' \otimes_R S \to M \otimes_R S) =
\bigoplus_{i \in I'} Q_i$ for some subset $I' \subset I$.
\end{lemma}

\begin{proof}
Let $N'_0 = N$.  We construct by induction an increasing sequence of countably
generated submodules $N'_{\ell} \subset M$ for $\ell = 0, 1, 2, \ldots$
such that: if $I'_{\ell}$ is the set of $i \in I$ such that the projection of
$\Im(N'_{\ell} \otimes_R S \to M \otimes_R S)$ onto $Q_i$ is
nonzero, then $\Im(N'_{\ell + 1} \otimes_R S \to M \otimes_R
S)$ contains $Q_i$ for all $i \in I'_{\ell}$.  To construct $N'_{\ell + 1}$
from $N'_\ell$, let $Q$ be the sum of (the countably many) $Q_i$ for
$i \in I'_{\ell}$, choose $P$ as in Lemma
\ref{algebra-lemma-lift-countably-generated-submodule}, and then let $N'_{\ell + 1} =
N'_{\ell} + P$.  Having constructed the $N'_{\ell}$, just take $N' =
\bigcup_{\ell} N'_{\ell}$ and $I' = \bigcup_{\ell} I'_{\ell}$.
\end{proof}


\begin{lemma}
\label{algebra-lemma-lift-countably-generated-submodule}
Let $R \to S$ be a ring map, let $M$ be an $R$-module, and let $Q$ be a
countably generated $S$-submodule of $M \otimes_R S$.  Then there exists a
countably generated $R$-submodule $P$ of $M$ such that
$\Im(P \otimes_R S \to M \otimes_R S)$ contains $Q$.
\end{lemma}

\begin{proof}
Let $y_1, y_2, \ldots$ be generators for $Q$ and write $y_j = \sum_k x_{jk}
\otimes s_{jk}$ for some $x_{jk} \in M$ and $s_{jk} \in S$. Then take $P$ be
the submodule of $M$ generated by the $x_{jk}$.
\end{proof}


\begin{lemma}
\label{algebra-lemma-ffdescent-ML}
\begin{reference}
Email from Juan Pablo Acosta Lopez dated 12/20/14.
\end{reference}
Let $R \to S$ be a faithfully flat ring map. Let $M$ be an $R$-module. If the
$S$-module $M \otimes_R S$ is Mittag-Leffler, then $M$ is Mittag-Leffler.
\end{lemma}

\begin{proof}
Write $M = \colim_{i\in I} M_i$ as a directed colimit
of finitely presented $R$-modules $M_i$.
Using Proposition \ref{algebra-proposition-ML-characterization}, we see that we
have to prove that for each $i \in I$ there exists $i \leq j$, $j\in I$
such that $M_i\rightarrow M_j$ dominates $M_i\rightarrow M$.

\medskip\noindent
Take $N$ the pushout
$$
\xymatrix{
M_i  \ar[r] \ar[d] & M_j \ar[d] \\
M \ar[r] & N
}
$$
Then the lemma is equivalent to the existence of $j$ such that
$M_j\rightarrow N$ is universally injective, see
Lemma \href{https://stacks.math.columbia.edu/tag/0AUM}{lemma domination universally injective (Tag 0AUM)}.
Observe that the tensorization by $S$
$$
\xymatrix{
M_i\otimes_R S  \ar[r] \ar[d] & M_j\otimes_R S \ar[d] \\
M\otimes_R S \ar[r] & N\otimes_R S
}
$$
Is a pushout diagram.  So because
$M \otimes_R S = \colim_{i\in I} M_i \otimes_R S$
expresses $M\otimes_R S$ as a colimit of $S$-modules of
finite presentation, and $M\otimes_R S$ is Mittag-Leffler,
there exists $j \geq i$ such that $M_j\otimes_R S\rightarrow N\otimes_R S$
is universally injective. So using that $R\rightarrow S$ is faithfully flat
we conclude that $M_j\rightarrow N$ is universally injective too.
\end{proof}


\begin{lemma}
\label{algebra-lemma-ffdescent-countable}
Let $R \to S$ be a faithfully flat ring map. Let $M$ be an $R$-module. If
the $S$-module $M \otimes_R S$ is countably generated, then $M$
is countably generated.
\end{lemma}

\begin{proof}
Say $M \otimes_R S$ is generated by the elements
$y_i$, $i = 1, 2, 3, \ldots$. Write $y_i = \sum_{j = 1, \ldots, n_i}
x_{ij} \otimes s_{ij}$ for some $n_i \geq 0$, $x_{ij} \in M$
and $s_{ij} \in S$. Denote $M' \subset M$ the submodule generated
by the countable collection of elements $x_{ij}$. Then
$M' \otimes_R S \to M \otimes_R S$ is surjective as the
image contains the generators $y_i$. Since $S$ is faithfully flat
over $R$ we conclude that $M' = M$ as desired.
\end{proof}


\begin{lemma}
\label{algebra-lemma-ffdescent-countable-projectivity}
Let $R \to S$ be a faithfully flat ring map.  Let $M$ be an $R$-module.
 If the $S$-module $M \otimes_R S$ is countably generated and projective,
then $M$ is countably generated and projective.
\end{lemma}

\begin{proof}
Follows from Lemmas \href{https://stacks.math.columbia.edu/tag/03C4}{lemma descend properties modules (Tag 03C4)},
\ref{algebra-lemma-ffdescent-ML}, and \ref{algebra-lemma-ffdescent-countable} and
Theorem \href{https://stacks.math.columbia.edu/tag/059Z}{theorem projectivity characterization (Tag 059Z)}.
\end{proof}


\begin{theorem}
\label{algebra-theorem-kaplansky-direct-sum}
Suppose $M$ is a direct sum of countably generated $R$-modules.  If $P$ is a
direct summand of $M$, then $P$ is also a direct sum of countably generated
$R$-modules.
\end{theorem}

\begin{proof}
Write $M = P \oplus Q$.  We are going to construct a Kaplansky d\'evissage
$(M_{\alpha})_{\alpha \in S}$ of $M$ which, in addition to the defining
properties (0)-(4), satisfies:
\begin{enumerate}
\item[(5)] Each $M_{\alpha}$ is a direct summand of $M$;
\item[(6)] $M_{\alpha} = P_{\alpha} \oplus Q_{\alpha}$, where $P_{\alpha} =P
\cap M_{\alpha}$ and $Q_\alpha = Q \cap M_{\alpha}$.
\end{enumerate}
(Note: if properties (0)-(2) hold, then in fact property (3) is equivalent to
property (5).)

\medskip\noindent
To see how this implies the theorem, it is enough to show that
$(P_{\alpha})_{\alpha \in S}$ forms a Kaplansky d\'evissage of $P$.  Properties
(0), (1), and (2) are clear.  By (5) and (6) for $(M_{\alpha})$, each
$P_{\alpha}$ is a direct summand of $M$.  Since $P_{\alpha} \subset P_{\alpha +
1}$, this implies $P_{\alpha}$ is a direct summand of $P_{\alpha + 1}$; hence
(3) holds for $(P_{\alpha})$.  For (4), note that
$$
M_{\alpha + 1}/M_{\alpha} \cong P_{\alpha + 1}/P_{\alpha} \oplus
Q_{\alpha + 1}/Q_{\alpha},
$$
so $P_{\alpha + 1}/P_{\alpha}$ is countably generated because this is true of
$M_{\alpha + 1}/M_{\alpha}$.

\medskip\noindent
It remains to construct the $M_{\alpha}$.  Write $M = \bigoplus_{i \in I} N_i$
where each $N_i$ is a countably generated $R$-module.  Choose a well-ordering
of $I$.  By transfinite recursion we are going to define an increasing family
of submodules $M_{\alpha}$ of $M$, one for each ordinal $\alpha$, such that
$M_{\alpha}$ is a direct sum of some subset of the $N_i$.

\medskip\noindent
For $\alpha = 0$ let $M_{0} = 0$.  If $\alpha$ is a limit ordinal and
$M_{\beta}$ has been defined for all $\beta < \alpha$, then define $M_{\alpha}
= \bigcup_{\beta < \alpha} M_{\beta}$.  Since each $M_{\beta}$ for $\beta <
\alpha$ is a direct sum of a subset of the $N_i$, the same will be true of
$M_{\alpha}$.  If $\alpha + 1$ is a successor ordinal and $M_{\alpha}$ has been
defined, then define $M_{\alpha + 1}$ as follows.  If $M_{\alpha} = M$, then let
$M_{\alpha + 1} = M$.  If not, choose the smallest $j \in I$ such that $N_j$ is
not contained in $M_{\alpha}$.  We will construct an infinite matrix $(x_{mn}),
m, n = 1, 2, 3, \ldots$ such that:
\begin{enumerate}
\item $N_j$ is contained in the submodule of $M$ generated by the entries
$x_{mn}$;
\item if we write any entry $x_{k\ell}$ in terms of its $P$- and
$Q$-components, $x_{k\ell} = y_{k\ell} + z_{k\ell}$, then the matrix $(x_{mn})$
contains a set of generators for each $N_i$ for which $y_{k\ell}$ or
$z_{k\ell}$ has nonzero component.
\end{enumerate}
Then we define $M_{\alpha + 1}$ to be the submodule of $M$ generated by
$M_{\alpha}$ and all $x_{mn}$; by property (2) of the matrix $(x_{mn})$,
$M_{\alpha + 1}$ will be a direct sum of some subset of the $N_i$.
To construct the matrix $(x_{mn})$, let $x_{11}, x_{12}, x_{13}, \ldots$
be a countable set of generators for $N_j$. Then if
$x_{11} = y_{11} + z_{11}$ is the decomposition into $P$- and
$Q$-components, let $x_{21}, x_{22}, x_{23}, \ldots$ be a countable
set of generators for the sum of the $N_i$ for which $y_{11}$ or $z_{11}$ have
nonzero component.  Repeat this process on $x_{12}$ to get elements $x_{31},
x_{32}, \ldots$, the third row of our matrix.  Repeat on $x_{21}$ to get the
fourth row, on $x_{13}$ to get the fifth, and so on, going down along
successive anti-diagonals as indicated below:
$$
\left(
\vcenter{
\xymatrix@R=2mm@C=2mm{
x_{11} & x_{12} \ar[dl] & x_{13} \ar[dl] & x_{14} \ar[dl] & \ldots  \\
x_{21} & x_{22} \ar[dl] & x_{23} \ar[dl] & \ldots  \\
x_{31} & x_{32} \ar[dl] & \ldots \\
x_{41} & \ldots \\
\ldots
}
}
\right).
$$

\medskip\noindent
Transfinite induction on $I$ (using the fact that we constructed
$M_{\alpha + 1}$
to contain $N_j$ for the smallest $j$ such that $N_j$ is not contained in
$M_{\alpha}$) shows that for each $i \in I$, $N_i$ is contained in some
$M_{\alpha}$.  Thus, there is some large enough ordinal $S$ satisfying: for
each $i \in I$ there is $\alpha \in S$ such that $N_i$ is contained in
$M_{\alpha}$.  This means $(M_{\alpha})_{\alpha \in S}$ satisfies property (1)
of a Kaplansky d\'evissage of $M$.  The family $(M_{\alpha})_{\alpha \in S}$
moreover satisfies the other defining properties, and also (5) and (6) above:
properties (0), (2), (4), and (6) are clear by construction;  property (5) is
true because each $M_{\alpha}$ is by construction a direct sum of some $N_i$;
and (3) is implied by (5) and the fact that $M_{\alpha} \subset M_{\alpha + 1}$.
\end{proof}


\begin{theorem}
\label{theorem-flattening-map}
In
Situation \ref{situation-iso}
assume
\begin{enumerate}
\item $f$ is of finite presentation,
\item $\mathcal{F}$ is of finite presentation, flat over $S$, and
pure relative to $S$, and
\item $u$ is surjective.
\end{enumerate}
Then $F_{iso}$ is representable by a closed immersion $Z \to S$.
Moreover $Z \to S$ is of finite presentation if $\mathcal{G}$ is
of finite presentation.
\end{theorem}

\begin{proof}
We will use without further mention that $\mathcal{F}$ is universally pure
over $S$, see
Lemma \ref{lemma-finite-type-flat-pure-along-fibre-is-universal}.
By
Lemma \ref{lemma-iso-sheaf}
and
Descent, Lemmas \href{https://stacks.math.columbia.edu/tag/03I0}{lemma closed immersion (Tag 03I0)} and
\href{https://stacks.math.columbia.edu/tag/02W5}{lemma descent data sheaves (Tag 02W5)}
the question is local for the \'etale topology on $S$.
Hence it suffices to prove, given $s \in S$, that there exists
an \'etale neighbourhood of $(S, s)$ so that the theorem holds.

\medskip\noindent
Using
Lemma \ref{lemma-finite-presentation-flat-along-fibre}
and after replacing $S$ by an elementary \'etale neighbourhood of $s$
we may assume there exists a commutative diagram
$$
\xymatrix{
X \ar[dr] & & X' \ar[ll]^g \ar[ld] \\
& S &
}
$$
of schemes of finite presentation over $S$,
where $g$ is \'etale, $X_s \subset g(X')$, the schemes $X'$ and $S$ are affine,
$\Gamma(X', g^*\mathcal{F})$ a projective $\Gamma(S, \mathcal{O}_S)$-module.
Note that $g^*\mathcal{F}$ is universally pure over $S$, see
Lemma \ref{lemma-affine-locally-projective-pure}.
Hence by
Lemma \ref{lemma-criterion}
we see that the open $g(X')$ contains the points of
$\text{Ass}_{X/S}(\mathcal{F})$ lying over $\Spec(\mathcal{O}_{S, s})$.
Set
$$
E = \{t \in S \mid \text{Ass}_{X_t}(\mathcal{F}_t) \subset g(X') \}.
$$
By
More on Morphisms,
Lemma \href{https://stacks.math.columbia.edu/tag/05KR}{lemma relative assassin constructible (Tag 05KR)}
$E$ is a constructible subset of $S$. We have seen that
$\Spec(\mathcal{O}_{S, s}) \subset E$. By
Morphisms, Lemma \href{https://stacks.math.columbia.edu/tag/05LW}{lemma constructible containing open (Tag 05LW)}
we see that $E$ contains an open neighbourhood of $s$. Hence after
replacing $S$ by a smaller affine neighbourhood of $s$ we may assume that
$\text{Ass}_{X/S}(\mathcal{F}) \subset g(X')$.

\medskip\noindent
Since we have assumed that $u$ is surjective we have $F_{iso} = F_{inj}$. From
Lemma \ref{lemma-universally-separating}
it follows that $u : \mathcal{F} \to \mathcal{G}$ is injective if and only if
$g^*u : g^*\mathcal{F} \to g^*\mathcal{G}$ is injective, and the same remains
true after any base change. Hence we have reduced to the case where,
in addition to the assumptions in the theorem, $X \to S$ is a morphism of
affine schemes and $\Gamma(X, \mathcal{F})$ is a projective
$\Gamma(S, \mathcal{O}_S)$-module. This case follows immediately from
Lemma \ref{lemma-flattening-module-map}.

\medskip\noindent
To see that $Z$ is of finite presentation if $\mathcal{G}$ is of finite
presentation, combine
Lemma \ref{lemma-iso-sheaf} part (4)
with
Limits, Remark \href{https://stacks.math.columbia.edu/tag/05LX}{remark limit preserving (Tag 05LX)}.
\end{proof}
\end{document}
